\documentclass[11pt]{article}

%%% LATEX COMMON Packages
\usepackage[bookmarks,colorlinks,breaklinks]{hyperref}  % PDF hyperlinks, with coloured links
\hypersetup{linkcolor=blue,citecolor=blue,filecolor=blue,urlcolor=blue} % all blue links
\usepackage{fullpage}
\usepackage{amsmath,amsfonts,amsthm,amssymb,xspace,bm, verbatim}
%\usepackage{natbib}
\usepackage{graphicx}
%\usepackage{bibspacing}
\usepackage{url,algorithm,algorithmic}

%Some options.
%\setlength{\bibspacing}{1mm}
%\setlength{\textwidth}{6.5 in}
%\setlength{\textheight}{9 in}
%\setlength{\oddsidemargin}{0in}
%\setlength{\topmargin}{.5in}
%\addtolength{\voffset}{-.5in}

%%%% Reference Definitions
\newcommand{\newref}[2][]{\hyperref[#2]{#1~\ref*{#2}}}
%\newcommand{\newref}[2][]{#1~\ref{#2}}
\renewcommand{\eqref}[1]{\hyperref[#1]{(\ref*{#1})}}
%\numberwithin{equation}
  
%%%%% THEOREM STYLE DEFINITIONS
\theoremstyle{plain}
\newtheorem{theorem}{Theorem}[section]
\newtheorem{lemma}[theorem]{Lemma}
%\newtheorem{claim}[theorem]{Claim}
%\newtheorem{proposition}[theorem]{Proposition}
%\newtheorem{corollary}[theorem]{Corollary}
%\newtheorem{conjecture}[theorem]{Conjecture}
\newtheorem{definition}[theorem]{Definition}
%\newtheorem{fact}[theorem]{Fact}
\theoremstyle{definition}
%\newtheorem{question}[theorem]{Question}
%\newtheorem{remark}[theorem]{Remark}

%Standard notations
\DeclareMathOperator*{\pr}{\mathsf{Pr}} 
\renewcommand{\iff}{\Leftrightarrow}
\DeclareMathOperator*{\ex}{\mathbb{E}}
\newcommand{\rgta}{\rightarrow}
\newcommand{\iprod}[2]{\langle #1, #2 \rangle}   

% %Real Stuff
\newcommand{\rn}{\mathbb{R}^n}
\newcommand{\rnn}{\mathbb{R}^{n \times n}}
\newcommand{\rmn}{\mathbb{R}^{m \times n}}
\newcommand{\reals}{\mathbb{R}}

\newcommand{\e}{\bm{e}}
\newcommand{\dpm}{\{1,-1\}}
\newcommand{\dpmn}{\dpm^n}
\newcommand{\sign}{\mathsf{sign}}
\newcommand{\sgn}{\mathsf{sgn}}
\newcommand{\hwt}{H_{w,\theta}}
\newcommand{\infl}{\mathbb{I}}
\newcommand{\AS}{\mathbb{AS}}
\newcommand{\hh}{\mathcal{H}}
\newcommand{\pset}{2^{[n]}}
\newcommand{\sumd}{\sum_{I,|I| \leq d}}
\newcommand{\gsn}{\mathcal{N}(0,1)^n}
\newcommand{\calC}{{\mathcal C}}
\newcommand{\Stab}{{\mathbb S}}
\newcommand{\GNS}{\mathbb{GNS}}
\newcommand{\NS}{\mathbb{NS}}
\newcommand{\apx}[1]{\stackrel{#1}{\approx}}
\newcommand{\apxl}[1]{\stackrel{#1}{\leq}}
\newcommand{\eps}{\epsilon}
%\renewcommand{\epsilon}{\varepsilon}
\newcommand{\Tr}{\operatorname{Tr}}
\newcommand{\Var}{\operatorname{Var}}
\newcommand{\Cov}{\operatorname{Cov}}
\newcommand{\Per}{\operatorname{Per}}
\newcommand{\Diff}{\operatorname{D}}
%\newcommand{\sgn}{\mathrm{sgn}}
\newcommand{\note}[1]{\marginpar{\tiny *note in TeX*}}
%\newcommand{\note}[1]{}
\newcommand{\ignore}[1]{}
\newcommand{\Avg}{\mathrm{Avg}}
\newcommand{\mmmatrix}{{\begin{array}{cc}
1 & \eta \\
\eta & 1\end{array} }}
\newcommand{\calD}{{\cal D}}
\newcommand{\calN}{{\cal N}}
\newcommand{\sqrtn}{\sqrt{n}}
\newcommand{\tO}{\tilde{O}}
\renewcommand{\phi}{\varphi}

\newcommand{\bits}{\{-1,1\}}
\newcommand{\bn}{\bits^n}
\newcommand{\isafunc}{: \bn \rightarrow \bits}
\newcommand{\R}{\mathbb{R}}
\newcommand{\Z}{\mathbb{Z}}
\newcommand{\N}{\mathbb{N}}
\newcommand{\NN}{{{\cal N}^n}}
%\newcommand{\E}{\operatorname{{\bf E}}}
%\newcommand{\Ex}{\mathop{{\bf E}\/}}
%\renewcommand{\Pr}{\operatorname{{\bf Pr}}}
%\newcommand{\Prx}{\mathop{{\bf Pr}\/}}
\newcommand{\littlesum}{\mathop{\textstyle \sum}}
\newcommand{\littlefrac}{\textstyle \frac}
\newcommand{\wt}[1]{\widetilde{#1}}
\newcommand{\poly}{\mathrm{poly}}
\newcommand{\Maj}{\mathrm{Maj}}
\newcommand{\Inf}{\mathrm{Inf}}
\newcommand{\half}{{\textstyle \frac 1 2}}
\newcommand{\wh}{\widehat}
\newcommand{\eqdef}{\stackrel{\textrm{def}}{=}}
\newcommand{\la}{\langle}
\newcommand{\ra}{\rangle}
\newcommand{\capp}{\mathrm{cap}}
\newcommand{\perim}{\mathrm{\Gamma}}
\newcommand{\surf}{\mathrm{\Gamma}}
\newcommand{\dist}{\mathrm{dist}}
\newcommand{\vol}{\mathrm{vol}}
\newcommand{\opt}{\mathsf{opt}}
\newcommand{\bdry}{\partial}
\newcommand{\ind}[1]{\mathbf{1}_{#1}}
\newcommand{\sphere}{{\cal S}}
\newcommand{\C}{{{\cal C}}}
\newcommand{\D}{{{\cal D}}}
\newcommand{\X}{{\cal X}}

%%%Commands for this document.
\newcommand{\mtx}[1]{{\mathsf{#1}}}
\newcommand{\frob}[1]{\left\| {#1} \right\|_\text{F}}
\newcommand{\onenorm}[1]{\left\| {#1} \right\|_1}
\newcommand{\twonorm}[1]{\left\| {#1} \right\|_2}
\newcommand{\infnorm}[1]{\left\| {#1} \right\|_{\infty}}
\DeclareMathOperator{\trace}{Tr}
\DeclareMathOperator*{\argmin}{argmin}
\DeclareMathOperator*{\argmax}{argmax}
\newtheorem{conj}{Conjecture}[section]
\newcommand{\vct}[1]{\bm{#1}}
\newcommand{\OX}{\overline{\mtx{X}}}
\newcommand{\OZ}{\overline{\mtx{Z}}}
\newcommand{\calA}{\mathcal{A}}
\newcommand{\nubar}{\overline{\bm{\nu}}}
\newcommand{\fro}[1]{\|#1\|_F}
\newcommand{\E}{\mathbb{E}}
\newcommand{\expec}[1]{\mathbb{E}\left[#1\right]}
\newcommand{\prob}[1]{\mathbb{P}\left[#1\right]}
\newcommand{\aff}{{\mathcal A}}
\newcommand{\Vt}{V^{t+1}}
\newcommand{\Vht}{\widehat{V}^{t+1}}
\newcommand{\Ut}{U^{t}}
\newcommand{\Uht}{\widehat{U}^{t}}
\newcommand{\Rt}{R^{(t+1)}}
\newcommand{\Vo}{V^*}
\newcommand{\So}{\Sigma^*}
\newcommand{\Uo}{U^*}
\newcommand{\Un}{U^N}
\newcommand{\Vn}{V^N}
\newcommand{\Sn}{\Sigma^N}
\newcommand{\sn}{\sigma^N}
\newcommand{\un}{u^N}
\newcommand{\so}{\sigma^*}
\newcommand{\vto}{v^{t+1}}
\newcommand{\ut}{u^t}
\newcommand{\utt}{(u^t)^\dag }
\newcommand{\uo}{u^* }
\newcommand{\uot}{(u^*)^\dag }
\newcommand{\vo}{v^* }
\newcommand{\vot}{(v^*)^\dag }
\newcommand{\vtto}{(v^{t+1})^\dag }
\newcommand{\Ai}{A_i}
\newcommand{\Ait}{A_i^\dag }
\newcommand{\wvto}{\widehat{v}^{t+1}}
\newcommand{\wvtto}{(\widehat{v}^{t+1})^\dag}
\newcommand{\ip}[2]{\langle #1, #2 \rangle}
\newcommand{\vt}{\vto}
\newcommand{\vht}{\wvto}
\newcommand{\order}[1]{O\left({#1}\right)}
\newcommand{\as}{AltMinSense\ }
\newcommand{\ase}{AltMinSense}
\newcommand{\amc}{AltMinComplete\ }
\newcommand{\amce}{AltMinComplete}
\newcommand{\Vw}{\widehat{V}}
\newcommand{\Uw}{\widehat{U}}
\begin{document}

\title{ Low-rank Matrix Completion using Alternating Minimization}
%\subtitle{[Extended Abstract]}
%
% You need the command \numberofauthors to handle the 'placement
% and alignment' of the authors beneath the title.
%
% For aesthetic reasons, we recommend 'three authors at a time'
% i.e. three 'name/affiliation blocks' be placed beneath the title.
%
% NOTE: You are NOT restricted in how many 'rows' of
% "name/affiliations" may appear. We just ask that you restrict
% the number of 'columns' to three.
%
% Because of the available 'opening page real-estate'
% we ask you to refrain from putting more than six authors
% (two rows with three columns) beneath the article title.
% More than six makes the first-page appear very cluttered indeed.
%
% Use the \alignauthor commands to handle the names
% and affiliations for an 'aesthetic maximum' of six authors.
% Add names, affiliations, addresses for
% the seventh etc. author(s) as the argument for the
% \additionalauthors command.
% These 'additional authors' will be output/set for you
% without further effort on your part as the last section in
% the body of your article BEFORE References or any Appendices.
%\numberofauthors{3} %  in this sample file, there are a *total*
% of EIGHT authors. SIX appear on the 'first-page' (for formatting
% reasons) and the remaining two appear in the \additionalauthors section.
%
\author{
% You can go ahead and credit any number of authors here,
% e.g. one 'row of three' or two rows (consisting of one row of three
% and a second row of one, two or three).
%
% The command \alignauthor (no curly braces needed) should
% precede each author name, affiliation/snail-mail address and
% e-mail address. Additionally, tag each line of
% affiliation/address with \affaddr, and tag the
% e-mail address with \email.
%
% 1st. author
Prateek Jain\\
{Microsoft Research India, Bangalore}\\
{prajain@microsoft.com}
% 2nd. author
\and
Praneeth Netrapalli\\
{The University of Texas at Austin}\\
{praneethn@utexas.edu}
% 3rd. author
\and
Sujay Sanghavi\\
{The University of Texas at Austin}\\
{sanghavi@mail.utexas.edu}
}
% There's nothing stopping you putting the seventh, eighth, etc.
% author on the opening page (as the 'third row') but we ask,
% for aesthetic reasons that you place these 'additional authors'
% in the \additional authors block, viz.
% Just remember to make sure that the TOTAL number of authors
% is the number that will appear on the first page PLUS the
% number that will appear in the \additionalauthors section.

\maketitle
%\onecolumn
\begin{abstract}
%In this paper, we investigate the performance of Alternating Minimization, a popular empirical approach to low-rank recovery. 
Alternating minimization represents a widely applicable and empirically successful approach for finding low-rank matrices that best fit the given data. %, and has enjoyed immense empirical success. %popular empirical approach for finding low-rank matrices that best fit data, in a variety of applications. 
For example, for the problem of low-rank matrix completion, this method is believed to be one of the most accurate and efficient, and formed a major component of the winning entry in the Netflix Challenge \cite{Koren09,KorenBV09}. 

In the alternating minimization approach, the low-rank target matrix is written in a {\em bi-linear form}, i.e. $X = UV^\dag$; the algorithm then alternates between finding the best $U$ and the best $V$. Typically, each alternating step in isolation is convex and tractable. However the overall problem becomes non-convex and  there has been almost no theoretical understanding of when this approach yields a good result.

In this paper we present first theoretical analysis of  the performance of alternating minimization for matrix completion, and the related problem of matrix sensing. For both these problems, celebrated recent results have shown that they become well-posed and tractable once certain (now standard) conditions are imposed on the problem. We show that alternating minimization {\em also} succeeds under similar conditions. Moreover, compared to existing results, our paper shows that alternating minimization guarantees faster (in particular, geometric) convergence to the true matrix, while allowing a simpler analysis. %Alternating minimization is a natural fit for matrix completion, and has been empirically popular; ours is the only analysis of when it succeeds. 

%In this paper, we consider the problem of low-rank matrix completion where the goal is to recover back a low-rank matrix using a small number of random entries from the matrix. In particular, we study the popular alternating minimization method for matrix completion. Alternating minimization procedure characeterizes a low-rank matrix $X$ as $X=UV^T$ and alternately solves for $U$ and $V$. However, as the procedure is optimizing a non-convex problem, it was not apriori known if the procedure always converges to the true matrix. However, extensive empirical studies show that the procedure is one of the most accurate as well as efficient algorithm for matrix completion. 
%
%In this paper, we show that, despite optimizing a non-convex problem, alternating minimization indeed provably solves low-rank matrix completion problem under standard assumptions  \cite{CandesR2007, KeshavanOM2009}. Our analysis is significantly simpler than the analysis for other known  methods for matrix completion \cite{CandesR2007, KeshavanOM2009}. Furthermore, unlike existing methods, our analysis provides linear convergence rate to the true matrix. 
%
%Moreover, we also study a related and more general problem of low-rank matrix sensing where the goal is to recover a low-rank matrix using a small number of linear measurements. Here also, we show that the alternating minimization procedure recovers back the underlying low-rank matrix under similar assumptions (of restricted isometry property) as existing methods \cite{RechtFP2007, JainMD10}, while also providing geometric convergence rate. 

\end{abstract}

% A category with the (minimum) three required fields
%\category{H.4}{Information Systems Applications}{Miscellaneous}
%A category including the fourth, optional field follows...
%\category{D.2.8}{Software Engineering}{Metrics}[complexity measures, performance measures]

%\terms{Theory}

%\keywords{ACM proceedings, \LaTeX, text tagging} % NOT required for Proceedings

\section{Introduction}
Finding a low-rank matrix to fit / approximate observations is a fundamental task in data analysis. In a slew of applications, a popular empirical approach has been to represent the target rank $k$ matrix $X\in \mathbb{R}^{m\times n}$  in a {\em bi-linear form} $X = UV^\dag$, where $U\in \mathbb{R}^{m\times k}$ and $V\in \mathbb{R}^{n\times k}$.  Typically, this is done for two  reasons: \\
{\em (a) Size and computation:} If the rank $k$ of the target matrix (to be estimated) is much smaller than $m,n$, then  $U,V$ are significantly smaller than  $X$ and hence are more efficient to optimize for. This is crucial for several practical applications, e.g., recommender systems where one routinely encounters matrices with billions of entries. \\%Several modern applications involve estimating matrices with millions or billions of entries \cite{Koren09}. Such matrices are often impossible to load in memory, and sometimes even hard to store on disk. However, if the rank $k$ of the matrix is much smaller than $m,n$, then  $U,V$ are significantly smaller than  $X$, and could be easier to optimize. For example, in recommendation systems one may want to find a low-rank ``ratings" matrix between millions of users and hundreds  of thousands of items. Here, storing and optimizing a full ratings matrix $X$ would be infeasible, but finding factors $U$, $V$ from the given sparse ratings data can still be practical \cite{KorenBV09}. \\%While the ratings data given is sparse, and the final low-rank matrix may be efficiently representable as $UV^\dag$, it may be hard to store and optimize a full matrix $X$. \\
{\em (b) Modeling:} In several applications, one would like to impose extra constraints on the target matrix, besides just low rank. Oftentimes, these constraints might be easier and more natural to impose on factors $U$, $V$.  For example, in Sparse PCA \cite{ZouHT06}, one looks for a low-rank $X$ that is the product of {\em sparse} $U$ and $V$. %Similarly, non-negative matrix factorization \cite{KimPar08} seeks $U$ and $V$ that are non-negative. Such constraints might be hard to impose in a problem parameterized by $X$, but can be naturally imposed in a problem formulation based on the bi-linear ($X=UV^\dag$) form (see e.g. \cite{KimPar08, ZouHT06}).%, but are hard to impose in a problem parameterized by an overall $X$.  \\

Due to the above two reasons, in several applications, the target matrix $X$ is parameterized by $X=UV^\dag$. For example, clustering \cite{KimPark08b}, sparse PCA \cite{ZouHT06} etc. 

Using the bi-linear parametrization of the target matrix $X$, the task of estimating $X$ now reduces to finding  $U$ and $V$ that, for example, minimize an error metric. The resulting problem is typically non-convex due to bi-linearity. Correspondingly, a popular approach has been to use {\bf alternating minimization:} iteratively keep one of $U,V$ fixed and optimize over the other, then switch and repeat, see e.g. \cite{KorenBV09}. While the overall problem is non-convex, each sub-problem is typically convex and can be solved efficiently.
%Infact, in many cases each intermediate problem becomes a least squares problem \cite{TakYouDeL77,KroDeL80}. Hence this approach is also popularly known as {\bf alternating least squares (ALS)}.

Despite wide usage of bi-linear representation and alternating minimization, there has been to date almost no theoretical understanding of when such a formulation works. Motivated by this disconnect between theory and practice in the estimation of  low-rank matrices, in this paper we provide the {\bf first guarantees for performance of alternating minimization}, for two low-rank matrix recovery problems: matrix completion, and matrix sensing. 

{\em Matrix completion} involves completing a low-rank matrix, by observing only a few of its elements. Its recent popularity, and primary motivation, comes from recommendation systems \cite{KorenBV09}, where the task is to complete a user-item ratings matrix using only a small number of ratings. As elaborated in Section~\ref{sec:results}, alternating minimization becomes particularly appealing for this problem as it provides a fast, distributed algorithm that can exploit both sparsity of ratings as well as the low-rank bi-linear parametrization of $X$. %, as now both the data and the matrix representation are of a size much smaller than the full ratings matrix.
%{\em Matrix completion} involves completing a low-rank matrix, by observing only a few of its elements. Its recent popularity, and primary motivation, comes from recommendation systems \cite{KorenBV09}. In the typical setup, the rows of the matrix represent the users of a recommendation system, the columns represent the items, and each element represents the level of preference (rating) that user shows for that item. It is thus routine to encounter matrices with millions of rows and columns. However, the observed ratings data is typically quite sparse: very few user-item pairs are actually rated. As elaborated in Section~\ref{sec:results}, alternating minimization becomes particularly appealing as a fast, distributed algorithm in this setting, as now both the data and the matrix representation are of a size much smaller than the full ratings matrix. 

{\em Matrix sensing} refers to the problem of recovering a low-rank matrix $M\in \mathbb{R}^{m\times n}$ from affine equations. That is, given $d$ linear measurements $b_i = tr(A_i^\dag M)$ and measurement matrices $A_i$'s, the goal is to recover back $M$. This problem is particularly interesting in the case of $d\ll mn$ and was first studied in \cite{RechtFP2007} and subsequently in \cite{JainMD10, LeeB10}. In fact,  matrix completion is a special case of this problem, where each observed entry in the matrix completion problem  represents one single-element measurement matrix $A_i$. 
%\langle A_i,M\rangle$, \footnote{Here $\langle A_i,M\rangle = tr(A_i^\dag M)$ refers to the standard matrix inner product.} 
%, we want to recover $M$ given $\{b_i,A_i\}_{i=1}^d$, in the case where $d$ is smaller than $mn$; this was first studied in {\bf \tt cite@@}, and subsequently in {\bf \tt cite@@}. Stated this way, matrix sensing is more general than matrix completion, which then refers to the case where each sample $A_i$ is a single-element matrix. 

Without any extra conditions, both  matrix sensing and matrix completion are ill-posed problems, with potentially multiple low-rank solutions, and are in general NP hard \cite{MekaJCD2008}. Current work on these problems thus impose some extra conditions, which  makes the problems both well defined, and amenable to solution via the respective proposed algorithms \cite{RechtFP2007, CandesR2007}; see Section \ref{sec:related} for more details. In this paper, we show that {\bf under similar conditions to the ones used by the existing methods,} alternating minimization also guarantees recovery of the true matrix; we also show that it requires only a small number of computationally cheap iterations and hence, as observed empirically, is computationally much more efficient than the existing methods. \\
{\bf Notations}: We represent a matrix by capital letter (e.g. $M$) and a vector by small letter ($u$). $u_i$ represents $i$-th element of $u$ and $U_{ij}$ denotes $(i,j)$-th entry of $U$. $U_i$ represents $i$-th column of $U$ and $U^{(i)}$ represents $i$-th row of $U$. $A^\dag$ denotes matrix transpose of $A$. $u=vec(U)$ represents vectorized $U$, i.e., $u=[U_1^\dag\ U_2^\dag\ \dots\ U_k^\dag]^\dag$.  $\|u\|_p$ denotes $L_p$ norm of $u$, i.e., $\|u\|_p=(\sum_i|u_i|^p)^{1/p}$. By default, $\|u\|$ denotes $L_2$ norm of $u$. $\|A\|_F$ denotes Frobenius norm of $A$, i.e., $\|vec(A)\|_2$. $\|A\|_2=\max_{x, \|x\|_2=1} \|Ax\|_2$ denotes spectral norm of $A$. $tr(A)$ denotes the trace (sum of diagonal elements) of square matrix $A$. Typically, $\widehat{U}$, $\widehat{V}$ represent factor matrices (i.e., $\widehat{U}\in \mathbb{R}^{m\times k}$ and $\widehat{V}\in \mathbb{R}^{n\times k}$) and $U$, $V$ represent their orthonormal basis.
\section{Our Results}\label{sec:results}
In this section, we will  first define the matrix sensing problem, and present our results for it. Subsequently, we will do the same for matrix completion. The matrix sensing setting -- i.e. recovery of any low-rank matrix from linear measurements that satisfy matrix RIP -- represents an easier analytical setting than matrix completion, but still captures several  key properties of the problem that helps us in developing an analysis for matrix completion.%  carries over to the 
We note that for either problem, ours represent the first global optimality guarantees for alternating minimization based algorithms. 
\subsection*{Matrix Sensing via Alternating Minimization}
Given $d$ linear measurements $b_i = \langle M,A_i\rangle = tr(A_i^\dag M),\ 1\leq i\leq d$ of an {\em unknown} rank-$k$ matrix $M\in \mathbb{R}^{m\times n}$ and the sensing matrices $A_i, 1\leq i\leq d$, the goal in matrix sensing is to recover back $M$. 
%The matrix sensing problem is the following: there is an unknown rank-$k$ matrix $M\in \mathbb{R}^{m\times n}$, of which we have $d$ linear measurements $b_i = \langle M,A_i\rangle = tr(A_i^\dag M)$, for $i=1,\ldots,d$. The task is to recover $M$ using the given measurement coefficients $\{b_i,A_i\}, i = 1,\ldots,d$. 
In the following we  collate these coefficients, so that $b\in \mathbb{R}^{d}$ is the vector of $b_i$'s, and $\aff(\cdot):\mathbb{R}^{m\times n}\rightarrow d$ is the corresponding linear map, with $b=\aff(M)$. With this notation, the Low-Rank Matrix Sensing problem is:
\renewcommand{\theequation}{LRMS}
\begin{equation}\label{eq:lrms}
 \hspace*{-10pt}\text{Find }X\in \rmn, \;\; s.t \;\; {\mathcal A}(X) = b,\ \  rank(X)\leq k.
\end{equation}
\renewcommand{\theequation}{\arabic{equation}}
\addtocounter{equation}{-1}
As in the existing work \cite{RechtFP2007} on this problem, we are interested in the  {\bf under-determined case}, where $d < mn$. Note that this problem is a strict generalization of the popular compressed sensing problem \cite{CandesT2005}; compressed sensing represents the case when $M$ is restricted to be a diagonal matrix.%, then this problem specializes to the popular compressed sensing problem \cite{CandesT05}. %Hence it is a generalization of the compressed sensing problem. 

For matrix sensing, alternating minimization approach involves representing $X$ as a product of two matrices $U\in \mathbb{R}^{m\times k}$ and $V\in \mathbb{R}^{n\times k}$, i.e., $X=UV^\dag$. If $k$ is (much) smaller than $m,n$, these matrices will be (much) smaller than $X$. With this bi-linear representation, alternating minimization can be viewed as an approximate way to solve the following non-convex optimization problem: 
\[
\min_{U\in \mathbb{R}^{m\times k}, V\in \mathbb{R}^{n\times k}}  \quad  \|\aff(UV^\dag) - b\|^2_2 
\]
As mentioned earlier, alternating minimization algorithm for matrix sensing  now alternately solves for $U$ and $V$ while fixing the other factor. See Algorithm~\ref{algo:sensing-altmin} for a pseudo-code of  \as algorithm that we analyze. 

We note  two key properties of \as: a) Each minimization -- over $U$ with $V$ fixed, and vice versa -- is a simple least-squares problem, which can be solved in time $O(dn^2k^2+n^3k^3)$\footnote{Throughout this paper, we assume $m\leq n$.}, b) We initialize $U^0$ to be the top-$k$ left singular vectors of  $\sum_i A_i b_i$ (step 2 of Algorithm~\ref{algo:sensing-altmin}). As we will see later in Section~\ref{sec:sensing}, this provides a good initialization point for the sensing problem which is crucial; if the first iterate $\widehat{U}^0$ is orthogonal, or almost orthogonal, to the true $U^*$ subspace, \as  may never converge to the true space (this is easy to see in the simplest case, when the map is identity, i.e. $\aff(X) = X$ -- in which case \as just becomes the power method). %Section \ref{sec:sensing} has further comments on the algorithm.

%We now state the algorithm.%Each minimization -- over $U$ with $V$ fixed, and vice versa -- is a simple least-squares problem, which can be solved in {\bf \tt @@complexity} time. 
%In particular, note that if we hold $U$ fixed, the minimization over $V$ {\em is} convex; in fact it is simply a least-squares problem. The same is true if we fix $V$ and minimize over $U$. One can thus alternate between minimizing one and then the other; this (plus an initializing step) is our alternating minimization algorithm, formally specified below.

\begin{algorithm}[th]
\caption{{\bf \as:} Alternating minimization for matrix sensing}
\label{algo:sensing-altmin}
\begin{algorithmic}[1]
\STATE Input $b,\aff$
\STATE Initialize $\widehat{U}^0$ to be the top-$k$ left singular vectors of $\sum_i A_i b_i$
\FOR{$t=0,\cdots,T-1$}
	\STATE $\widehat{V}^{t+1}\leftarrow \argmin_{V\in \R^{n\times k}} \  \|\aff(\widehat{U}^t \,V^\dag )-b\|_2^2$
%	\STATE $V^{t+1}\leftarrow QR(\widehat{V}^{t+1})\ \ \text{ (Orthogonalization using QR decomposition) }$
	\STATE $\widehat{U}^{t+1}\leftarrow \argmin_{U\in \mathbb{R}^{m\times k}} \  \|\aff(U\,(\widehat{V}^{t+1})^\dag )-b\|_2^2$
%        \STATE $U^{t+1}\leftarrow QR(\widehat{U}^{t+1})\ \ \text{ (Orthogonalization using QR decomposition) }$
%        \STATE $\Sigma^{t+1}\leftarrow \argmin_{\Sigma\in\mathbb{R}^{k\times k}} \|\aff(U^{t+1}\Sigma(V^{t+1})^\dag )-b\|_2^2$
\ENDFOR
\STATE Return $X=\widehat{U}^T(\widehat{V}^T)^\dag $
\end{algorithmic}
\end{algorithm}

%{\bf Comments on algorithm:} {\em (a)} Each minimization -- over $U$ with $V$ fixed, and vice versa -- is a simple least-squares problem, which can be solved in {\bf \tt @@complexity} time. {\em (b)} We require the top-$k$ singular vectors of the $m\times n$ matrix $\sum_i A_i b_i$ for initialization. This is because if the first iterate $\widehat{U}^0$ starts off orthogonal, or almost orthogonal, to the true $U^*$ subspace, it may may never converge to the true space (this is easy to see in the simplest case, when the map is identity, i.e. $\aff(X) = X$ -- in which case our method just becomes the power method). Section \ref{sec:sensing} has further comments on the algorithm.

%We now analyze the performance of \ase. 
In general, since $d<mn$, problem \eqref{eq:lrms} is not well posed as there can be multiple rank-$k$ solutions that satisfy $\calA(X)=b$. However, inspired by a similar condition in compressed sensing \cite{CandesT2005}%\footnote{In particular, this matrix RIP when specialized to diagonal $X$ is exactly the RIP property of sparse vector compressed sensing \cite{CandesT2005}.}
, Recht et al. \cite{RechtFP2007} showed that if the linear map $\aff$ satisfies a {\em (matrix) restricted isometry property} (RIP), then a trace-norm based convex relaxation of \eqref{eq:lrms} leads to exact recovery. This property is defined below. 
\begin{definition}\cite{RechtFP2007}\label{defn:rip}
A linear operator $\aff(\cdot):\rmn\rightarrow \mathbb{R}^d$ is said to satisfy $k$-RIP, with $\delta_{k}$ RIP constant, if for all $X\in \rmn$ s.t. rank$(X)\leq k$, the following holds:
\begin{align}\label{eqn:sensing-2k-RIP1}
(1-\delta_k) \left\|X\right\|_F^2 \leq \left\|\calA(X)\right\|_2^2 \leq (1+\delta_k) \left\|X\right\|_F^2.
\end{align}
\end{definition}
Several random matrix ensembles with sufficiently many measurements ($d$) satisfy  matrix RIP \cite{RechtFP2007}. For example, if $d=\Omega(\frac{1}{\delta_k^2}kn\log n)$ and each entry of $A_i$ is sampled i.i.d. from a $0$-mean sub-Gaussian distribution then $k$-RIP is satisfied with RIP constant  $\delta_k$. 

We now present our main result for \ase.
\begin{theorem}\label{thm:sensing_err}
 Let $M=U^*\Sigma^*V^{*^\dag }$ be a rank-$k$ matrix with non zero singular values $\sigma_1^*\geq \sigma_2^*\dots\geq \sigma_k^*$. Also, let the linear measurement operator $\calA(\cdot):\rmn \rightarrow \mathbb{R}^d$ satisfy $2k$-RIP with RIP constant $\delta_{2k}<\frac{(\so_k)^2}{(\so_1)^2}\,\frac{1}{100k}$.  Then, in the \as algorithm (Algorithm \ref{algo:sensing-altmin}), for all $T> 2\log (\|M\|_F/\epsilon)$,  the iterates $\widehat{U}^{T}$ and $\widehat{V}^T$ satisfy:%\footnote{Here $\|\cdot\|_F^2$ represents the Frobenius norm of the matrix, i.e. the square root of the sum of squares of all its elements.}: 
 $$\|M-\widehat{U}^{T}(\widehat{V}^T)^\dag\|_F\leq \epsilon.$$
%where $\epsilon>0$ is any constant. 
%Then, the $(t+1)$-th iterates $U^{t+1}$, $V^{t+1}$ of the AltMin algorithm  satisfy: 
%with $\dist\left(U^{(0)},U^*\right) \leq \frac{1}{2}$. Then $U^{(t)}$ and $V^{(t)}$ satisfy:
%\begin{align*}
%  \dist\left(V^{t+1},V^*\right) &\leq \alpha_V\cdot\dist\left(U^{t}, U^*\right)  \mbox{ and} \\
%  \dist\left(U^{t+1},U^*\right) &\leq \alpha_U\cdot\dist\left(V^{t+1}, V^*\right)
%\end{align*}
%where $\alpha_V, \alpha_U < 1$ are some constants depending only on $\delta_{2k}$. $\dist\left(U, W\right)$ denotes the principal angle based distance between subspaces spanned by the columns of $U\in \mathbb{R}^{m\times k}$ and $W\in \mathbb{R}^{m\times k}$ (See Definition~\ref{defn:dist}). 
\end{theorem}
The above theorem establishes geometric convergence (in $O(\log(1/\epsilon))$ steps) of \as to the optimal solution of \eqref{eq:lrms} under standard RIP assumptions. %Furthermore, we demonstrate geometric convergence ($O(\log (1/\epsilon))$ steps) to the optimal solution. 
This is in contrast to existing iterative methods for trace-norm minimization all of which require at least $O(\frac{1}{\sqrt{\epsilon}})$ steps; interior point methods for trace-norm minimization converge to the optimum in $O(\log(1/\epsilon))$ steps but require storage of the full $m\times n$ matrix and require $O(n^5)$ time per step, which makes it infeasible for even moderate sized problems.

Recently, several projected gradient based methods have been developed for matrix sensing \cite{JainMD10, LeeB10} that also guarantee convergence to the optimum in $O(\log (1/\epsilon))$ steps. But each iteration in these algorithms requires computation of the top $k$ singular components of an $m\times n$ matrix, which is typically significantly slower  than solving a least squares problem (as required by each iteration of \ase).\\
{\bf Stagewise AltMinSense Algorithm}: A drawback of our analysis for \as is the dependence of $\delta_{2k}$ on the condition number ($\kappa=\frac{\so_1}{\so_k}$) of $M$, which implies that the number of measurements $d$ required by \as grows quadratically with $\kappa$. We address this issue by using a stagewise version of \as (Algorithm~\ref{algo:sensing-altmin-improved}) for which we are able to obtain near optimal measurement requirement. 

The key idea behind our stagewise algorithm is that if one of the singular vectors of $M$ is very dominant, then we can treat the underlying matrix as a rank-$1$ matrix plus noise and approximately recover the top singular vector. Once we remove this singular vector from the measurements, we will have a relatively well-conditioned problem. Hence, at each stage of Algorithm~\ref{algo:sensing-altmin-improved}, we seek to remove the remaining most dominant singular vector of $M$. See Section~\ref{sec:sam} for more details; here we state the corresponding theorem regarding the performance of Stage-AltMin. %In Section \ref{sec:sam}, we provide a pseudocode of stage wise modification of \as that removes this dependence. The algorithm is formally stated in that section; here we state the corresponding theorem regarding its performance.
\begin{theorem}\label{thm:sam_err}
Let $M=U^*\Sigma^*V^{*^\dag }$ be a rank-$k$ incoherent matrix with non zero singular values $\sigma_1^*\geq \sigma_2^*\dots\geq \sigma_k^*$. Also, let $\calA(\cdot):\rmn \rightarrow \mathbb{R}^d$ be a linear measurement operator that satisfies $2k$-RIP with RIP constant $\delta_{2k}<\frac{1}{3200k^2}$.  Suppose, Stage-AltMin (Algorithm \ref{algo:sensing-altmin-improved}) is supplied inputs $\aff$, $b=\aff(M)$. Then, the $i$-th stage iterates $\Uw^T_{1:i}$, $V^T_{1:i}$ satisfy: $$\|M-\Uw^T_{1:i}\left(V^T_{1:i}\right)^{\dag}\|_F^2\leq \max(\epsilon, 16k(\so_{i+1})^2),$$% + (\so_i)^2\left(\frac{1}{4}\right)^t.$$
%That is, the $T$-th step iterates of the $k$-th stage, satisfy: $$\|M-\Uw^T_{1:k}V^T_{1:k}\|_F^2\leq \epsilon,$$
where $T= \Omega\left(\log(\|M\|_F^2/\epsilon)\right)$. That is, the $T$-th step iterates of the $k$-th stage, satisfy: $\|M-\Uw^T_{1:k}\left(V^T_{1:k}\right)^{\dag}\|_F^2\leq \epsilon.$
\end{theorem}
The above theorem guarantees exact recovery using $O(k^4 n \log n)$ measurements which is only $O(k^3)$ worse than the information theoretic lower bound. We also note that for simplicity of analysis, we did not optimize the constant factors in $\delta_{2k}$.% which should improve with a little more careful analysis.
%{\bf Comments on the theorem:}
%{\em (a)} We establish the success of alternating minimization under assumptions that are typical of work on this problem \cite{RechtFP2007,LeeB2009,JainMD10}.
%Our theorem shows geometric convergence of the iterates to the true matrix $M$. 
%{\bf \tt @@ add comments here. 0. that alt min works once we make assumptions similar to what is reqd for other methods to solve under-determined case. 1. that we get $\epsilon$ error within $\log 1/\epsilon$ iterations. compare to other methods for low-rank, - e.g. that iterative methods to solve convex programs need many more SVDs (write precise scaling) to get within $\epsilon$, and also cite the same. 2. dependence on condition number, comparison of values of $\delta_{2k}$ to those needed by other algorithms.}
%In the following, we first present a simple proof for the rank-1 case, and then for the general rank-$k$ case.
%{\bf \tt @@ also finish writing about stage wise. leave this algorithm in that section, though.}
\subsection*{Matrix Completion via Alternating Minimization}
The matrix completion problem is the following: there is an unknown rank-$k$ matrix $M\in \mathbb{R}^{m\times n}$, of which we know a set $\Omega \subset [m]\times [n]$ of elements; that is, we know the values of elements $M_{ij}$, for $(i,j)\in\Omega$. The task is to recover $M$. Formally, the Low-Rank Matrix Completion problem is:\renewcommand{\theequation}{LRMC}
\begin{equation}
\mbox{Find rank-}k \mbox{ matrix } X \mbox{ s.t. }  P_\Omega(X)=P_\Omega(M),\label{eq:matcomp}
 \end{equation} \renewcommand{\theequation}{\arabic{equation}} \addtocounter{equation}{-1}
where for any matrix $S$ and a set of elements $\Omega \subset [m]\times [n]$ the matrix $P_\Omega(S)\in \mathbb{R}^{m\times n}$ is as defined below:
\begin{equation}
  P_\Omega(S)_{ij}=\begin{cases}S_{ij}&\mbox{ if }(i,j)\in \Omega,
\\ 0&\mbox{ otherwise.}
\end{cases}\label{eq:pomega}
\end{equation}
We are again interested in the under-determined case; in fact, for a fixed rank $k$, as few as $O(n\log n)$ elements may be observed. This problem is a special case of matrix sensing, with the measurement matrices $A_i=e_je_\ell^\dag$ being non-zero only in single elements; however, such matrices do not satisfy  matrix RIP conditions like (\ref{eqn:sensing-2k-RIP1}). For example, consider a low-rank $M=e_1e_1^\dag$ for which a uniformly random $\Omega$ of size $O(n\log n)$ will most likely miss the non-zero entry of $M$.

Nevertheless, like matrix sensing, matrix completion has been shown to be possible once additional conditions are applied to the low-rank matrix $M$ and the observation set $\Omega$. Starting with the first work \cite{CandesR2007}, the typical assumption has been to have $\Omega$ generated uniformly at random, and $M$ to satisfy a particular incoherence property that, loosely speaking, makes it very far from a sparse matrix. In this paper, we show that {\em once} such assumptions are made, alternating minimization {\em also} succeeds. We now restate, and subsequently use, this incoherence definition.

\begin{definition}\cite{CandesR2007}\label{defn:sampling-incoherentmatrices}
A matrix $M\in \rmn$ is incoherent with parameter $\mu$ if:%  An $n\times k$ matrix $V \in \R^{n\times k}$ is said to be incoherent with parameter $\mu$ if we have:
\begin{align}\label{eqn:sampling-incoherentmatrices}
  \twonorm{u^{(i)}} \leq \frac{\mu \sqrt{k}}{\sqrt{m}} \;\forall\;i \in [m],\  \twonorm{v^{(j)}} \leq \frac{\mu \sqrt{k}}{\sqrt{n}} \;\forall\;j \in [n],% \ \ \ \ \twonorm{u^{(j)}} \leq \frac{\mu \sqrt{k}}{\sqrt{m}} \;\;\;\;\forall\;\;j \in [n]
\end{align}
where $M=U\Sigma V^T$ is the SVD of $M$ and $u^{(i)}$, $v^{(j)}$ denote the $i^{\textrm{th}}$ row of $U$
and the $j^{\textrm{th}}$ row of $V$ respectively.
\end{definition}
The alternating minimization algorithm can be viewed as an approximate way to solve the following non-convex problem:
\[
\min_{U,V\in \mathbb{R}^{n\times k}}  \quad  \| P_\Omega(UV^\dag) - P_\Omega(M)\|^2_F 
\]
Similar to \ase, the altmin procedure proceeds by alternatively solving for $U$ and $V$.  As noted earlier, this approach has been popular in practice and has seen several variants and extensions being used in practice \cite{Zhouetal08,KorenBV09,Koren09,chen2012matrix}. However, for ease of analysis, our algorithm further modifies the standard alternating minimization method. In particular, we introduce  partitioning of the observed set $\Omega$, so that we use different partitions of $\Omega$ in each iteration. See Algorithm~\ref{algo:completion-altmin} for a pseudo-code of our variant of the alternating minimization approach.
%One could now proceed along the lines of the \as algorithm above; indeed doing so may provide good results empirically \cite{Zhouetal08,KorenBV09,Koren09,chen2012matrix}{\bf \tt @@cite various papers by yehuda, recht etc.} 
%However, proving the success of alternating minimization in the matrix completion setup (with incoherence and random $\Omega$) is even more complicated than the sensing problem. We hence develop a modified algorithm, where the main innovation is a partitioning and stage-wise use of the observed set $\Omega$, so that we use different partitions in each iteration.
\begin{algorithm}[th]
\caption{{\bf AltMinComplete:} Alternating minimization for matrix completion}
\label{algo:completion-altmin}
\begin{algorithmic}[1]
\STATE Input: observed set $\Omega$, values $P_{\Omega}(M)$
\STATE Partition $\Omega$ into $2T+1$ subsets $\Omega_0,\cdots,\Omega_{2T}$ with each element of $\Omega$ belonging to one of the $\Omega_t$ with equal probability (sampling with replacement)
\STATE $\widehat{U}^{0}=SVD(\frac{1}{p}P_{\Omega_0}(M), k)$ i.e., top-$k$ left singular vectors of $\frac{1}{p}P_{\Omega_0}(M)$
\STATE Clipping step : Set all elements of $\widehat{U}^{0}$ that have magnitude greater than $\frac{2\mu \sqrt{k}}{\sqrt{n}}$ to zero and orthonormalize the columns of $\widehat{U}^0$
\FOR{$t=0,\cdots,T-1$}
	\STATE $\widehat{V}^{t+1}\leftarrow \argmin_{V\in \R^{n\times k}} \  \|P_{\Omega_{t+1}}(\widehat{U}^t V^\dag - M)\|_F^2$
%	\STATE $V^{t+1}\leftarrow QR(\widehat{V}^{t+1})$ (QR decomposition)
	\STATE $\widehat{U}^{t+1}\leftarrow \argmin_{U\in \mathbb{R}^{m\times k}} \  \|P_{\Omega_{T+t+1}}(U\left(\widehat{V}^{t+1}\right)^\dag - M) \|_F^2$
%        \STATE $U^{t+1}\leftarrow QR(\widehat{U}^{t+1})$ (QR decomposition)
%         \STATE $\Sigma^{t+1}\leftarrow \argmin_{\Sigma\in\mathbb{R}^{k\times k}} \|\aff(U^{t+1}\Sigma(V^{t+1})^\dag )-b\|_2^2$
\ENDFOR
\STATE Return $X=\widehat{U}^T(\widehat{V}^T)^\dag $
\end{algorithmic}
\end{algorithm}

We now present our main result for \eqref{eq:matcomp}:
%The following theorem is an analogue of Theorem \ref{thm:sensing-altmin-geomconv} for the case of matrix completion.
\begin{theorem}\label{thm:sampling-altmin-main}
Let $M=U^*\Sigma^*V^{*^\dag } \in \mathbb{R}^{m\times n}$ ($n\geq m$) be a rank-$k$ incoherent matrix, i.e., both $U^*$ and $V^*$ are $\mu$-incoherent (see Definition~\ref{defn:sampling-incoherentmatrices}).  Also, let each entry of $M$ be observed uniformly and independently with probability, 
$$p>C\frac{\left(\frac{\sigma_1^*}{\sigma_k^*}\right)^2\mu^2k^{2.5}\log n \log \frac{k\|M\|_F}{\epsilon}}{m\delta_{2k}^2},$$
where $\delta_{2k}\leq \frac{\so_k}{12k\so_1}$ and $C>0$ is a global constant. 
%$$p>C\frac{\left(\frac{\sigma_1^*}{\sigma_k^*}\right)^4\mu^2k^{4.5}\log \frac{n}{\epsilon}}{m}$$
%(assume $n\geq m$) for a global constant $C>0$. Let $\Omega$ be the indices of observed entries of $M$. 
 Then w.h.p. for $T = C'\log\frac{\frob{M}}{\epsilon}$, the outputs $\widehat{U}^{T}$ and $V^{T}$ of  Algorithm \ref{algo:completion-altmin}, with input $(\Omega, P_\Omega(M))$ (see Equation \eqref{eq:pomega}) satisfy: 
%Then, outputs $U^{t}$ and $V^{t}, \forall t\geq 1$ of  Algorithm \ref{algo:sensing-altmin}, with input $b=\calA\left(M\right)$ and $\aff$, satisfy: %and $U^{0}$ with $\dist\left(U^{0},U^*\right) \leq \frac{1}{2}$, then $U^{t}$ and $V^{t}, \forall t\geq 1$ satisfy:
%\begin{align*}
$  \frob{M - \widehat{U}^{T}\left(V^{T}\right)^{\dag}} \leq \epsilon.$
%   \dist\left(U^{t},U^*\right) &\leq \left(\frac{1}{4}\right)^t \mbox{ , }   \dist\left(V^{t},V^*\right) \leq \left(\frac{1}{4}\right)^t,
%\end{align*}
% where $\dist(\cdot, \cdot)$ is the maximum principle angle based distance between subspaces spanned by the columns of $U\in \mathbb{R}^{m\times k}$ and $W\in \mathbb{R}^{m\times k}$ (See Definition~\ref{defn:dist}).
%where $c_1$ and $c_2 < 1$ are some constants depending only on $\delta$.
\end{theorem}
The above theorem implies that by observing $\left|\Omega\right|=O\left((\frac{\so_1}{\so_k})^4k^{4.5}n \log n\log(k\|M\|_F/\epsilon)\right)$ random entries of an incoherent $M$, \amc can recover $M$ in $O(\log(1/\epsilon))$ steps. In terms of sample complexity ($|\Omega|$), our results show alternating minimization may require a bigger $\Omega$ than convex optimization, as our result has $|\Omega|$ depend on the condition number, required accuracy ($\epsilon$) and worse dependence on $k$ than known  bounds. In contrast, trace-norm minimization based methods require $O(kn \log n)$ samples only. 

Empirically however, this is not seen to be the case \cite{JainMD10} and we leave further tightening of the sample complexity bounds for matrix completion as an open problem. 

In terms of time complexity, we show that \amc needs  time $O(|\Omega| k^2 \log (1/\epsilon))$. This is in contrast to popular trace-norm minimization based methods that need $O(1/\sqrt{\epsilon})$ steps \cite{CaiCS2008} and total time complexity of $O(|\Omega|n/\sqrt{\epsilon})$; note that the latter can be potentially quadratic in  $n$. Furthermore, each step of such methods requires computation of the SVD of an $m\times n$ matrix. As mentioned earlier, interior point methods for trace-norm minimization also converge in $O(\log (1/\epsilon))$ steps but each iteration requires $O(n^5)$ steps and need storage of the entire $m\times n$ matrix $X$. 


%Another appealing aspect of \amc is the ability to parallelize it across elements. In particular, consider for example the minimization w.r.t. $U$ for a fixed $\widehat{V}$. Each row $U^{(i)} \in \mathbb{R}^{1\times k}$ of $U$ can now be minimized in parallel, independent of other rows. This has lead to highly scalable algorithms for recommendation systems \cite{Zhouetal08,Koren09}.  %furthermore, this minimization depends only on those rows $j$ of $\widehat{V}$ for which the element $(i,j)$ is observed, i.e. $(i,j)\in \Omega$. This is another reason for the empirical interest in this algorithm \cite{Zhouetal08,Koren09}.% that the minimization steps naturally decompose
%{\bf \tt @@comments etc as in matrix sensing}
%
% \subsection*{Further Results for Matrix Sensing}
%
% We now improve our result on alternating minimization for matrix sensing on two fronts: extending to the case when we have noisy measurements, and a new algorithm -- StagewiseAltMin -- that removes the dependence of the RIP constant on the condition number of the matrix to be recovered. 
%
% {\bf \tt @@ put in theorem on noisy matrix sensing in terms of Err < eps}
%
% {\bf Stagewise algorithm and condition number:} In Theorem \ref{thm:sensing_err} above, the value of the matrix RIP constant $\delta_{2k}$ required for success depended on the condition number $\frac{(\so_k)}{(\so_1)}$ of the matrix $M$ to be recovered. Such a dependence does not exist for other algorithms based on convex optimization \cite{RechtFP2007} or singular value projection \cite{JainMD10}. In Section \ref{sec:sam}, we provide a stage wise modification of \as that removes this dependence. The algorithm is formally stated in that section; here we state the corresponding theorem regarding its performance.
% \begin{theorem}\label{thm:sam_err}
% Let $M=U^*\Sigma^*V^{*^\dag }$ be a rank-$k$ incoherent matrix with non zero singular values $\sigma_1^*\geq \sigma_2^*\dots\geq \sigma_k^*$. Also, let $\calA(\cdot):\rmn \rightarrow \mathbb{R}^d$ be a linear measurement operator that satisfies $2k$-RIP with RIP constant $\delta_{2k}<\frac{1}{3200k^2}$.  Suppose, Stage-AltMin (Algorithm \ref{algo:sensing-altmin-improved}) is supplied inputs $\aff$, $b=\aff(M)$. Then, the $i$-th stage iterates $\Uw^T_{1:i}$, $V^T_{1:i}$ satisfy: $$\|M-\Uw^T_{1:i}\left(V^T_{1:i}\right)^{\dag}\|_F^2\leq \max(\epsilon, 16k(\so_{i+1})^2),$$% + (\so_i)^2\left(\frac{1}{4}\right)^t.$$
% %That is, the $T$-th step iterates of the $k$-th stage, satisfy: $$\|M-\Uw^T_{1:k}V^T_{1:k}\|_F^2\leq \epsilon,$$
% where $T= \Omega\left(\log(\|M\|_F^2/\epsilon)\right)$. That is, the $T$-th step iterates of the $k$-th stage, satisfy: $$\|M-\Uw^T_{1:k}\left(V^T_{1:k}\right)^{\dag}\|_F^2\leq \epsilon.$$
% \end{theorem}
%
% {\bf Comments on the theorem:}
% {\em (a)}
%
% {\bf \tt @@ finish}
%
%
%%% Local Variables: 
%%% mode: latex
%%% TeX-master: "altmin"
%%% End: 

%where, linear measurements $b$ are obtained by measuring a rank-$k$ matrix $M\in \rmn$, i.e., $b=\mathcal{A}(M).$ In several practical applications, measurements can be noisy in which case $b$ is given by $b= \mathcal{A}(M) +\eta$, where $\eta$ is bounded noise. 
%%% Local Variables: 
%%% mode: latex
%%% TeX-master: "altmin"
%%% End: 

\section{Related Work}\label{sec:related}
%Low-rank recovery, and matrix completion in particular, by alternating minimization have been the subject of much empirical effort; ours is the first analysis of when this succeeds. We now summarize three lines of related work.
{\bf Alternating Minimization}: Alternating minimization and its variants have been applied to several low-rank matrix estimation problems. For example, clustering \cite{KimPark08b}, sparse PCA \cite{ZouHT06}, non-negative matrix factorization \cite{KimPar08}, signed network prediction \cite{HsiehCD12} etc. There are three main reasons for such wide applicability of this approach: a) low-memory footprint and fast iterations, b) flexible modeling, c) amenable to parallelization. However, despite such empirical success, this approach has largely been used as a heuristic and has had no theoretical analysis other than the guarantees of convergence to the {\em local minima} \cite{Zangwill69}. Ours is the first analysis of this approach for two practically important problems: a) matrix completion, b) matrix sensing. 

\noindent {\bf Matrix Completion:} This is the problem of completing a low-rank matrix from a few sampled entries. Candes and Recht  \cite{CandesR2007} provided the first results on this problem, showing that under the random sampling and incoherence conditions (detailed above), $O(kn^{1.2}\log n)$ samples allow for recovery via convex trace-norm minimization; this was improved to $O(kn\log n)$ in \cite{CandesT2009}. For large matrices, this approach is not very attractive due to the need to store and update the entire matrix, and because iterative methods for trace norm minimization require  $O(\frac{1}{\sqrt{\epsilon}})$ steps to achieve additive error of $\epsilon$. Moreover, each such step needs to compute an SVD. 

Another  approach, in \cite{KeshavanOM2009}, involved taking a single SVD, followed by gradient descent  on a Grassmanian manifold. However,  {\em (a)} this is more expensive than alternating minimization as it needs to compute gradient over Grassmanian manifold which in general is a computationally intensive step, and {\em (b)} the analysis of the algorithm only guarantees asymptotic convergence, and in the worst case might take exponential time in the problem size. 

Recently, several other matrix completion type of problems have been studied in the literature. For example, robust PCA \cite{ChandrasekaranSPW11,CandesLMW11}, spectral clustering \cite{JalaliCSX11} etc. Here again, under additional assumptions, convex relaxation based methods have rigorous analysis but alternating minimization based algorithms continue to be  algorithms of choice in practice. 
 
\noindent {\bf Matrix Sensing:} The general problem of matrix sensing was first proposed by \cite{RechtFP2007}. They  established recovery via trace norm minimization, assuming the sensing operator satisfies ``restricted isometry" conditions. Subsequently, several other methods \cite{JainMD10, LeeB10}  were proposed for this problem that also recovers the underlying matrix with optimal number of measurements and can give an $\epsilon$-additive approximation in time $O(\log(1/\epsilon)$. But, similar to matrix completion, most of  these methods require computing SVD of a large matrix at each step and hence have poor scalability to large problems.  %{\bf \tt @@ add more refs, and the deltas they require}

We show that \as and AltMin-Completion provide more scalable algorithms for their respective problems. We  demonstrate that these algorithms have geometric convergence to the optima, while each iteration is relatively cheap. For this, we assume conditions similar to those required by existing algorithms; albeit, with one drawback: number of samples required by our analysis depend on the condition number of the underlying matrix $M$. For the matrix sensing problem, we remove this requirement by using a stagewise algorithm; we leave similar analysis for matrix completion as an open problem.%There is no other analysis (though plenty of empirical work) on alternating minimization for these problems.

%{\bf \tt @@ this para may need some change} Apart from explaining widespread success of such methods for low-rank matrix completion and sensing problems, our method also hints at two important aspects to take care off, while designing an alternating minimization method: a) good initialization point, b) stagewise solution of the problem so that even if the initial point is heavily biased towards the dominant components of the true rank-$k$ subspace, later stages can remove the dominant components and can recover the smaller components of the true subspace. 

%
%Estimating low-rank matrices that satisfy given  constraints is an important problem with several applications in recommendation systems \cite{BellK2007}, control theory \cite{FazelHB2001}, machine learning \cite{MekaJCD2008} etc. However, the problem is NP-hard in general and is even hard to approximate within a logarithmic factor \cite{MekaJCD2008}. 
%
%There are several cases where the problem is well studied and has provably efficient solutions \cite{RechtFP2007, CandesR2007, AroraMK2011}. One of the most well-known example here is low-rank approximation of a given matrix. It is well-known that the SVD of a given matrix is the best low-rank approximation and there exists several efficient algorithms for this task \cite{GolVan96}. 
%
%However, the problem becomes significantly harder if rather than the entire matrix, only a subset of its entries are provided. This problem is called low-rank matrix completion problem and has shown to be critical for several practical tasks such as recommender systems, 2-d map reconstruction \cite{CandesR2007} etc. For this problem, first positive results were by \cite{CandesR2007} who showed that if the matrix is not ``spiky'', i.e., is incoherent (see Definition~\ref{defn:incoherence}), then it can be recovered back from $O(kn^{1.2}\log n)$ randomly sampled entries using a trace-norm based method. Later, \cite{CandesT2009, Recht2009} improved upon this result and brought down the sample complexity to information theoretic optimal $O(kn\log n)$. However, this method does not scale well with problem size and as such is not very popular in practice as it requires to solve a trace-norm minimization for which convergence rates are typically only $O(\frac{1}{\sqrt{\epsilon}})$ to achieve additive error of $\epsilon$. Moreover, each step of the optimization algorithm needs to compute a SVD of the matrix which can be prohibitively expensive.  
%
%
%In an another line of work, \cite{RechtFP2007} studied the matrix-sensing problem where the goal is to recover a low-rank matrix using a small number of (possibly dense) measurements. Here again, they show that using $O(kn\log n)$ measurements and a trace-norm minimization method, a rank-$k$ matrix can be recovered exactly. 
%
%Apart from the above two mentioned problems, several other popular low-rank estimation problems have been studied. For example, non-negative matrix factoriazation, low-rank kernel learning, general affine constrained low-rank estimation problem, low-rank plus sparse decomposition problem etc. For most of these problems, only a few provably efficient algorithms are known and most of these algorithms some form of convex relaxation which in practice makes algorithm less scalable  to large problem sizes. 
%
%For most of the above mentioned low-rank estimation problems, one method that has consistently been proven to be accurate and robust in practice is the method of alternating minimization. This method decomposes a low-rank matrix as $M=UV^\dag$ and then does co-ordinate descent (minimization) over $U$ and $V$ alternately. Typically, the problem obtained by fixing $U$ (or $V$) is a strongly convex problem and can be solved efficiently. For example, for matrix completion or sensing problems, the resulting problem is a least sqaures problem that can be solved efficiently using conjugate gradient type of methods. Furthermore, typically such methods avoid SVD style computations which in practice are much more expensive then least squares compuation. However, decomposing $M$ in this bi-linear form results in a non-convex problem which cannot be analyzed in general. In fact, to the best of our knowledge, only analysis for such methods is for the singular value decomposition problem, where such alternating mimization method essentially reduces to the well-known power-method.
%
%In this work, we provide first analysis of the above mentioned alternating mimization method for the matrix completion and sensing problems. For each of the problem, we show that alternating minimization when supplied with a reasonably good initial point leads to a geoemtric convergence to the optimal solution, i.e., $\epsilon$-additive error can be obtained in just $O(\log (1/\epsilon)$ iterations. However, our analysis requires measurements that scale up with higher condition no. of the underlying matrix. We show that we can alleviate this issue by using a stagewise alternating minimization procedure where in the $i$-th stage we solve a rank-$i$ problem. This analysis leads to recovery guarantees for sensing problem with near optimal measurement complexity. For the matrix completion problem, while alternating mimization would still lead to recovery with near optimal measureents, the ``goodness'' of intialization method that we use is an open problem which leave as a topic of future research. 

%%% Local Variables: 
%%% mode: latex
%%% TeX-master: "altmin"
%%% End: 

\section{Matrix Sensing}\label{sec:sensing}
In this section, we  study the matrix sensing problem \eqref{eq:lrms} and prove that if the measurement operator, $\aff$, satisfies RIP then \as (Algorithm \ref{algo:sensing-altmin}) recovers the underlying low-rank matrix {\em exactly} (see Theorem~\ref{thm:sensing_err}).

%Below we recall  the robust version of \eqref{eq:lrms}: 
%\begin{align}
%  \min_{X}&\ \  \|\aff(X)-b\|_2^2,\nonumber\\
%{\text s.t.}&\ \  X=\widehat{U} \widehat{V}^\dag , \widehat{U}\in \mathbb{R}^{m\times k},  \widehat{V}\in \mathbb{R}^{n\times k}, 
%\label{eq:altmin}
%\end{align}
%where $\aff:\rmn \rightarrow \mathbb{R}^d$ is a linear operator and $b = \aff\left(M\right)$ is the vector of measurements corresponding to an underlying low rank matrix $M$.
% WLOG, we can assume that $b_i=\Tr(A_iM), \forall i$ where $A_i$ are the measurement matrices corresponding to $\aff\left(\cdot\right)$.

%Let $M=U^*\Sigma^*V^{*T}$ be the SVD of $M$. $U^*\in \mathbb{R}^{m\times k}$, $V^*\in \mathbb{R}^{n\times k}$ are orthonormal singular vectors and $\Sigma^*\in \mathbb{R}^{k\times k}$ is a diagonal matrix with $\sigma_i^*$ being the $i$-th singular value. Also, let $\sigma_1^*\geq \sigma_2^*\dots\geq \sigma_k^*$. 

%We now describe the standard alternating minimization algorithm (AltMin), but with initial iterate selected using one step of the Singular Value Projection (SVP) method for \eqref{eq:lrms} \cite{JainMD10}. 

% % % % % % Recall that, each step of the AltMin algorithm maintains iterates $\widehat{U}^t, \widehat{V}^t$ ($t\geq 0$). Let $\widehat{U}^t$ be the iterate after the $t$-th step, then $\Vht$ is obtained by solving the least squares problem obtained by fixing $U$ to be $\Uht$  in \eqref{eq:altmin}. That is, $$\widehat{V}^{t+1}=\argmin_{\widehat{V}} \  \|\aff(\Uht\widehat{V}^\dag )-b\|_2^2,\ \  \widehat{V}^{t+1}=V^{t+1}R_V^{t+1}.$$%%\ \ \text{ (QR decomposition) }.$$
% % % % % % Similarly, $\widehat{U}^{t+1}$ is obtained by fixing $V$ to be $\Vht$:% and finally, $\Sigma^{t+1}$ is obtained by fixing both $U=U^{t+1}$ and $V=V^{t+1}$, i.e., 
% % % % % %  $$\widehat{U}^{t+1}=\argmin_{\widehat{U}} \  \|\aff(\widehat{U}(\Vht)^\dag )-b\|_2^2,\ \  \widehat{U}^{t+1}=U^{t+1}R_U^{t+1}.$$%%\ \ \text{ (QR decomposition) }.$$
% % % % % % Note that unlike standard alternating minimization algorithms, we compute QR decomposition of both $\widehat{V}^{t+1}$ and $\widehat{U}^{t+1}$. QR decomposition of both $U$ and $V$ can be computed in time $O(k^2 n)$. Furthermore, we stress that our analysis  holds even without the QR decomposition step; introduction of QR step helps us to simplify notations and proofs. 

% $$\Sigma^{t+1}=\argmin_{\Sigma} \  \|\aff(U^{t+1}\Sigma V^{t+1})-b\|_2^2.$$
%See Algorithm~\ref{algo:sensing-altmin} for a pseudo code of the AltMin algorithm for matrix sensing. Note that only the first step in the algorithm requires finding SVD; all the other steps involves solving a least squares problem. Also, 
% \begin{algorithm}[th]
% \caption{Alternating minimization for matrix sensing}
% \label{algo:sensing-altmin}
% \begin{algorithmic}[1]
% \STATE Input $b,\aff$
% \STATE $U^0=SVD(\aff^\dag  b, k)$ i.e., top-$k$ left singular vectors of $\aff^\dag  b=\sum_i A_i b_i$
% \FOR{$t=0,\cdots,T-1$}
% 	\STATE $\widehat{V}^{t+1}\leftarrow \argmin_{\widehat{V}\in \R^{n\times k}} \  \|\aff(\Ut\widehat{V}^\dag )-b\|_2^2$
% 	\STATE $V^{t+1}\leftarrow QR(\widehat{V}^{t+1})\ \ \text{ (QR decomposition) }$
% 	\STATE $\widehat{U}^{t+1}\leftarrow \argmin_{\widehat{U}\in \mathbb{R}^{m\times k}} \  \|\aff(\widehat{U}(V^{t+1})^\dag )-b\|_2^2$
%         \STATE $U^{t+1}\leftarrow QR(\widehat{U}^{t+1})\ \ \text{ (QR decomposition) }$
% \ENDFOR
% \STATE Return $X=\widehat{U}^{T}(V^{T})^\dag $
% \end{algorithmic}
% \end{algorithm}

%As mentioned earlier, although the alternating minimization algorithm is popular in practice, proof of its convergence to the global optimum of \eqref{eq:altmin} is an open problem.
%The following theorem shows that if $\aff(\cdot)$ satisfies RIP then AltMinSense (Algorithm \ref{algo:sensing-altmin})
%converges to the global optimum of \eqref{eq:altmin} at linear rate. 

% % % % % % \noindent Now, we present our result for AltMin algorithm (Algorithm~\ref{algo:sensing-altmin}) when applied to the RIP-based matrix sensing:
%\begin{theorem}\label{thm:sensing_err}
% Let $M=U^*\Sigma^*V^{*^\dag }$ be a rank-$k$ matrix with non zero singular values $\sigma_1^*\geq \sigma_2^*\geq \cdots\geq \sigma_k^*$. Also, let $\calA(\cdot):\rmn \rightarrow \mathbb{R}^d$ be a linear measurement operator that satisfies $2k$-RIP with RIP constant $\delta_{2k}<\frac{(\so_k)^2}{(\so_1)^2}\cdot \frac{1}{100k}$.  Suppose, AltMinSense (Algorithm \ref{algo:sensing-altmin}) is supplied inputs $\aff$, $b=\aff(M)$. Then, $\forall T> 2\log (\|M\|_F/\epsilon)$,  iterates $\widehat{U}^{T}$ and $V^T$ satisfy: 
%$$\|M-\widehat{U}^T(V^T)^\dag\|_F\leq \epsilon.$$
%where $\epsilon>0$ is any constant. 
%Then, the $(t+1)$-th iterates $U^{t+1}$, $V^{t+1}$ of the AltMin algorithm  satisfy: 
%with $\dist\left(U^{(0)},U^*\right) \leq \frac{1}{2}$. Then $U^{t}$ and $V^{t}$ satisfy:
%\begin{align*}
%  \dist\left(V^{t+1},V^*\right) &\leq \alpha_V\cdot\dist\left(U^{t}, U^*\right)  \mbox{ and} \\
%  \dist\left(U^{t+1},U^*\right) &\leq \alpha_U\cdot\dist\left(V^{t+1}, V^*\right)
%\end{align*}
%where $\alpha_V, \alpha_U < 1$ are some constants depending only on $\delta_{2k}$. $\dist\left(U, W\right)$ denotes the principal angle based distance between subspaces spanned by the columns of $U\in \mathbb{R}^{m\times k}$ and $W\in \mathbb{R}^{m\times k}$ (See Definition~\ref{defn:dist}). 
%\end{theorem}
At a high level, we prove Theorem \ref{thm:sensing_err} by showing that  the ``distance'' between subspaces spanned by $\widehat{V}^t$ (iterate at time $t$) and $V^*$ decreases exponentially with $t$. This done based on the observation that once the (standard) matrix RIP condition (Definition \ref{defn:rip}) holds, alternating minimization can be viewed, and analyzed, as a {\bf perturbed version of the power method.} This is easiest to see for the rank-1 case below; we detail this proof, and then the more general rank-$k$ case.

In this paper, we use the following definition of distance between subspaces:
\begin{definition}\cite{GolVan96}\label{defn:dist}
  Given two matrices $\widehat{U},\widehat{W}\in \mathbb{R}^{m\times k}$, the (principal angle) distance between the subspaces spanned by the columns of $\widehat{U}$ and $\widehat{W}$ is given by:
\begin{align*}
  \dist\left(\widehat{U},\widehat{W}\right) \eqdef \twonorm{U_{\perp}^\dag W} = \twonorm{W_{\perp}^\dag U}
\end{align*}
where $U$ and $W$ are orthonormal bases of the spaces $\textrm{Span}\left(\widehat{U}\right)$ and $\textrm{Span}\left(\widehat{W}\right)$, respectively. Similarly, $U_{\perp}$ and $W_{\perp}$ are any orthonormal bases of the perpendicular spaces $\textrm{Span}\left(U\right)^{\perp}$ and $\textrm{Span}\left(W\right)^{\perp}$, respectively.
\end{definition}
{\bf Note:} {\em (a)} The distance depends only on the spaces spanned by the columns of $\widehat{U},\widehat{W}$, {\em (b)}
if the ranks of $\widehat{U}$ and $\widehat{W}$ (i.e. the dimensions of their spans) are not equal, then $\dist\left(\widehat{U},\widehat{W}\right) = 1$, and {\em (c)} $\dist\left(\widehat{U},\widehat{W}\right) = 0$ if and only if they span the same subspace of $\mathbb{R}^m$.

We now present a theorem that bounds the distance between the subspaces spanned by $\widehat{V}^t$ and $V^*$ and show that it decreases exponentially with $t$.
%The following theorem shows that the distance between the subspaces spanned by $V^t$ and $V^*$ 
\begin{theorem}\label{thm:sensing}
Let $b=\aff(M)$ where  $M$ and $\aff$ satisfy assumptions given in Theorem~\ref{thm:sensing_err}. Then, the $(t+1)$-th iterates $\widehat{U}^{t+1}$, $\widehat{V}^{t+1}$ of \as satisfy:
%with $\dist\left(U^{(0)},U^*\right) \leq \frac{1}{2}$. Then $U^{t}$ and $V^{t}$ satisfy:
\begin{align*}
  \dist\left(\widehat{V}^{t+1},V^*\right) &\leq \frac{1}{4}\cdot\dist\left(\widehat{U}^{t}, U^*\right)  \mbox{ , } \\
  \dist\left(\widehat{U}^{t+1},U^*\right) &\leq \frac{1}{4}\cdot\dist\left(\widehat{V}^{t+1}, V^*\right)
\end{align*}
where %$\alpha_V, \alpha_U < 1$ are some constants depending only on $\delta_{2k}$. 
$\dist\left(U, W\right)$ denotes the principal angle based distance (see Definition~\ref{defn:dist}).
\end{theorem}
Using Theorem \ref{thm:sensing}, we are now ready to prove Theorem \ref{thm:sensing_err}.
\begin{proof}[Proof Of Theorem~\ref{thm:sensing_err}]
Assuming correctness of Theorem~\ref{thm:sensing}, Theorem~\ref{thm:sensing_err} follows by using the following set of inequalities:
\begin{align*}
  \|M-\widehat{U}^T(\widehat{V}^T)^\dag\|_F^2&\stackrel{\zeta_1}{\leq} \frac{1}{1-\delta_{2k}}\|\aff(M-\widehat{U}^T(\widehat{V}^T)^\dag)\|_2^2,\\
&\stackrel{\zeta_2}{\leq} \frac{1}{1-\delta_{2k}}\|\aff(M(I-V^T(V^T)^\dag))\|_2^2,\\
&\stackrel{\zeta_3}{\leq} \frac{1+\delta_{2k}}{1-\delta_{2k}}\|U^*\Sigma^*(V^*)^\dag(I-V^T(V^T)^\dag))\|_F^2,\\
&\stackrel{\zeta_4}{\leq} \frac{1+\delta_{2k}}{1-\delta_{2k}} \|M\|_F^2\dist^2\left(V^T, V^*\right)\stackrel{\zeta_5}{\leq} \epsilon,
\end{align*}
where $V^T$ is an orthonormal basis of $\widehat{V}^T$, $\zeta_1$ and $\zeta_3$ follow by RIP, $\zeta_2$ holds as $\widehat{U}^T$ is the least squares solution,
$\zeta_4$ follows from the definition of $\dist(\cdot,\cdot)$ and finally $\zeta_5$ follows from Theorem~\ref{thm:sensing} and by setting $T$ appropriately.
\end{proof}
%$$\leq \frac{1}{1-\delta_{2k}}\|\aff(M(I-V^T(V^T)^\dag))\|_2^2\leq \frac{1+\delta_{2k}}{1-\delta_{2k}}\|\Sigma^*(V^*)^\dag(I-V^T(V^T)^\dag))\|_F^2\leq \frac{1+\delta_{2k}}{1-\delta_{2k}} \|M\|_F^2\dist(V^T, V^*).$$
To complete the proof of Theorem~\ref{thm:sensing_err}, we now need to prove Theorem~\ref{thm:sensing}. In the next section, we illustrate the main ideas of the proof of Theorem~\ref{thm:sensing} by applying it to a rank-$1$ matrix i.e., when $k = 1$. We then provide a proof of Theorem \ref{thm:sensing} for arbitrary $k$ in Section~\ref{sec:sensing_k}.
\subsection{Rank-$1$ Case}\label{sec:sensing_1}
In this section, we provide a proof of Theorem~\ref{thm:sensing} for the special case of $k=1$. That is, let $M=u^*\sigma^*(v^*)^\dag $ s.t. $u^*\in\mathbb{R}^{m},\ \|u^*\|_2=1$ and $v^*\in\mathbb{R}^{n}, \|v^*\|_2=1$. Also note that when $\widehat{u}$ and $\widehat{w}$ are vectors, $\dist(\widehat{u}, \widehat{w})=1-(u^\dag w)^2$, where $u=\widehat{u}/\|\widehat{u}\|_2$ and $w=\widehat{w}/\|\widehat{w}\|_2$. 

Consider the $t$-th update step in the \as procedure. As $\widehat{v}^{t+1}=$\\
$\argmin_{\widehat{v}}\sum_{i=1}^d\left(\widehat{u}^{t\dag} A_i^\dag \widehat{v}-\so u^{*^\dag} A_i^\dag v^*\right)^2$,
setting the gradient of the above objective function to $0$, we obtain:
\begin{align*}
  \left(\sum_{i=1}^d A_iu^t(u^t)^\dag A_i^\dag \right)\|\widehat{u}^t\|_2\wvto&=\sigma^*\left(\sum_{i=1}^d A_iu^tu^{*^\dag} A_i^\dag \right)\vo,%\vto&=\frac{\wvto}{\|\wvto\|} %\label{eq:r1update}
\end{align*}
where $u^t=\widehat{u}^t/\|\widehat{u}^t\|_2$. 
Now, let $B=\sum_{i=1}^d A_iu^t(u^t)^\dag A_i^\dag$ and $C=\sum_{i=1}^d A_iu^t(u^*)^\dag A_i^\dag$. Then, 
\begin{align}
  \|\widehat{u}^t\|_2\wvto&=\sigma^*B^{-1}C\vo,\nonumber\\
&=\underbrace{\ip{\uo}{\ut}\sigma^*\vo}_{\text{Power Method}}-\underbrace{B^{-1}\left(\ip{\uo}{\ut}B-C\right)\sigma^*\vo}_{\text{Error Term}}. \label{eq:error_1}
\end{align}
Note that the first term in the above expression is the power method iterate (i.e., $M^{\dag} u^t$). The second term is an error term and the goal is to show that it becomes smaller as $\ut$ gets closer to $\uo$. Note that when $\ut=\uo$, the error term is $0$ {\em irrespective} of the measurement operator $\aff$.

Below, we provide a precise bound on the error term:
\begin{lemma}
\label{lem:error_1}
Consider the error term defined in \eqref{eq:error_1} and let $\aff$ satisfy $2$-RIP with constant $\delta_{2}$. Then, 
$$\|B^{-1}\left(\ip{\uo}{\ut}B-C\right)\vo\|\leq \frac{3\delta_{2}}{1-3\delta_{2}}\sqrt{1-\ip{\ut}{\uo}^2}$$
\end{lemma}
See Appendix~\ref{app:sensing1} for a detailed proof of the above lemma. 

Using the above lemma, we now finish the proof of Theorem~\ref{thm:sensing}: 
\begin{proof}[Proof of Rank-$1$ case of Theorem~\ref{thm:sensing}]
%Here, we show that {\em assuming} $\ip{\ut}{\uo}$ to be large enough, the rank-$1$ case update \eqref{eq:r1update} leads to decrease in distance between current $v$ and $\vo$. 
%Similar, arguments can be made to show decrease in $\dist(u, \uo)$. 
Let $\vt=\wvto/\|\wvto\|_2$. Now, using \eqref{eq:error_1} and Lemma~\ref{lem:error_1}:, 
\begin{align*}
  &\ip{\vto}{\vo}=\frac{\ip{\wvto}{\vo}}{\|\wvto\|}=\frac{\ip{\wvto/\sigma^*}{\vo}}{\|\wvto/\sigma^*\|} \nonumber\\
%&=\frac{\ip{\uo}{\ut}- (\vo)^\dag B^{-1}\left(\ip{\uo}{\ut}B-C\right)\vo}{ \sqrt{ \left(\ip{\uo}{\ut}-(\vo)^\dag B^{-1}\left(\ip{\uo}{\ut}B-C\right)\vo\right)^2+\left(v_\perp^\dag B^{-1}\left(\ip{\uo}{\ut}B-C\right)\vo\right)^2 } }, \nonumber \\
&\leq \frac{\ip{\uo}{\ut}-\widehat{\delta}_2\sqrt{1-\ip{\uo}{\ut}^2}}{\sqrt{\left(\ip{\uo}{\ut}-\widehat{\delta}_2\sqrt{1-\ip{\uo}{\ut}^2}\right)^2+\widehat{\delta}_2^2\left(1-\ip{\uo}{\ut}^2\right)}},
\end{align*}
where $\widehat{\delta}_2= \frac{3\delta_2}{1-3\delta_2}$. That is, %, $v_\perp$ is the component of $\wvto$ that is perpendicular to $\vo$. The last inequality follows using Lemma~\ref{lem:error_1}. 
%Theorem~\ref{thm:sensing} follows directly using the above expression. 
\begin{align*}&\dist^2(\vto, \vo) \leq \frac{\widehat{\delta}_2^2(1-\ip{\uo}{\ut}^2)}{(\ip{\uo}{\ut}-\widehat{\delta}_2\sqrt{1-\ip{\uo}{\ut}^2})^2+\widehat{\delta}_2^2(1-\ip{\uo}{\ut}^2)},\end{align*}
Hence, {\em assuming} $\ip{\uo}{\ut}\geq 5\widehat{\delta}_2$, $\dist(\vto, \vo)\leq \frac{1}{4}\dist(\ut, \uo).$
As $\dist(u^{t+1}, \uo)$ and $\dist(v^{t+1}, \vo)$ are decreasing with $t$ (from the above bound), we only need to show that $\ip{u^0}{\ut}\geq 5\widehat{\delta}_2$. 
Recall that $\widehat{u}^0$ is obtained by using one step of SVP algorithm \cite{JainMD10}. Hence, using Lemma 2.1 of \cite{JainMD10} (see Lemma~\ref{lem:jmd}): 
$$\|\sigma_1^*(I-u^0(u^0)^\dag)\uo)\|_2^2\leq \|M-\widehat{u}^0(\widehat{v}^0)^\dag\|_F^2\leq 2\delta_{2}\|M\|_F^2.$$
Therefore, $\ip{u^0}{\uo}\geq \sqrt{1-2\delta_2}\geq 5\widehat{\delta}_2$ assuming $\delta_{2}\leq \frac{1}{100}$. 
\end{proof}
% We have the following lemma regarding linear operators satisfying $\delta_{2k}$ RIP.
% Given the observations $\calA(X)$ corresponding to a rank-$k$ matrix $X \eqdef U^*V^{*^\dag }$, we wish to reconstruct $X$.
% For this purpose, we will use the alternating minimization algorithm\cite{}. We start with an initial estimate $U^{(0)}$.
% We obtain $V^{(0)}$ by minimizing the error $\sum_i \left(\Tr\left(A_i X\right) - \Tr\left(A_i U^{(0)}V^\dag \right)\right)^2$
% over all $V\in \R^{n\times k}$. We then fix $V^{(0)}$ and minimize $\sum_i \left(\Tr\left(A_i X\right) -
% \Tr\left(A_i UV^{(0)^\dag }\right)\right)^2$ to obtain $U^{(1)}$. We then continue this procedure for some number of steps.
% The algorithm is formally given in Algorithm \ref{algo:sensing-altmin}. The main result of this section is the following
% theorem which says that if $\calA(\cdot)$ satisfies $2k$-RIP for some small enough constant $\delta$ then the error between
% $\left(U^{t},V^{t}\right)$ obtained by the algorithm at the end of step $t$ and $\left(U^*,V^*\right)$, the truth goes
% to $0$ geometrically in $t$. To formalize the notion of distance between subspaces, we have the following definition from
% \cite{GolVan96}:
% \begin{theorem}\label{thm:sensing-altmin-geomconv}
%   Suppose $\calA(\cdot)$ satisfies $2k$-RIP with $\delta$.
% Let $X=U^*V^{*^\dag }$ be a rank-$k$ matrix with non zero singular values between $\gamma$ and $\Gamma$.
% Assume further that $\delta < \frac{\gamma}{10k\Gamma}$.
% Suppose Algorithm \ref{algo:sensing-altmin} is run with input $\calA\left(\right)$ and $U^{(0)}$
% with $\dist\left(U^{(0)},U^*\right) \leq \frac{1}{2}$. Then $U^{t}$ and $V^{t}$ satisfy:
% \begin{align*}
%   \dist\left(U^{t},U^*\right) &\leq c_1 c_2^t \mbox{ and} \\
%   \dist\left(V^{t},V^*\right) &\leq c_1 c_2^t
% \end{align*}
% where $c_1$ and $c_2 < 1$ are some constants depending only on $\delta$.
% \end{theorem}


% If $\calA$ satisfies $2k$-RIP with $\delta=0$ then alternating minimization is the same as orthogonal iteration.
% So, if $\dist\left(U^{(0)},U^*\right) < 1$ then we have $\dist\left(V^{(1)},V^*\right)=0$ i.e., Algorithm
% \ref{algo:sensing-altmin} converges after just one step to the truth. However, since we have $\delta > 0$,
% Algorithm \ref{algo:sensing-altmin} has an error term compared to the orthogonal iteration. We will prove
% Theorem \ref{thm:sensing-altmin-geomconv} by analyzing the error term and showing that it decreases with
% every iteration under the theorem's assumptions. As a first step in this process we will now characterize
% $V^{t}$ as the solution of a set of linear equations.
\subsection{Rank-$k$ Case}
\label{sec:sensing_k}
In this section, we present the proof of Theorem \ref{thm:sensing} for arbitrary $k$, i.e., when $M$ is a rank-$k$ matrix (with SVD $U^*\Sigma^* \left(V^*\right)^{\dag})$.

Similar to the analysis for the rank-$1$ case (Section \ref{sec:sensing_1}), we show that even for arbitrary $k$, the updates of AltMinSense are essentially power-method type updates but with a bounded error term whose magnitude decreases with each iteration.

However, directly analyzing iterates of \as is a bit tedious due to non-orthonormality of intermediate iterates $\widehat{U}$. Instead, {\bf for analysis only} we consider the iterates of a modified version of \ase, where we explicitly orthonormalize each iterate using the QR-decomposition\footnote{The QR decomposition factorizes a matrix into an orthonormal matrix (a basis of its column space) and an upper triangular matrix; that is given $\widehat{S}$ it computes $\widehat{S} = SR$ where $S$ has orthonormal columns and $R$ is upper triangular. If $\widehat{S}$ is full-rank, so are $S$ and $R$.}. In particular, suppose we replace steps 4 and 5 of \ase with the following 
\begin{align}
\widehat{U}^{t}&= \, U^{t}R_U^{t}\ \ \text{ (QR decomposition), }\nonumber\\
\widehat{V}^{t+1}&\leftarrow \, \argmin_{V} \  \|\aff(U^t V^\dag )-b\|_2^2,\nonumber\\ 
\widehat{V}^{t+1}&= \, V^{t+1}R_V^{t+1}\ \ \text{ (QR decomposition) } \nonumber  \\
\widehat{U}^{t+1}&\leftarrow \, \argmin_{U} \  \|\aff(U (V^{t+1})^\dag )-b\|_2^2 \label{eq:update_k1}
\end{align}
In our algorithm, in each iterate both $\widehat{U}^{t},\widehat{V}^{t}$ remain full-rank because $\dist\left(U^t,U^*\right)<1$; with this, the following lemma implies that the spaces spanned by the iterates in our \ase ~ algorithm are {\em exactly the same} as the respective ones by the iterates of the above modified version (and hence the distances $\dist(\widehat{U}^{t},U^*)$ and $\dist(\widehat{V}^{t},V^*)$ are also the same for the two algorithms).

\begin{lemma}\label{lemma:QR-irrelevance}
Let $\widehat{U}^t$ be the $t^{th}$ iterate of our AltMinSense algorithm, and $\widetilde{U}^t$ of the modified version stated above. Suppose also that both $\widehat{U}^t,\widetilde{U}^t$ are full-rank, and span the same subspace. Then the same will be true for the subsequent iterates for the two algorithms, i.e. $Span(\widehat{V}^{t+1}) = Span(\widetilde{V}^{t+1})$, $Span(\widehat{U}^{t+1}) = Span(\widetilde{U}^{t+1})$, and all matrices at iterate $t+1$ will be full-rank.
\end{lemma}

The proof of the above lemma can be found in Appendix \ref{app:sensingk}. In light of this, we will now prove Theorem~\ref{thm:sensing} {\bf with} the new QR-based iterates \eqref{eq:update_k1}.

\begin{lemma}\label{lemma:sensing-updateeqn1}
  Let $\Uht$ be the $t$-th step iterate of \as and let $\Ut, \widehat{V}^{t+1}$ and $\Vt$ be obtained by Update \eqref{eq:update_k1}. Then, 
\begin{align}
  \widehat{V}^{t+1} = \underbrace{V^*\Sigma^*{U^*}^\dag  U^{t}}_{\substack{\text{Power-method}\\ \text{Update}}} - \underbrace{F}_{\substack{\text{Error}\\ \text{Term}}},\ \ 
  V^{t+1}=\widehat{V}^{t+1}(R^{(t+1)})^{-1},\label{eq:update_k}
\end{align}
where $F$ is an error matrix defined in \eqref{eq:F_k} and $R^{(t+1)}$ is a triangular matrix obtained using QR-decomposition of $\widehat{V}^{t+1}$. 
%  and $\widehat{V}^{t}$ in time step $t$ of Algorithm \ref{algo:sensing-altmin}.
%Let  and $\widehat{V}^{t}=\left[\widehat{v}_1^{t},\cdots,
%\widehat{v}_k^{t}\right]$ where each $u_p^{t}$ and $\widehat{v}_p^{t} \in \R^n$. Then $\widehat{v}_q^{t}$
%satisfy the following set of equations:
%\begin{align}\label{eqn:sensing-updateeqn}
%  \sum_{q=1}^k \left(\sum_{i=1}^d A_i u_p^{t} u_q^{(t)^\dag } A_i^\dag \right)\widehat{v}_q^{t}
%	= \sum_{q=1}^k \left(\sum_{i=1}^d A_i u_p^{t} u_q^{*^\dag } A_i^\dag \right) v_q^* \;\;\;\;\; \forall \; p\in[k].
%\end{align}
\end{lemma}
See Appendix~\ref{app:sensing} for a detailed proof of the above lemma. 

%The above lemma shows that alternating minimization update at each step essentially performs a QR-iteration type update (used for SVD computation) but we also have an extra error term given by $F$.
Before we give an expression for the error matrix $F$, we define the following notation.
Let $v^* \in \R^{nk}$ be given by: $v^*=vec(V^*)$, i.e., $v^*=\left[v_1^{*\dag} v_2^{*\dag}\dots v_k^{*\dag}\right]^\dag$. 
%and $S\in \R^{nk\times nk}$ be given by
%$$v^*=\left[\begin{matrix}v_1^*\\\vdots\\v_k^*\end{matrix}\right], S=\left[\begin{matrix}\sigma_1^*I_{n}&\dots&0_{n}\\\vdots&\ddots&\vdots\\0_n&\dots&\sigma_k^*I_n\end{matrix}\right],$$
%i.e., $v^*$ is a vectorized version of the matrix $V^*$ by stacking all its columns in a single column and $S$ is a diagonal matrix. We denote the above vectorizing operation by $v^* = vec\left(V^*\right)$.
Define $B$, $C$, $D$, $S$ as follows: {\small 
\begin{align}\label{eq:BCD_k}
& B \eqdef \left[\begin{array}{ccc}
                 B_{11} & \cdots & B_{1k} \\
		 \vdots  	& \ddots &\vdots \\
                 B_{k1} & \cdots & B_{kk} \\
                \end{array} \right]
\mbox{ , }
 C \eqdef \left[\begin{array}{ccc}
                 C_{11} & \cdots & C_{1k} \\
		 \vdots  	& \ddots &\vdots \\
                 C_{k1} & \cdots & C_{kk} \\
                \end{array} \right]\mbox{ , }\nonumber\\
&D\eqdef \left[\begin{array}{ccc}
                 D_{11}  & \cdots & D_{1k} \\
		 \vdots 	& \ddots &\vdots \\
                 D_{k1} & \cdots & D_{kk} \\
                \end{array} \right]\mbox{ , }
S\eqdef \left[\begin{matrix}\sigma_1^*I_{n}&\dots&0_{n}\\\vdots&\ddots&\vdots\\0_n&\dots&\sigma_k^*I_n\end{matrix}\right]. 
\end{align}}
where , for $1\leq p,q\leq k$: $B_{pq} \eqdef \sum_{i=1}^d A_i u_p^{t}{u_q^{t}}^\dag  A_i^\dag $,\\ $C_{pq} \eqdef \sum_{i=1}^d A_i u_p^{t}{u_q^*}^\dag  A_i^\dag, $ and, $D_{pq} \eqdef \iprod{u_p^{t}}{u_q^*}\mathbb{I}_{n\times n}.$ Recall that, $ u_p^{t}$ is the $p$-th column of $U^{t}$ and $u_q^*$ is the $q$-th left singular vector of the underlying matrix $M=U^*\Sigma^* (V^*)^\dag$. 
Finally $F$ is obtained by ``de-stacking'' the vector\\ $B^{-1}\left(BD-C\right)Sv^*$ i.e., the $i^{\textrm{th}}$ column of $F$ is given by:{\small
\begin{align}
  F_i&\eqdef\left[\begin{matrix} \left(B^{-1}\left(BD-C\right)Sv^*\right)_{ni+1} \\ \left(B^{-1}\left(BD-C\right)Sv^*\right)_{ni+2} \\ \vdots \\ \left(B^{-1}\left(BD-C\right)Sv^*\right)_{ni+n}
            \end{matrix}\right], F \eqdef \left[F_1 \; F_2 \;\cdots \;F_k \right]. \label{eq:F_k}
\end{align}}
Note that the notation above should have been $B^t,C^t$ and so on. We suppress the dependence on $t$ for notational simplicity. 
% To make the proof more manageable, we introduce the following notation. Let $p,q\in[k]$. Define the
% $n\times n$ matrices $B_{pq} \eqdef \sum_{i=1}^d A_i u_p^{t}{u_q^{t}}^\dag  A_i^\dag $ and $C_{pq} \eqdef \sum_{i=1}^d A_i u_p^{t}{u_q^*}^\dag  A_i^\dag $.
% Using this notation, \eqref{eqn:sensing-updateeqn} can be written as follows:
% \begin{align*}
%   \sum_{q=1}^k B_{pq} \widehat{v}_q^{t} = \sum_{q=1}^k C_{pq} v_q^* \;\;\;\; \forall p \in [k]
% \end{align*}
% Now define two $nk\times nk$ matrices $B$ and $C$ as follows :
% \begin{align*}
%  B \eqdef \left[\begin{array}{cccc}
%                  B_{11} & B_{12} & \cdots & B_{1k} \\
%                  B_{21} & B_{22} & \cdots & B_{2k} \\
% 		 \vdots & 	& \ddots &\vdots \\
%                  B_{k1} & B_{k2} & \cdots & B_{kk} \\
%                 \end{array} \right]
% \mbox{ and }
%  C \eqdef \left[\begin{array}{cccc}
%                  C_{11} & C_{12} & \cdots & C_{1k} \\
%                  C_{21} & C_{22} & \cdots & C_{2k} \\
% 		 \vdots & 	& \ddots &\vdots \\
%                  C_{k1} & C_{k2} & \cdots & C_{kk} \\
%                 \end{array} \right]
% \end{align*}
% where the blocks $B_{pq}$ and $C_{pq}$ are the $n\times n$ matrices defined above. Also define the vectors $\widehat{v}^{t}$
% and $v^* \in \R^{nk}$ as the vectors obtained by stacking the vectors $\widehat{v}_q^{t}$ and $v_q^*$ respectively i.e.,
% \begin{align*}
%   \widehat{v}^{t} \eqdef \left[\begin{array}{c}
% 				    \widehat{v}_1^{t} \\
% 				    \widehat{v}_2^{t} \\
% 				    \vdots \\
% 				    \widehat{v}_k^{t} \\
%                                  \end{array}\right]
% \mbox{ and }
%   v^* \eqdef \left[\begin{array}{c}
% 		      v_1^* \\
% 		      v_2^* \\
% 		      \vdots \\
% 		      v_k^* \\
% 		    \end{array}\right].
% \end{align*}
% Using all the above notation, \eqref{eqn:sensing-updateeqn} can be concisely written as follows:
% \begin{align}\label{eqn:sensing-updateeqn-concise}
%   B\widehat{v}^{t} = Cv^*.
% \end{align}
% As we already mentioned, \eqref{eqn:sensing-updateeqn-concise} is very similar to the first step in the orthogonal
% iteration (please refer \cite{GolVan96} for details). To capture this connection, we introduce a block diagonal $nk\times nk$
% matrix $D$:
% \begin{align*}
%  D \eqdef \left[\begin{array}{cccc}
%                  D_{11} & D_{12} & \cdots & D_{1k} \\
%                  D_{21} & D_{22} & \cdots & D_{2k} \\
% 		 \vdots & 	& \ddots &\vdots \\
%                  D_{k1} & D_{k2} & \cdots & D_{kk} \\
%                 \end{array} \right]
% \end{align*}
% where each block $D_{pq} \eqdef \iprod{u_p^{t}}{u_q^*}\mathbb{I}_{n\times n}$. \eqref{eqn:sensing-updateeqn-concise}
% can now be modified as follows:
Now, from Update \eqref{eq:update_k}, we have
\begin{align}
  &V^{t+1} = \widehat{V}^{t+1}{R^{(t+1)}}^{-1} = \left(V^*\Sigma^*{U^*}^\dag  U^{t} - F\right){R^{(t+1)}}^{-1} \nonumber\\
  \Rightarrow & {V_{\perp}^*}^\dag V^{t+1} = -{V_{\perp}^*}^\dag  F {R^{(t+1)}}^{-1}.\label{eqn:sensing-orthogonaliter-error}
\end{align}
where $V_{\perp}^*$ is an orthonormal basis of  $\textrm{Span}\left(v_1^*,v_2^*,\cdots,v_k^*\right)^{\perp}$. Therefore, 
\begin{multline*}
dist(V^*, V^{t+1})=\|{V_{\perp}^*}^\dag V^{t+1}\|_2=\|{V_{\perp}^*}^\dag  F {R^{(t+1)}}^{-1}\|_2\leq \|F(\Sigma^*)^{-1}\|_2\|\Sigma^*{R^{(t+1)}}^{-1}\|_2.\end{multline*}
Now, we break down the proof of Theorem \ref{thm:sensing} into the following two steps: 
\begin{itemize}\setlength{\topsep}{0pt}\setlength{\itemsep}{0pt}
%  \item the minimum singular value of $B$ is not very small (Lemma \ref{lemma:sensing-minsingval-B}),
  \item show that $\twonorm{F(\Sigma^*)^{-1}}$ is  small (Lemma \ref{lemma:sensing-twonorm-F}) and
  \item show that $\|\Sigma^*{R^{(t+1)}}^{-1}\|_2$ is small(Lemma \ref{lemma:sensing-minsingval-Rt}).
\end{itemize}

%Note that if RIP constant $\delta_{2k}=0$ then $F=0$ and Update \eqref{eq:update_k} is just the standard orthogonal iteration update and hence $V^{(t+1)}=V^*$. 
%We would like to stress again that if $\calA(\cdot)$ satisfied $2k$-RIP with $\delta=0$, then
%we would have $F=0$ and 
%\eqref{eqn:sensing-orthogonaliter-error} would just be the orthogonal iteration.
% For $\delta$ small enough, we will show that $F$ is small (in an appropriate norm) and hence the procedure is very
%close to the orthogonal iteration giving us the required error bound. The proof of Theorem \ref{thm:sensing-altmin-geomconv}
%can now be broken down into three important steps: show that
%\begin{itemize}
%  \item the minimum singular value of $B$ is not very small (Lemma \ref{lemma:sensing-minsingval-B}),
%  \item $\twonorm{F}$ is very small (Lemma \ref{lemma:sensing-twonorm-F}) and
%  \item the minimum singular value of $R^{t}$ is not very small (Lemma \ref{lemma:sensing-minsingval-Rt}).
%\end{itemize}
We will now state the two corresponding lemmas.
Complete proofs can be found in Appendix \ref{app:sensingk}
% We first show that the minimum eigenvalue of $B$ is lower-bounded by $1-\delta_{2k}$ where $\delta_{2k}$ is the RIP constant of $\aff$ for rank-$2k$ matrices. This also justifies inversion of $B$ in \eqref{eq:F_k}. 
%To simplify the notation in the following, we denote $u^{t}, u_p^{t}, U^{t}$ etc. by $u, u_p, U$ etc. respectively and
%$v^{t}, v_p^{t}, V^{t}$ etc. by $v, v_p, V$ etc. respectively.
%Recall that the non zero singular values of $X$ lie between $\gamma$ and $\Gamma$. So we can assume without loss of generality that
%\begin{align}
%  \sqrt{\gamma} \leq \twonorm{u_p^*} \leq \sqrt{\Gamma} \label{eqn:sensing-up*-norm} \\
%  \sqrt{\gamma} \leq \twonorm{v_p^*} \leq \sqrt{\Gamma} \label{eqn:sensing-vp*-norm}
%\end{align}
The first lemma bounds the spectral norm of $F(\Sigma^*)^{-1}$.
%Let $\lambda \eqdef \sqrt{1-\dist\left(U,U^*\right)^2}$. $\lambda$ quantifies how aligned the subspaces $U$ and $U^*$ are -the higher the $\lambda$ the more aligned $U$ and $U^*$ are.
\begin{lemma}\label{lemma:sensing-twonorm-F}
Let linear measurement $\aff$ satisfy RIP for all $2k$-rank matrices and let $b=\aff(M)$ with $M\in \rmn$ being a rank-$k$ matrix. Then, spectral norm of error matrix $F(\Sigma^*)^{-1}$ (see Equation~\ref{eq:update_k}) after $t$-th iteration update satisfy:
%We have the following bound on $\twonorm{F}$:
\begin{align}\label{eqn:sensing-twonorm-F}
  \twonorm{F(\Sigma^*)^{-1}} \leq \frac{\delta_{2k} k }{1-\delta_{2k}}dist(U^t, U^*).
\end{align}
\end{lemma}
The following lemma bounds the spectral norm of $\Sigma^*R^{(t+1)^{-1}}$.
\begin{lemma}\label{lemma:sensing-minsingval-Rt}
Let linear measurement $\aff$ satisfy RIP for all $2k$-rank matrices and let $b=\aff(M)$ with $M\in \rmn$ being a rank-$k$ matrix. Then, 
\begin{equation}\label{eqn:sensing-minsingval-Rt}
  \|\Sigma^*(R^{(t+1)})^{-1}\|_2\leq \frac{\sigma^*_1/\sigma^*_k}{\sqrt{1-\dist^2\left(U^{t},U^*\right)}-\frac{(\sigma^*_1/\sigma^*_k)\delta_{2k}k\dist(U^{t},U^*)}{1-\delta_{2k}}}.%\sigma_{\textrm{min}}\left(R^{t}\right) \geq \gamma \lambda - \twonorm{F}.
\end{equation}
\end{lemma}
With the above two lemmas, we now prove Theorem \ref{thm:sensing}.
\begin{proof}[Proof Of Theorem \ref{thm:sensing}]
  Using \eqref{eqn:sensing-orthogonaliter-error}, \eqref{eqn:sensing-twonorm-F} and \eqref{eqn:sensing-minsingval-Rt}, we
obtain the following:
\begin{align}
  \dist\left(V^{t+1},V^*\right) &= \twonorm{{V_{\perp}^*}^\dag V^{t+1}},\nonumber\\
	&\leq \twonorm{{V_{\perp}^*}^\dag F(\Sigma^*)^{-1}\Sigma^*{R^{(t+1)}}^{-1}}, \nonumber\\
	&\leq \twonorm{V_{\perp}^*}\twonorm{F(\Sigma^*)^{-1}}\twonorm{\Sigma^*{R^{(t+1)}}^{-1}} \nonumber\\
%	&\leq 1 \cdot \frac{\delta k  \dist\left(U^{t},U^*\right)}{1-\delta} \cdot \frac{1}{\sigma_{\textrm{min}}\left(R^{t}\right)} \\
	&\leq \frac{(\sigma^*_1/\sigma^*_k)\delta_{2k} k \cdot \dist\left(U^{t},U^*\right)}{(1-\delta_{2k})L},\label{eq:distb_sensing}		
\end{align}
where $L= \sqrt{1-\dist\left(U^{t},U^*\right)^2} - 
			\frac{(\sigma^*_1/\sigma^*_k)\delta_{2k} k \dist\left(U^{t},U^*\right)}{1-\delta_{2k}}$. 
Also, note that $U^{0}$ is obtained using SVD of $\sum_i A_i b_i$. Hence, using Lemma~\ref{lem:jmd}, we have:
\begin{align}
& \|\aff(U^{0}\Sigma^{0}V^{0}-U^*\Sigma^*(V^*)^\dag\|_2^2\leq 4\delta_{2k}\|\aff(U^*\Sigma^*(V^*)^\dag)\|_2^2,\nonumber\\
\Rightarrow& \|U^{0}\Sigma^{0}V^{0}-U^*\Sigma^*(V^*)^\dag\|_F^2\leq 4\delta_{2k}(1+3\delta_{2k})\|\Sigma^*\|_F^2,\nonumber\\
\Rightarrow&\|U^{0}(U^{0})^\dag U^*\Sigma^*(V^*)^\dag-U^*\Sigma^*(V^*)^\dag\|_F^2\leq 6\delta_{2k}\|\Sigma^*\|_F^2,\nonumber\\
\Rightarrow&(\sigma_k^*)^2\|(U^{0}(U^{0})^\dag-I)U^*\|_F^2\leq 6\delta_{2k}k(\sigma_1^*)^2,\nonumber\\
\Rightarrow&\dist(U^{0}, U^*)\leq \sqrt{6\delta_{2k}k}\left(\frac{\sigma_1^*}{\sigma_k^*}\right)<\frac{1}{2}.
\end{align}
Using \eqref{eq:distb_sensing} with  $\dist\left(U^{0},U^*\right) < \frac{1}{2}$ and $\delta_{2k} < \frac{1}
{24(\sigma^*_1/\sigma^*_k)^2k}$, we obtain: $\dist\left(V^{t},V^*\right) < \frac{1}{4} \dist\left(U^{t},U^*\right)$.  
Similarly we can show that $\dist\left(U^{t+1},U^*\right) < \frac{1}{4} \dist\left(V^{t},V^*\right)$. 
\end{proof}
% @@@@To see if this sort of statement makes sense and is not confusing.
% If yes, we can keep it else we can discard it and keep our existing algorithm.@@@@
% The above lemma also shows that $U\left(\wh{V}\right)^{\dag} = \wh{U} \left(\wt{V}\right)^{\dag}$ and that the column spaces of $\wh{V}$ and $\wt{V}$ are equal.
%\subsection{Noisy Case}
%\label{sec:sensing_noisy}
%%% Local Variables: 
%%% mode: latex
%%% TeX-master: "altmin"
%%% End: 

\newcommand{\uht}{\widehat{u}^t}
\section{Matrix Completion}
In this section, we study the Matrix Completion problem \eqref{eq:matcomp} and show that, assuming $k$ and $\frac{\so_1}{\so_k}$ are constant, \amc (Algorithm~\ref{algo:completion-altmin}) recovers the underlying matrix $M$ using only $O(n \log n)$ measurements (i.e., we prove Theorem~\ref{thm:sampling-altmin-main}).

As mentioned, while observing elements in $\Omega$ constitutes a linear map, matrix completion is different from matrix sensing because the map does not satisfy RIP. The (now standard) approach is to assume incoherence of the true matrix $M$, as done in Definition \ref{defn:sampling-incoherentmatrices}. With this, and the random sampling of $\Omega$, matrix completion exhibits similarities to matrix sensing. For our analysis, we can again use the fact that incoherence allows us to view alternating minimization as a perturbed power method, whose error we can control.

However, there are important differences between the two problems, which make the analysis of completion more complicated. Chief among them is the fact that we need to establish {\em the incoherence of each iterate}. For the first initialization $\widehat{U}^0$, this necessitates the ``clipping" procedure (described in step 4 of the algorithm). For the subsequent steps, this requires the partitioning of the observed $\Omega$ into $2T+1$ sets (as described in step 2 of the algorithm).

As in the case of matrix sensing, we prove our main result for matrix completion (Theorem~\ref{thm:sampling-altmin-main}) by first establishing a geometric decay of the distance between the subspaces spanned by $\widehat{U}^{t}, \widehat{V}^{t}$ and $U^*,V^*$ respectively.
\begin{theorem}\label{thm:sampling-altmin-geomconv}
 Under the assumptions of Theorem \ref{thm:sampling-altmin-main}, the $(t+1)^{\textrm{th}}$ iterates $\Uw^{t+1}$ and $\Vw^{t+1}$
satisfy the following property w.h.p.:
\begin{align*}
 \dist\left(\Vw^{t+1}, V^*\right) &\leq \frac{1}{4} \dist\left(\Uw^{t}, U^*\right)  \mbox{ and } \\
 \dist\left(\Uw^{t+1}, U^*\right) &\leq \frac{1}{4} \dist\left(\Vw^{t+1}, V^*\right),\;\;\;\; \forall \;\;1 \leq t \leq T.
\end{align*}
\end{theorem}
We use the above result along with incoherence of $M$ to prove Theorem~\ref{thm:sampling-altmin-main}. See Appendix~\ref{app:mc} for a detailed proof. 

Now, similar to the matrix sensing case, alternating minimization needs an initial iterate that is close enough to $U^*$ and $V^*$, from
where it will then converge. To this end, Steps $3-4$ of Algorithm \ref{algo:completion-altmin}  use SVD of $P_\Omega(M)$ followed by clipping to initialize $\widehat{U}^0$. While the SVD step guarantees that  $\widehat{U}^0$ is close enough to $U^*$, it might not remain incoherent. To maintain incoherence, we introduce an extra clipping step which guarantees incoherence of $\widehat{U}^0$ while also ensuring that $\widehat{U}^0$ is close enough to $U^*$ (see Lemma~\ref{lemma:svd-clipping-jointguarantee})
%Now, since  $M$ is incoherent, the top $k$ singular components of $\frac{1}{p}P_{\Omega}(M)$ are close to the top $k$ singular components of $M$
%(Lemma \ref{lemma:first-step-svd}). However the top $k$ singular components of $\frac{1}{p}P_{\Omega}(M)$ need not be incoherent.
%Lemma \ref{lemma:clipping-incoherence-distance} then guarantees that clipping the resulting vectors to make them incoherent
%does not increase the distance between $U^0$ and $U^*$ by much. We combine the above two results into the following lemma about  ``goodness'' of $U^0$. %Combining the above results, we have the following formal result. In what follows, we assume that $M,\Omega,p$ etc. are as in the statement of Theorem \ref{thm:sampling-altmin-main}.
\begin{lemma}\label{lemma:svd-clipping-jointguarantee}
Let $M,\Omega,p$ be as defined in Theorem~\ref{thm:sampling-altmin-main}. Also, let $U^0$ be the initial iterate obtained by  step $4$ of Algorithm \ref{algo:completion-altmin}. Then, w.h.p. we have% we have,% the following results regarding $U^{0}$:
\begin{itemize}\setlength{\topsep}{-5pt}\setlength{\itemsep}{-5pt}
  \item	$\dist\left(U^{0},U^*\right) \leq \frac{1}{2}$ and
  \item	$U^{0}$ is incoherent with parameter $4\mu \sqrt{k}$.
\end{itemize}
\end{lemma}
%We will prove the following lemma which says that if $V^*$ is incoherent and we have an orthonormal column matrix $V$ such
%that $\dist\left(V,V^*\right) \leq d$, then we can obtain a orthonormal column matrix $\widetilde{V}$ such that:
%\begin{itemize}
%  \item	$\dist\left(\widetilde{V},V^*\right) \leq d + \epsilon$ and
%  \item	$\widetilde{V}$ is incoherent.
%\end{itemize}
The above lemma guarantees a ``good'' starting point for alternating minimization. Using this, we now present a proof of Theorem~\ref{thm:sampling-altmin-geomconv}.  %we can focus our attention on the analysis of alternating minimization.
Similar to the sensing section, we first explain key ideas of our proof using rank-$1$ example. Then in Section~\ref{sec:comp_general} we extend our proof to  general rank-$k$ matrices. 
\subsection{Rank-$1$ Case}
Consider the rank-$1$ matrix completion problem where $M=\so\uo(\vo)^\dag$.
Now, the $t$-th step iterates $\vht$ of Algorithm \ref{algo:completion-altmin} are given by: 
\begin{align*}
  \vht=\argmin_{\widehat{v}} \sum_{(i,j)\in \Omega}(M_{ij}-\widehat{u}^t_i\widehat{v}_j)^2.%,\ \ \ \vt=\frac{\vht}{\|\vht\|_2}. 
\end{align*}
Let $\ut=\uht/\|\uht\|_2$. Then, $\forall j$:{\small
\begin{align}
\hspace*{-10pt}&\|\uht\|_2\sum_{i:(i,j)\in \Omega}(\ut_i)^2\vht_j =\so\sum_{i:(i,j)\in\Omega}\ut_i\uo_i\vo_j \nonumber\\
&\Rightarrow \|\uht\|_2\vht_j = \frac{\so}{\sum_{i:(i,j)\in \Omega}(\ut_i)^2}\sum_{i:(i,j)\in\Omega}\ut_i\uo_i\vo_j\nonumber\\
&=\sigma^*\ip{\ut}{\uo}\vo_j-\frac{\sigma^*(\ip{\ut}{\uo}\sum_{i:(i,j)\in\Omega}(\ut_i)^2\vo_j-\sum_{i:(i,j)\in\Omega}\ut_i\uo_i\vo_j)}{\sum_{i:(i,j)\in\Omega}(\ut_i)^2}.\label{eq:r1compupdate}
\end{align}}
Hence,
\begin{align}
\hspace*{-5pt}  \|\uht\|_2\vht=\underbrace{\ip{\uo}{\ut}\sigma^*\vo}_{\text{Power Method}}-\underbrace{\sigma^*B^{-1}\left(\ip{\ut}{\uo}B-C\right)\vo}_{\text{Error Term}},%\left(\vo(\uo)^\dag-\frac{B^{-1}}{p} P_\Omega\left(\vo(\uo)^\dag(\ut(\ut)^\dag-I)\right)\right)\ut.
\label{eq:compvtupdate}
\end{align}
where $B, C\in \mathbb{R}^{n\times n}$ are diagonal matrices, such that, 
\begin{equation}
  \label{eq:bc_comp}
  B_{jj}=\frac{\sum_{i:(i,j)\in \Omega}(\ut_i)^2}{p}, \ \ \ C_{jj}=\frac{\sum_{i:(i,j)\in\Omega}\ut_i\uo_i}{p}. 
\end{equation}
Note the similarities between the update \eqref{eq:compvtupdate} and  the rank-$1$ update \eqref{eq:error_1} for the sensing case. Here again, it is essentially a power-method update (first term) along with a bounded error term (see Lemma~\ref{lem:errb_comp1}). Using this insight, we now prove Theorem~\ref{thm:sampling-altmin-geomconv} for the special case of rank-$1$ matrices. 
%\begin{proof}[of Theorem~\ref{thm:sampling-altmin-geomconv} for rank-$1$ case]
Our proof can be divided in three major steps:
\begin{itemize}\setlength{\topsep}{-5pt}\setlength{\itemsep}{-5pt}
\item {\em Base Case}: Show that $u^0=\widehat{u}^0/\|\widehat{u}^0\|_2$ is incoherent and have small distance to $\uo$ (see Lemma~\ref{lemma:svd-clipping-jointguarantee}).  
\item {\em Induction Step (distance)}: Assuming $\ut=\widehat{u}^t/\|\uht\|_2$ to be incoherent and that $\ut$ has a small distance to $\uo$, $\vt$ decreases distances to $\vo$ by at least a constant factor. 
\item {\em Induction Step (incoherence)}: Show incoherence of $\vt$, while assuming incoherence of $\ut$ (see Lemma~\ref{lem:incoherence_vt})
\end{itemize}
%The first step of our proof follows directly from Lemma~\ref{lemma:svd-clipping-jointguarantee}. Also, the last step of our proof follows from Lemma~\ref{lem:incoherence_vt}. We stress that $\vt$ does not increase the incoherence parameter ($\mu_1$) when compared to that of $\ut$.
We first prove the second step of our proof. To this end, we provide the following lemma that bounds the error term. See Appendix~\ref{app:rank1_comp} for a proof of the below given lemma. 
\begin{lemma}\label{lem:errb_comp1}
Let $M$, $p$, $\Omega$, $\ut$ be as defined in Theorem~\ref{thm:sampling-altmin-main}. Also, let $\ut$ be a unit vector with incoherence parameter $\mu_1=\frac{6(1+\delta_2)\mu}{1-\delta_2}$.%=\sigma^*\uo\vo$ be a $\mu$-incoherent matrix and let $\Omega$ be sampled uniformly and independently of $\ut, \uo, \vo$, with $p\geq \frac{C\mu_1^2\log n}{n\delta_2^2}$ ($C>0$ is a global constant). 
%Consider the error term defined in \eqref{eq:r1compupdate} and let samples $\Omega$ and underlying matrix $M=\sigma^*\uo\vo$ satisfy assumptions mentioned in Theorem~\ref{thm:sampling-altmin-geomconv}. 
%Also, let $\ut$ be a unit $L_2$ norm vector with incoherence parameter $\mu_1$. 
Then, w.p. at least $1-\frac{1}{n^3}$: 
$$\|B^{-1}\left(\ip{\uo}{\ut}B-C\right)\vo\|_2\leq \frac{\delta_{2}}{1-\delta_{2}}\sqrt{1-\ip{\ut}{\uo}^2}.$$
\end{lemma}
%Our proof for rank-$1$ matrix completion proceeds similarly to the corresponding proof for the matrix sensing problem. That is, at each step we show that the error term is bounded and show a geometric decay in the distance to the optimal vectors $\vo, \uo$. 

%However, there is one additional challenge in the matrix completion case, which is: showing incoherence of each iterate. For this purpose, we need independence of the iterate from samples, so that we can  use standard concentration bounds. Hence, in Algorithm \ref{algo:completion-altmin}, we sample a different $\Omega$ in each step of alternating minimization, that is {\em independent} of the current iterates. As we will show later,  to obtain an error of $\epsilon$, the algorithm needs to run for $O\left(\log \frac{\|M\|_F}{\epsilon}\right)$ iterations.
%So the total number of samples needed increases by a multiplicative factor of  $O\left(\log \frac{\|M\|_F}{\epsilon}\right)$.

%We first show that assuming $t$-th step iterate $\ut$ to be incoherent and having a small distance to $\uo$, $\vt$ decreases distance to $\vo$ by at least a constant factor. Then, we show incoherence of updated iterate $\vt$ in Lemma~\ref{lem:incoherence_vt} and finally, we show that the initial vector $u^0$ satisfies the required incoherence and distance assumptions.
%Note that in Algorithm \ref{algo:completion-altmin}, we sample a different $\Omega$ in each step of alternating minimization - this is to ensure that the indices in $\Omega$ are independent of the current
%iterates which is necessary for our analysis. Since the error decays geometrically, to obtain an error of $\epsilon$, the algorithm needs to run for $O\left(\log \frac{1}{\epsilon}\right)$ iterations.
%So the total number of samples needed is equal to the number of samples needed in each iteration multiplied by $O\left(\log \frac{1}{\epsilon}\right)$.
Multiplying \eqref{eq:compvtupdate} with $\vo$ and using Lemma~\ref{lem:errb_comp1}, we get:
%Using \eqref{eq:compvtupdate} and Lemma~\ref{lem:errb_comp1}, 
\begin{equation}
  \label{eq:ipvtvo_comp}
  \|\uht\|_2\ip{\widehat{v}^{t+1}}{\vo}\geq \so\ip{\ut}{\uo}-2\so\delta_2\sqrt{1-\ip{\ut}{\uo}^2},
\end{equation}
where $\delta_2<\frac{1}{12}$ is a constant defined in the Theorem statement and is similar to the RIP constant in Section \ref{sec:sensing}.%) which we need to be smaller than $\frac{1}{12}$ (i.e., we need $\delta_2 < \frac{1}{12}$).

Similarly, by multiplying \eqref{eq:compvtupdate} with $v_{\perp}$ (where $\ip{v_{\perp}^*}{\vo}=0$ and $\|v_{\perp}^*\|_2=1$) and using Lemma~\ref{lem:errb_comp1}:
$$\|\uht\|_2\ip{\widehat{v}^{t+1}}{v_{\perp}^*}\leq 2\so\delta_2\sqrt{1-\ip{\ut}{\uo}^2}.$$
%Similarly, $\forall v_{\perp}^*$ where $\ip{v_{\perp}^*}{\vo}=0$ and $\|v_{\perp}^*\|_2=1$, $\ip{\widehat{v}^{t+1}}{v_{\perp}^*}\leq 2\so\delta_2\sqrt{1-\ip{\ut}{\uo}^2}$. 
Using the above two equations: 
\begin{equation*}1-\ip{\vt}{\vo}^2\leq \frac{4\delta_2^2( 1-\ip{\ut}{\uo}^2)}{(\ip{\ut}{\uo}-2\delta_2\sqrt{1-\ip{\ut}{\uo}^2})^2+(2\delta_2\sqrt{1-\ip{\ut}{\uo}^2})^2}.\end{equation*}
Assuming, $\ip{\vt}{\vo}\geq 6\delta_2$, $$\dist(\vt, \vo)=\sqrt{1-\ip{\vt}{\vo}^2}\leq \frac{1}{4}\sqrt{1-\ip{\ut}{\uo}^2}.$$
Using same arguments, we can show that, $\dist(u^{t+1}, \uo)\leq \dist(\vt, \vo)/4$. Hence, after $O(\log(1/\epsilon))$ iterations, $\dist(\ut, \uo)\leq \epsilon $ and $\dist(\vt, \vo)\leq \epsilon$. This proves our second step. 

We now provide the following lemma to prove the third step. We stress that $\vt$ does not increase the incoherence parameter ($\mu_1$) when compared to that of $\ut$. %We prove the third step, i.e., incoherence of $\vt$ assuming incoherence of $\ut$, in Lemma~\ref{lem:incoherence_vt}. 
\begin{lemma}  \label{lem:incoherence_vt}
Let $M$, $p$, $\Omega$ be as defined in Theorem~\ref{thm:sampling-altmin-main}. Also, let $\ut$ be a unit vector with incoherence parameter $\mu_1=\frac{6(1+\delta_2)\mu}{1-\delta_2}$. 
%Let $M=\sigma^*\uo\vo$ be $\mu$-incoherent  and let $\Omega$ be sampled uniformly and independently of $\ut, \uo, \vo$, with $p\geq \frac{C\mu_1^2\log n}{n\delta_2^2}$ ($C>0$ is a global constant). 
 Then, w.p. at least $1-\frac{1}{n^3}$, $\vt$ is also $\mu_1$ incoherent. 
\end{lemma}
See Appendix~\ref{app:rank1_comp} for a detailed proof of the lemma. 

Finally, for the base case we need that $u^0$ is $\mu_1$ incoherent and also $\ip{u^0}{\uo}\geq 6\delta_2$. This follows directly by using Lemma~\ref{lemma:svd-clipping-jointguarantee} and the fact that $\delta_2\leq 1/12$. 

Note that, to obtain an error of $\epsilon$, \amc needs to run for $O\left(\log \frac{\|M\|_F}{\epsilon}\right)$ iterations. Also, we need to sample a fresh $\Omega$ at each iteration of \amce. Hence, the total number of samples needed by \amc is $O\left(\log \frac{\|M\|_F}{\epsilon}\right)$ larger than the  number of samples required per step.
% However, above we assumed that each $\ut$ is incoherent and $\ip{\ut}{\uo}\geq 6\delta_2$. We prove incoherence of $\vht$ by induction, that is assming $\ut$ is incoherent we show $\vt=\vht/\|\vht\|_2$ is incoherent. A detailed proof for the same can be found in Lemma~\ref{lem:incoherence_vt}. Lemma~\ref{lemma:svd-clipping-jointguarantee} shows the base case of induction : incoherence of $u^0$ as well as $\ip{u^0}{\uo}\geq 6\delta_2$. $\ip{\ut}{\uo}\geq 6\delta_2$ follows directly by using the fact that $\ip{\ut}{\uo}$ increases monotonically and the fact that  $\ip{u^0}{\uo}\geq 6\delta_2$.
%\end{proof}
%We now provide a lemma to show incoherence of each $\vt$, assuming $\ut$ to be incoherent. We stress that $\vt$ does not increase the incoherence parameter ($\mu_1$) when compared to that of $\ut$.
%In the following lemmas, we assume that $\delta_2<1/12$ is a constant, independent of $m$, $n$.% (i.e., does not scale with $m$ or $n$) less than $\frac{1}{12}$.
%We now show that, assuming $\ut$ to be {\em incoherent},  the above update is able to decrease distance of $\vt$ to $\vo$ at a geometric rate. 
%\begin{lemma}
%Let observed samples $\Omega$ and underlying matrix $M=\sigma^*\uo(\vo)^\dag$ satisfy assumptions mentioned in Theorem~\ref{thm:sampling-altmin-geomconv}.% Also, let $\ut$ be a unit $L_2$ norm vector with incoherence parameter $\mu_1$. Then, 
%$$\dist(\vt, \vo)\leq \frac{1}{4}\dist(\ut, \uo).$$
%\end{lemma}
%\begin{proof}
%Using \eqref{eq:r1compupdate}, component of $\vt$ along $\vo$ is given by:
%\begin{align}
%  \ip{\vt}{\vo}=\ip{\ut}{\uo}-(\vo)^\dag B^{-1}\vo
%\end{align}
%\end{proof}
%We now present proof of Theorem~\ref{thm:sampling-altmin-geomconv} for the above rank-$1$ updates. 
%We now present a lemma to bound the error term in \eqref{eq:compvtupdate}. 
%We provide proofs of  the above lemmas in Appendix~\ref{app:rank1_comp}. 
% \begin{corollary}\label{coro:twonormdist-frobdist}
%   $\wh{U}^T(\wh{V}^T)^{\dag}$ returned by Algorithm \ref{algo:completion-altmin} satisfies:
% \begin{align*}
%   \frob{\wh{U}^T(\wh{V}^T)^{\dag} - U^*\Sigma^* (V^*)^{\dag} } \leq @@@
% \end{align*}
% \end{corollary}
% \begin{proof}
%   We can bound the frobenius norm of the error as follows:
% \begin{align*}
%   \frob{\wh{U}^T(\wh{V}^T)^{\dag} - U^*\Sigma^* (V^*)^{\dag} }^2
% 	&= \frob{\wh{U}^T(\wh{V}^T)^{\dag} - U^*\Sigma^* (V^*)^{\dag}\wh{V}^T\left(\wh{V}^T\right)^{\dag}
% 		+ U^*\Sigma^* (V^*)^{\dag}\wh{V}^T\left(\wh{V}^T\right)^{\dag} - U^*\Sigma^* (V^*)^{\dag} }^2 \\
% 	&= \frob{\left(\wh{U}^T - U^*\Sigma^* (V^*)^{\dag}\wh{V}^T\right)\left(\wh{V}^T\right)^{\dag}
% 		+ U^*\Sigma^* (V^*)^{\dag}\left(\wh{V}^T\left(\wh{V}^T\right)^{\dag}-I\right) }^2 \\
% 	&= \frob{\left(\wh{U}^T - U^*\Sigma^* (V^*)^{\dag}\wh{V}^T\right)\left(\wh{V}^T\right)^{\dag}}^2
% 		+ \frob{U^*\Sigma^* (V^*)^{\dag}\left(\wh{V}^T\left(\wh{V}^T\right)^{\dag}-I\right) }^2 \\
% 	&= \frob{\wh{U}^T - U^*\Sigma^* (V^*)^{\dag}\wh{V}^T}^2
% 		+ \frob{\Sigma^*\left(V^*\right)^{\dag}\left(\wh{V}^T\left(\wh{V}^T\right)^{\dag}-I\right) }^2 \\
% 	&\stackrel{(\zeta_1)}{\leq} \left(\frac{k\delta_2}{1-\delta_2}\right)^2 \dist\left(\wh{V}^T,V^*\right)^2
% 		+ k \left(\sigma_1^*\right)^2 \dist\left(\wh{V}^T,V^*\right)^2 \\
% 	&= \left(\left(\frac{k\delta_2}{1-\delta_2}\right)^2 
% 		+ k \left(\sigma_1^*\right)^2\right) \dist\left(V^T,V^*\right)^2 \\
% 	&\leq @@@,
% \end{align*}
% where $(\zeta_1)$ follows from \eqref{}.@@@
% \end{proof}
\subsection{Rank-$k$ case}\label{sec:comp_general}
We now extend our proof of Theorem~\ref{thm:sampling-altmin-geomconv}  to matrices with arbitrary rank. Here again, we show that the \amc algorithm reduces to power method with bounded perturbation at each step. %However, bounds on perturbation hold only when each iterate generated by \amc is incoherent. So, in addition to geometric decay in the distance to the optimal singular vectors, we also show incoherence of each iterate. 

Similar to the matrix sensing case, we analyze the following QR decomposition based update instead of directly analyzing the updates of Algorithm~\ref{algo:completion-altmin}:
\begin{align}
\widehat{U}^{t}&= \, U^{t}R_U^{t}\ \ \text{ (QR decomposition), }\nonumber\\
\widehat{V}^{t+1}&=\argmin_{\widehat{V}} \  \|P_{\Omega}(U^t\widehat{V}^\dag)-P_{\Omega}(M)\|_F^2,\nonumber\\ 
\widehat{V}^{t+1}&=V^{t+1}R_V^{t+1}.\ \ \text{ (QR decomposition) },\nonumber\\ 
\widehat{U}^{t+1}&=\argmin_{\widehat{U}} \  \|P_{\Omega}(\widehat{U}(V^{t+1})^\dag)-P_{\Omega}(M)\|_F^2.\label{eq:updatecomp_k1}
\end{align}
Here again, we would stress that the updates output exactly the same matrices at the end of each iteration and we prefer QR-based updates due to notational ease. 

Now, as matrix completion is a special case of matrix sensing, Lemma~\ref{lemma:sensing-updateeqn1} characterizes the updates of the \amc algorithm (see Algorithm~\ref{algo:completion-altmin}). That is, 
\begin{align}
  \widehat{V}^{t+1} &= \underbrace{V^*\Sigma^*{U^*}^\dag  U^{t}}_{\text{Power-method Update}} - \underbrace{F}_{\text{Error Term}}, \nonumber\\
  V^{t+1}&=\widehat{V}^{t+1}(R^{(t+1)})^{-1},\label{eq:updatecomp_k}
\end{align}
where $F$ is the error matrix defined in \eqref{eq:F_k} and $R^{(t+1)}$ is a upper-triangular matrix obtained using QR-decomposition of $\widehat{V}^{t+1}$. See \eqref{eq:BCD_k} for the definition of  $B, C$, $D$, and $S$. 

Also, note that for the special case of matrix completion, $B_{pq}, C_{pq}, 1\leq p,q\leq k$ are {\em diagonal matrices} with $$(B_{pq})_{jj}=\frac{1}{p}\sum_{i:(i,j)\in \Omega}\Ut_{ip}\Ut_{iq},\ \ \ (C_{pq})_{jj}=\frac{1}{p}\sum_{i:(i,j)\in \Omega}\Ut_{ip}\Uo_{iq}.$$ We use this structure to further simplify the update equation. We first define matrices $B^j, C^j, D^j\in \mathbb{R}^{k\times k}, 1\leq i\leq n$: {\small $$B^j=\frac{1}{p}\sum_{i:(i,j)\in \Omega} (\Ut)^{(i)}{(\Ut)^{(i)}}^\dag , \ \ C^j=\frac{1}{p}\sum_{i:(i,j)\in\Omega} (\Ut)^{(i)}{(\Uo)^{(i)}}^\dag,$$} and $D^j=(\Ut)^\dag \Uo$. 
 Using the above notation, \eqref{eq:updatecomp_k} decouples into $n$ equations of the form ($1\leq j\leq n$):
 \begin{align}\label{eqn:sampling-updateeqn-decoupled}
    (\Vt)^{(j)} = (\Vo)^{(j)} (D^j-(B^j)^{-1}(B^jD^j-C^j)) (R^{(t+1)})^{-1},
 \end{align}
 where $(\Vt)^{(j)}$ and $(\Vo)^{(j)}$ denote the $j^{\textrm{th}}$ rows of $\Vt$ and $\Vo$ respectively.

Using the above notation, we now provide a proof of Theorem~\ref{thm:sampling-altmin-geomconv} for the general rank-$k$ case. 
\begin{proof}[Proof of Theorem~\ref{thm:sampling-altmin-geomconv}]
Multiplying the update equation \eqref{eq:updatecomp_k} on the left by $\left(\Vo_{\perp}\right)^{\dag}$, we get:\\
$(\Vo_{\perp})^{\dag}\Vht=-(\Vo_{\perp})^{\dag}F (\Rt)^{-1}.$ That is, 
\begin{align*}\dist(V^*, V^{t+1})&=\|{V_{\perp}^*}^\dag V^{(t+1)}\|_2=\|{V_{\perp}^*}^\dag  F {R^{(t+1)}}^{-1}\|_2\\&\leq \|F(\Sigma^*)^{-1}\|_2\|\Sigma^*{R^{(t+1)}}^{-1}\|_2.\end{align*}
Now, similar to the sensing case (see Section~\ref{sec:sensing_k}) we break down our proof into the following two steps:
\begin{itemize}
  \item Bound $\twonorm{F(\Sigma^*)^{-1}}$  (Lemma \ref{lemma:Fbound_comp}) and
  \item Bound $\|\Sigma^*{R^{(t+1)}}^{-1}\|_2$, i.e., the minimum singular value of $\left(\Sigma^*\right)^{-1}R^{(t+1)}$ (Lemma \ref{lemma:Rbound_comp}).
\end{itemize}
Using Lemma~\ref{lemma:Fbound_comp} and Lemma~\ref{lemma:Rbound_comp}, w.p. at least $1-1/n^3$, 
\begin{equation*}\dist(V^*, V^{t+1})\leq \|F(\Sigma^*)^{-1}\|_2\|\Sigma^*{R^{(t+1)}}^{-1}\|_2 \leq  \frac{(\sigma^*_1/\sigma^*_k)k(\delta_{2k}/(1-\delta_{2k}))  \cdot \dist\left(U^{(t)},U^*\right)}{\sqrt{1-\dist\left(U^{(t)},U^*\right)^2} - 
			\frac{(\sigma^*_1/\sigma^*_k)k\delta_{2k}  \dist\left(U^{(t)},U^*\right)}{1-\delta_{2k}}}.\end{equation*}
Now, using Lemma~\ref{lemma:svd-clipping-jointguarantee} we get: $\dist(U^t, U^*)\leq \dist(U^0, U^*)\leq \frac{1}{2}$. By selecting $\delta_{2k}<\frac{\so_k}{12k\so_1}$, i.e., $p\geq \frac{C(\so_1)^2k^4\log n}{m(\so_k)^2}$ and using above two inequalities:
$$\dist(\Vt, \Vo)\leq \frac{1}{4}\dist(\Ut, \Uo).$$
Furthermore, using Lemma~\ref{lemma:completion-incoherence-Vt} we get that $\Vt$ is $\mu_1$ incoherent. 
Hence, using similar arguments as above, we also get:
$\dist(U^{t+1}, \Uo)\leq \left(\frac{1}{4}\right)\dist(\Vt, \Vo).$ %, \dist(V^{t+1}, \Vo)\leq \left(\frac{1}{4}\right)^t.$
%Thus, using induction, we get: 
%$$\dist(U^{t+1}, \Uo)\leq \left(\frac{1}{4}\right)^t, \ \ \dist(V^{t+1}, \Vo)\leq \left(\frac{1}{4}\right)^t.$$
%That is, after $O(\log (1/\epsilon))$ iterations, we get: $\dist(\Ut, \Uo)\leq \epsilon,\ \dist(\Vt, \Vo)\leq \epsilon$. Hence Proved.
%This proves Theorem \ref{thm:sampling-altmin-geomconv}.
%Moreover, using  \eqref{eq:updatecomp_k}, the residual after $t$-th step is given by:
%$$M-\Ut(\Vht)^\dag=(I-\Ut(\Ut)^\dag)M-\Ut F^\dag.$$
%That is, $$\|M-\Ut(\Vht)^\dag\|_F\leq \sqrt{k}\sigma_1^*\epsilon+\sqrt{k}\sigma_1^*\epsilon\leq 2\sqrt{k}\|M\|_F\epsilon,$$
%where the first inequality follows using $\dist(\Ut, \Vo)\leq\epsilon$ and Lemma~\ref{lemma:Fbound_comp}.
%Theorem \ref{thm:sampling-altmin-main} now follows by setting $\epsilon'=2\sqrt{k}\|M\|_F\epsilon$.
%We again stress that,  a fresh $\Omega$ is assumed to be sampled at each step of \amce. % - this is to ensure that the indices in $\Omega$ are independent of the currentiterates. 
% Since the  number of iterations is bounded by $O(\log(\|M\|_F/\epsilon))$, the  total number of samples needed increases by a factor of $O(\log(\|M\|_F/\epsilon))$.%which is necessary for our analysis. As noted above, since the number of iterations is $O\left(\log \frac{k\|M\|_F}{\epsilon}\right)$,
%the total number of samples needed is equal to the number of samples needed in each iteration multiplied by $O\left(\log \frac{k\|M\|_F}{\epsilon}\right)$.
%Now, $\|\frac{1}{p}P_\Omega(M(I-\Ut(\Ut)^\dag))\Ut\|_F\leq \sqrt{k}\|\frac{1}{p}P_\Omega(M(I-\Ut(\Ut)^\dag))\|_2\leq \sqrt{k}\log n\|
\end{proof}
We now provide lemmas required by our above given proof. See Appendix~\ref{app:compk} for detailed proof of each of the lemmas.% are given in Appendix~\ref{app:compk}. 

We first provide a lemma to bound incoherence of $\Vt$, assuming incoherence of $\Ut$. 
\begin{lemma}\label{lemma:completion-incoherence-Vt}
Let $M, \Omega, p$ be as defined in Theorem~\ref{thm:sampling-altmin-main}. Also, let $U^t$ be the $t$-th step iterate obtained by \eqref{eq:updatecomp_k1}. Let $U^t$  be $\mu_1=\frac{16\sigma_1^*\mu\sqrt{k}}{\sigma_k^*}$ incoherent. Then, w.p. at least $1-1/n^3$, iterate $V^{(t+1)}$ is also $\mu_1$ incoherent.% with parameter $\mu_1 \eqdef $.
\end{lemma}
We now bound the error term ($F$) in  AltMin update \eqref{eq:updatecomp_k}. 
\begin{lemma}\label{lemma:Fbound_comp}
  Let $F$ be the error matrix defined by \eqref{eq:F_k} (also see \eqref{eq:updatecomp_k}) and let $\Ut$ be a $\mu_1$-incoherent orthonormal matrix obtained after $(t-1)^{\textrm{th}}$ update. Also, let $M$, $\Omega$, and $p$ satisfy assumptions of Theorem~\ref{thm:sampling-altmin-main}. Then, w.p. at least $1-1/n^3$: 
$$\twonorm{F(\Sigma^*)^{-1}} \leq \frac{\delta_{2k} k }{1-\delta_{2k}}dist(U^t, U^*).$$
\end{lemma}
Next, we present a lemma to bound $\|(\Rt)^{-1}\|_2$. 
\begin{lemma}\label{lemma:Rbound_comp}
  Let $\Rt$ be the lower-triangular matrix obtained by QR decomposition of $\Vht$ ( see \eqref{eq:updatecomp_k}) and let $\Ut$ be a $\mu_1$-incoherent orthonormal matrix obtained after $(t-1)^\textrm{th}$ update. Also, let $M$ and $\Omega$ satisfy assumptions of Theorem~\ref{thm:sampling-altmin-main}. Then,
\begin{align}\label{eqn:completion-minsingval-Rt}
  &\|\Sigma^*(R^{(t+1)})^{-1}\|_2\leq \frac{\sigma^*_1/\sigma^*_k}{\sqrt{1-\dist^2\left(U^{(t)},U^*\right)}-\frac{(\sigma^*_1/\sigma^*_k)\delta_{2k}k\dist(U^{(t)},U^*)}{1-\delta_{2k}}}%\sigma_{\textrm{min}}\left(R^{(t)}\right) \geq \gamma \lambda - \twonorm{F}.
\end{align}
\end{lemma}
\begin{proof}
Lemma follows by exactly the same proof as that of Lemma~\ref{lemma:sensing-minsingval-Rt} for the matrix sensing case. 
\end{proof}
%We now bound $\|B^{-1}\|_2$ and $\|C^j\|_2$, which is required by our bound for $F$ as well as for our incoherence proof. 



%\begin{lemma}[Matrix Completion RIP]
%Let $M_1=U_1\Sigma_1V_1^\dag$ be a rank-$k$, $\mu$-incoherent matrix and similarly let $M_2=U_2\Sigma_2V_2^\dag$ be another rank-$k$, $\mu$-incoherent mat%rix. Also let $p$ be the probability of sampling an entry of $M_1, M_2$. Let, $\Omega$ be the set of sampled entries and satisfy assumptions given in Theorem~\ref{thm:sampling-altmin-geomconv}. Then, 
%$$(1-\delta_{2k})\|M_1-M_2\|_F^2\leq \frac{1}{p}\|P_\Omega(M_1-M_2)\|_F^2\leq (1+\delta_{2k})\|M_1-M_2\|_F^2. $$ 
%\end{lemma}
%\begin{proof}
%Note that, $\|P_\Omega(M_1-M_2)\|_F^2=\|P_\Omega(M_1-U_1U_1^\dag M_2-(I-U_1U_1^\dag) M_2V_2V_2^T-(I-U_1U_1^\dag)M_2(I-V_2V_2^\dag)\|_F^2=\|P_\Omega(M_1-U_1U_1^\dag M_2\|_F^2+\|P_\Omega((I-U_1U_1^\dag) M_2)\|_F^2-2\ip{M_1-U_1U_1^\dag M_2}{P_\Omega((I-U_1U_1^\dag) M_2)}.$  
%\end{proof}

% % % \begin{lemma}
% % % Let $V\in \mathbb{R}^{n\times k}$ be a $\mu$-{\bf incoherent} orthonormal matrix and let $U\in \mathbb{R}^{m\times k}$ be {\bf any} orthonormal matrix. Also let $p\geq \frac{C\mu^2k^2\log n}{n\delta_{2k}^2}$ be the probability of sampling an entry of $M$, where $C>0$ is a global constant. Let, $\Omega$ be the set of sampled entries and satisfy assumptions given in Theorem~\ref{thm:sampling-altmin-geomconv}. Then, 
% % % $$\|U\Sigma V^\dag\|_F^2-\delta_{2k}\|U \Sigma V^\dag\|_F^2\leq \frac{1}{p}\|P_\Omega(U\Sigma V^\dag)\|_F^2\leq \|U\Sigma V^\dag\|_F^2+\delta_{2k}\|U\Sigma V^\dag\|_F^2. $$
% % % \end{lemma}
% % % \begin{proof}
% % % Note that, $\|P_\Omega(U\Sigma V^\dag)\|_F^2=\|P_\Omega(\sum_{\ell=1}^k\sigma_\ell u_\ell (v_l)^\dag)\|_F^2=\sum_{\ell_1, \ell_2}\sigma_{\ell_1}\sigma_{\ell_2}\ip{u_{\ell_1}v_{\ell_1}^\dag}{P_\Omega(u_{\ell_2}v_{\ell_2}^\dag)}$. Hence, by using Lemma~\ref{lemma:incoh1} and \ref{lemma:incoh2}, we get: 
% % % $$\frac{1}{p}\|P_\Omega(U\Sigma V^\dag)\|_F^2\leq \sum_{\ell}\sigma_\ell^2 + \delta_2\sum_{\ell_1, \ell_2}\sigma_{\ell_1}\sigma_{\ell_2} \leq (1+\delta_{2k})\sum_{\ell}\sigma_\ell^2=(1+\delta_{2k})\|U\Sigma V^\dag\|_F^2,$$ where $k\delta_2\leq \delta_{2k}$ and $p\geq \frac{C\mu^2k^2\log n}{n\delta_{2k}^2}$. Lower bound on $\|P_\Omega(U\Sigma V^\dag)\|_F^2$ also follows similarly. 
% % % %Let $Z=\frac{1}{p}\sum_{ij}\delta_{ij}\ip{U^{(i)}\Sigma}{V^{*(j)}}^2,$ where $\delta_{ij}$ be a Bernoulli random variable that indicates membership of index $i:(i,j)\in \Omega$. That is, $\delta_{ij}=1$ w.p. $p$ and $0$ otherwise. Note that, $\E[Z]=\|U\Sigma(\Vo)^\dag\|_F^2.$ Also, $\E[Z^2]=\frac{1}{p}\sum_{ij}\ip{U^{(i)}\Sigma}{V^{*(j)}}^4\leq \frac{\mu^2k^2}{np}\sum_{ij}\ip{U^{(i)}}{V^{*(j)}}^2=\frac{\mu^2k^2}{np}\|U(\Vo)^\dag\|_F^2.$ Similarly, $\max_{ij}\ip{U^{(i)}}{V^{*(j)}}^2\leq \frac{\mu^2k^2}{n}$. Finally, using Bernstein's inequality: 
% % % %$$\Pr(|Z-\E[Z]|\geq \delta\|U(\Vo)^\dag\|_F^2)\leq \exp\left(-\frac{\delta^2\|U(\Vo)^\dag\|_F^2}{1+\delta}\right).$$
% % % \end{proof}
% % % \begin{lemma}
% % % Let $v\in \mathbb{R}^{n}$ be a unit $\mu$-{\bf incoherent} vector and let $u\in \mathbb{R}^{m}$ be {\bf any} $m$-dimensional unit vector. Also let $p\geq \frac{6\mu^2\log n}{n\delta_2^2}$ be the probability of sampling an entry of $M$, where $0\leq \delta_2\leq 1$ is any constant. Let, $\Omega$ be the set of indices of the sampled entries that satisfies assumptions given in Theorem~\ref{thm:sampling-altmin-geomconv}. Then, 
% % % $$(1-\delta_2)\leq \frac{1}{p}\|P_\Omega(uv^\dag)\|_F^2\leq (1+\delta_2).$$\label{lemma:incoh1}
% % % \end{lemma}
% % % \begin{proof}
% % %   Let $Z=\frac{1}{p}\|P_\Omega(uv^\dag)\|_F^2=\frac{1}{p}\sum_{ij}\delta_{ij}(u_iv_j)^2$ where $\delta_{ij}$ be a Bernoulli random variable that indicates membership of index $i:(i,j)\in \Omega$. That is, $\delta_{ij}=1$ w.p. $p$ and $0$ otherwise. Note that, $\E[Z]=\sum_{ij}(u_iv_j)^2=1$, $\E[Z^2]=\frac{1}{p}\sum_{ij}(u_iv_j)^4\leq \frac{\mu^2}{np}$, and $\max_{ij}|(u_iv_j)^2|\leq \frac{\mu^2}{n}$. Now, using Bernstein's inequality:
% % % $$\Pr(|Z-\E[Z]|\geq \delta_2)\leq \exp\left(-\frac{\delta_2^2}{\frac{\mu^2}{np}+\frac{\mu^2}{np}\delta_2}\right).$$
% % % Lemma now follows by setting $p\geq \frac{6\mu^2\log n}{n\delta_2^2}$. 
% % % \end{proof}
% % % \begin{lemma}\label{lemma:incoh2}
% % % Let $v_1, v_2\in \mathbb{R}^{n}$ be two unit $\mu$-{\bf incoherent} vector  and let $u_1, u_2 \in \mathbb{R}^{m}$ be {\bf any} $m$-dimensional unit vectors such that $\ip{u_1}{u_2}=0$. Also let $p\geq \frac{3\mu^2\log n}{n\delta_2^2}$ be the probability of sampling an entry of $M$, where $0\leq \delta_2\leq 1$ is any constant. Let, $\Omega$ be the set of indices of the sampled entries that satisfies assumptions given in Theorem~\ref{thm:sampling-altmin-geomconv}. Then, 
% % % $$|\frac{1}{p}\ip{P_\Omega(u_1v_1^\dag)}{u_2v_2^\dag}|\leq \delta_2.$$
% % % \end{lemma}
% % % \begin{proof}
% % % Let $Z=\frac{1}{p}\ip{P_\Omega(u_1v_1^\dag)}{u_2v_2^\dag}=\frac{1}{p}\sum_{ij}\delta_{ij}(u_1)_i(u_2)_i(v_1)_j(v_2)_j,$ where $\delta_{ij}$ is a Bernoulli random variable defined in the previous lemma. Note that, $\E[Z]=\sum_i(u_1)_i(u_2)_i\sum_j(v_1)_j(v_2)_j=0$, $\E[Z^2]=\frac{1}{p}\sum_{ij}(u_1)_i^2(u_2)_i^2(v_1)_j^2(v_2)_j^2\leq \frac{\mu^2}{np}$ (using incoherence of $v_1$). Also, $\max_{ij}|(u_1)_i(u_2)_i(v_1)_j(v_2)_j|\leq \frac{\mu^2}{n}$. Hence, , using Bernstein's inequality:
% % % $$\Pr(|Z-\E[Z]|\geq \delta_2)\leq \exp(-\frac{\delta_2^2}{\frac{\mu^2}{np}+\frac{\mu^2}{np}\delta_2})\leq \frac{1}{n^3}.$$


%%% Local Variables: 
%%% mode: latex
%%% TeX-master: "altmin"
%%% End: 

%%% Local Variables: 
%%% mode: latex
%%% TeX-master: "altmin"
%%% End: 

%\subsection{Noisy Case}
\label{sec:sensing_noisycase}
We now consider an extension of the matrix sensing problem where measurements can be corrupted arbitrarily using a bounded noise. That is, we 
%We now extend Theorem \ref{thm:sensing} in the presence of noise. In particular, we consider the case where we 
observe $b=\aff\left(M+N\right)$, where $N$ is the noise matrix. % s.t.,
%\begin{align}\label{eqn:matrix-noisy}
%  M = U^* \Sigma^* {V^*}^{\dag},\ \ \ 
%  N = \Un \Sn (\Vn)^{\dag},
%\end{align}
%nd $N=\Un \Sn (\Vn)^{\dag}$ is the SVD of the noise matrix $N$. 
For this noisy case as well, we show that \as recovers $M$ upto an additive approximation depending on the Frobenius norm of  $N$. 
\begin{theorem}\label{thm:sensing-noisycase}
 Let $M$ and $\calA(\cdot)$ be as defined in Theorem~\ref{thm:sensing_err}. Suppose, \as algorithm (Algorithm \ref{algo:sensing-altmin}) is supplied inputs $\aff$, $b=\aff(M+N)$, where $N$ is the noise matrix s.t. $\left\|N\right\|_F  < \frac{1}{100}\so_k$.  Then, after $T=4\log(2/\epsilon)$ steps, iterates $\widehat{U}^{T}$, $\widehat{V}^{T}$ of \as  satisfy: 
%   Let $M=U^*\Sigma^*V^{*^\dag }$ be a rank-$k$ matrix with non zero singular values $\sigma_1^*\geq \sigma_2^*\dots\geq \sigma_k^*$ and $N=\Un \Sn (\Vn)^{\dag}$ be a matrix with singular values $\sn_1 \geq \sn_2 \geq
%\cdots \sn_n$ such that $\left\|N\right\|_F  < \frac{1}{100}\sigma_k$. Also, let $\calA(\cdot):\rmn \rightarrow \mathbb{R}^d$ be a linear measurement
%operator that satisfies $2k$-RIP with RIP constant $\delta_{2k}<\frac{(\sigma^*_k)^2}{100k(\sigma_1^*)^2}$.  Suppose, AltMin algorithm (Algorithm \ref{algo:sensing-altmin}) is supplied inputs $\aff$, $b=\aff(M+N)$. Then, after $T=4\log(2/\epsilon)$ steps, iterates $U^{T}$, $V^{T}$ of the AltMin algorithm  satisfy: 
%with $\dist\left(U^{(0)},U^*\right) \leq \frac{1}{2}$. Then $U^{(t)}$ and $V^{(t)}$ satisfy:
\begin{align*}
  \dist\left(\Vw^{T},V^*\right) \leq \frac{10\|N\|_F}{\so_k} + \epsilon, \ 
  \dist\left(\Uw^{T},U^*\right) \leq \frac{10\|N\|_F}{\so_k} + \epsilon.
\end{align*}
See Definition~\ref{defn:dist} for definition of $\dist\left(U, W\right)$. 
%where $\alpha_V, \alpha_U < 1$ are some constants depending only on $\delta_{2k}$. 
% denotes the principal angle based distance between subspaces spanned by the columns of $U\in \mathbb{R}^{m\times k}$ and $W\in \mathbb{R}^{m\times k}$ (See Definition~\ref{defn:dist}).
\end{theorem}
%\begin{theorem}\label{thm:sensing-noisy-main}
%   Let $M$ and $\calA(\cdot)$ be as defined in Theorem~\ref{thm:sensing_err}. 
%=U^*\Sigma^*V^{*^\dag }$ be a rank-$k$ matrix with non zero singular values $\sigma_1^*\geq \sigma_2^*\dots\geq \sigma_k^*$ and $N=\Un \Sn (\Vn)^{\dag}$ be the noise matrix. Also, let $\calA(\cdot):\rmn \rightarrow \mathbb{R}^d$ be a linear measurement operator that satisfies $2k$-RIP with RIP constant $\delta_{2k}<\frac{(\sigma^*_k)^2}{100k(\sigma_1^*)^2}$.  
%Suppose, \as algorithm (Algorithm \ref{algo:sensing-altmin}) is supplied inputs $\aff$, $b=\aff(M+N)$. Then, after $T=4\log(2/\epsilon)$ steps, iterates $\widehat{U}^{T}$, $\widehat{V}^{T}$ of \as  satisfy: 
%$$\|M-\widehat{U}\widehat{V}^T\|_F\leq C \|N\|_F,$$
%where $C=(\frac{\so_1}{\so_k}+10\delta_{2k})$. 
%\end{theorem}
At a high level, our proof for noisy case follows closely, the exact case proof given in Section~\ref{sec:sensing}. That is, we show that the update of \as algorithm is similar to power-method type update but with two errors terms: one due to incomplete measurements and another due to the noise matrix. A detailed proof of Theorem \ref{thm:sensing-noisycase} can be found in Appendix~\ref{app:noisy_sensing}.%One key difference is that top singular vectors of $M+N$ matrix are not exactly same as the singular vectors of $M$. Hence, we need to appropriately bound the distance of recovered matrices $\widehat{U}, \widehat{V}$ to the true singular vectors. 
%with $\dist\left(U^{(0)},U^*\right) \leq \frac{1}{2}$. Then $U^{(t)}$ and $V^{(t)}$ satisfy:
%\begin{align*}
%  \dist\left(V^{T},V^*\right) \leq \frac{10\|N\|_F}{\so_k} + \epsilon, \ \ \ 
%  \dist\left(U^{T},U^*\right) \leq \frac{10\|N\|_F}{\so_k} + \epsilon.
%\end{align*}
%See Definition~\ref{defn:dist} for definition of $\dist\left(U, W\right)$. 
%\end{theorem}
%The following theorem shows that if the Frobenius norm
%(sum of squared singular values) of the noise matrix $N$ is small enough, then alternating minimization gives a good approximation
%of $M$.
% We are now ready to prove Theorem \ref{thm:sensing-noisycase}.
% \begin{proof}(Theorem \ref{thm:sensing-noisycase})
%   From \eqref{eqn:sensing-orthogonaliter-error-noisycase}, \eqref{eqn:sensing-twonorm-noise-noisycase}, \eqref{eqn:sensing-twonorm-F},
% \eqref{eqn:sensing-twonorm-G-noisycase} and \eqref{eqn:sensing-minsingval-Rt-noisycase} we
% obtain the following:
% \begin{align*}
%   \dist\left(V^{(t)},V^*\right) = \twonorm{{V_{\perp}^*}^\dag V^{(t)}}
% 	&\leq \twonorm{{V_{\perp}^*}^\dag \left(Y \Lambda X^{\dag}U^{(t)}{R^{(t)}}^{-1} - {V_{\perp}^*}^\dag  F {R^{(t)}}^{-1} - {V_{\perp}^*}^\dag  G {R^{(t)}}^{-1}\right)} \\
% 	&\leq \twonorm{V_{\perp}^*}\left(\twonorm{Y \Lambda X^{\dag}U^{(t)}} + \twonorm{F} + \twonorm{G}\right)\twonorm{{R^{(t)}}^{-1}} \\
% 	&\leq 1 \cdot \left(\lambda_1 + \frac{\delta k \Gamma \dist\left(U^{(t)},U^*\right)}{1-\delta} + \frac{\delta \frob{\Lambda}}{1-\delta} \right) \cdot \frac{1}{\sigma_{\textrm{min}}\left(R^{(t)}\right)} \\
% 	&\leq \frac{ \lambda_1 + \frac{\delta k \Gamma \dist\left(U^{(t)},U^*\right)}{1-\delta} + \frac{\delta \frob{\Lambda}}{1-\delta} }
% 	{ \gamma \lambda - \lambda_1 - \frac{\delta k \Gamma \dist\left(U^{(t)},U^*\right)}{1-\delta} - \frac{\delta \frob{\Lambda}}{1-\delta} }.
% \end{align*}
% @@@complete this. The theorem follows from the above equation.
% \end{proof}
%%% Local Variables: 
%%% mode: latex
%%% TeX-master: "altmin"
%%% End: 

%\newcommand{\Uw}{\widehat{U}}
%\newcommand{\Vw}{\widehat{V}}
\newcommand{\soj}{(\so_j)^2}
\newcommand{\soi}{(\so_i)^2}
\newcommand{\soii}{(\so_{i+1})^2}
\newcommand{\Up}{U_{\perp}}
\begin{algorithm}[t]
\caption{{\bf Stage-AltMin}: Stagewise Alternating Minimization for Matrix Sensing }
\begin{algorithmic}[1]
\STATE  Input: $b,\aff$
\STATE $\widehat{U}^T \leftarrow []$, $\widehat{V}^T \leftarrow []$
\FOR{$i=1,\cdots,k$}
%\STATE $\widehat{U}^0 \leftarrow [\widehat{U}^T \;0]$, $\widehat{V}^0 \leftarrow [\widehat{V}^T \;0]$  where $ \widehat{U}^{0}\in \mathbb{R}^{m\times i}$, $\widehat{V}^{0}\in \mathbb{R}^{n\times i}$
{\STATE {$[\widehat{U}^0_{1:i}\ \widehat{V}^0_{1:i}]=\mbox{top i-singular vectors of }$ $\left(\widehat{U}^T_{1:i-1} (\widehat{V}_{1:i-1}^T)^\dag-\frac{3}{4}\aff^T(\aff(\widehat{U}_{1:i-1}^T(\widehat{V}^T_{1:i-1})^\dag)-b)\right)$ i.e., one step of SVP \cite{JainMD10}\label{stp:sam_svp}}}%\textrm{One-Step-SVP}\left(\aff,b,\widehat{U}^{T} (\widehat{V}^T)^{\dag}, i\right)$
\FOR{$t=0,\cdots,T-1$}
	\STATE {$\widehat{V}_{1:i}^{t+1}\leftarrow \argmin_{V\in \R^{n\times i}} \  \|\aff(\widehat{U}^t_{1:i}V^\dag )-b\|_2^2$\label{stp:vt}}
	\STATE {$\widehat{U}_{1:i}^{t+1}\leftarrow \argmin_{U\in \mathbb{R}^{m\times i}} \  \|\aff(U_{1:i}(\widehat{V}_{1:i}^{t+1})^\dag )-b\|_2^2$\label{stp:ut}}
\ENDFOR
\ENDFOR
% \STATE $V^{t+1}\leftarrow QR(\widehat{V}^{t+1})\ \ \text{ (Orthogonalization using QR decomposition) }$
\STATE Output: $X=\widehat{U}_{1:i}^T(\widehat{V}_{1:i}^T)^\dag $
\end{algorithmic}\label{algo:sensing-altmin-improved}
\end{algorithm}
\section{Stagewise AltMin Algorithm}\label{sec:sam}
In Section~\ref{sec:sensing}, we showed that if $\delta_{2k}\leq \frac{(\so_k)^2}{(\so_1)^2k}$ then \as (Algorithm~\ref{algo:sensing-altmin}) recovers the underlying matrix. This means that, $d=\frac{(\so_1)^4}{(\so_k)^4}k^2n\log  n$ random Gaussian measurements (assume $m\leq n$) are required to recover $M$. For matrices with large condition number ($\so_1/ \so_k$), this would be significantly larger than the information theoretic bound of $O(k n \log n/k)$ measurements. %For matrices with large condition number ($\so_1/ \so_k$), this measurement requirement can be significantly larger than the information theoretically optimal bound $O(k n \log n)$  

To alleviate this problem, we present a modified version of \as called Stage-AltMin. Stage-AltMin proceeds in $k$ stages where in the $i$-th stage, a rank-$i$ problem is solved. The goal of the $i$-th stage is to recover top $i$-singular vectors of $M$, up to $O(\so_{i+1})$  error. 

Specifically, we initialize the $i$-th stage of our algorithm using one step of the SVP algorithm \cite{JainMD10} (see Step~\ref{stp:sam_svp} of Algorithm~\ref{algo:sensing-altmin-improved}).  We then show that, if $\delta_{2k}\leq \frac{1}{10k}$, then Stage-AltMin (Steps~\ref{stp:vt}, \ref{stp:ut} of Algorithm~\ref{algo:sensing-altmin-improved}) decreases the error $\|M-\widehat{U}^T_{1:i}(\widehat{V}^T_{1:i})^\dag\|_F$ to $O(\so_{i+1})$. Hence, after $k$ steps, the error decreases to $O(\so_{k+1})=0$. Note that, $\widehat{U}^t_{1:i}\in \mathbb{R}^{m\times i}$ represents the $t$-th step iterate ($U$) in the $i$-th stage; $\widehat{V}^t_{1:i}\in \mathbb{R}^{n\times i}$ is also defined similarly.

Recall that, the main problem with our analysis of \as is that if $\sigma_i\gg \sigma_{i+1}$ (for some $i$) then $\delta_{2k}\leq \frac{(\so_{i+1})^2}{(\so_i)^2k}$ would need to be  small. However, in such a scenario, %if 
%as our $i$-th stage solves only a rank-$i$ problem, hence, if 
%The key idea behind our analysis is that if 
%$\so_{i}$ is much larger than $\so_{i+1}$ then 
 the $i$-th stage of Algorithm~\ref{algo:sensing-altmin-improved} can be thought of as solving a noisy sensing problem where the goal is to recover $M_i\eqdef\Uo_{1:i}\So_{1:i}(\Vo_{1:i})^\dag$ using noisy measurements $b=\aff(\Uo_{1:i}\So_{1:i}(\Vo_{1:i})^\dag+N)$ where noise matrix $N\eqdef\Uo_{i+1:k}\So_{i+1:k}(\Vo_{i+1:k})^\dag$.
Here $M_i$ and $N$ represent the top $i$ singular components and last $k-i$ singular components of $M$ respectively.
Hence, using noisy-case type analysis (see Section~\ref{app:noisy_sensing}) we show that the error $\|M-\widehat{U}^t(\widehat{V}^t)^\dag\|_F$ decreases to $O(\so_{i+1})$. 

We now formally present the proof of our main result (see Theorem~\ref{thm:sam_err}). 
%\begin{theorem}\label{thm:sam_err}
%Let $M=U^*\Sigma^*V^{*^\dag }$ be a rank-$k$ matrix with non zero singular values $\sigma_1^*\geq \sigma_2^*\dots\geq \sigma_k^*$. Also, let $\calA(\cdot):\rmn \rightarrow \mathbb{R}^d$ be a linear measurement operator that satisfies $2k$-RIP with RIP constant $\delta_{2k}<\frac{1}{3200k^2}$.  Suppose, Stage-AltMin (Algorithm \ref{algo:sensing-altmin-improved}) is supplied inputs $\aff$, $b=\aff(M)$. Then, the $i$-th stage iterates $\Uw^T_{1:i}$, $V^T_{1:i}$ satisfy: $$\|M-\Uw^T_{1:i}V^T_{1:i}\|_F^2\leq \max(\epsilon, 16k(\so_{i+1})^2),$$% + (\so_i)^2\left(\frac{1}{4}\right)^t.$$
%That is, the $T$-th step iterates of the $k$-th stage, satisfy: $$\|M-\Uw^T_{1:k}V^T_{1:k}\|_F^2\leq \epsilon,$$
%where $T= \Omega\left(\log(\|M\|_F^2/\epsilon)\right)$. That is, the $T$-th step iterates of the $k$-th stage, satisfy: $$\|M-\Uw^T_{1:k}V^T_{1:k}\|_F^2\leq \epsilon.$$
%\end{theorem}
\begin{proof}[Proof Of Theorem~\ref{thm:sam_err}]
We prove the theorem using mathematical induction. 

\noindent {\bf Base Case}: After the $0$-th step, error is: $\|M\|_F^2\leq \sum_{j=1}^k \sigma_j^2\leq k\sigma_1^2$. Hence, base case holds. 

\noindent {\bf Induction Step}: Here, assuming that the error bound holds for $(i-1)$-th stage, we prove the error bound for the $i$-th stage. 

Our proof proceeds in two steps. First, we show that the initial point $\Uw^0_{1:i}$, $\Vw^0_{1:i}$ of the $i$-th stage, obtained using Step~\ref{stp:sam_svp}, has $c(\so_{i})^2+O\left(k(\so_{i+1})^2\right)$ error, with $c<1$. In the second step,  we show that using the initial points  $\Uw^0_{1:i}$, $\Vw^0_{1:i}$, the AltMin algorithm iterations in the $i$-th stage (Steps~\ref{stp:vt}, \ref{stp:ut}) reduces the error to $\max(\epsilon,16k\sigma_{i+1}^2)$. 

We formalize the above mentioned first step in Lemma~\ref{lemma:sam_svp_err} and then prove the second step in Lemma~\ref{lemma:sam_altmin}. 
%Using Lemma~\ref{lemma:sam_svp_err}, for $i$-th stage, 
%\begin{align}
%  \|M-\Uw^0_{1:i}\Vw^0_{1:i}\|_F^2\leq \sum_{j=i+1}^k (\so_j)^2 + \frac{1}{100} (\so_i)^2. \label{eq:sam_err_svp1}
%\end{align}
%Now, we consider two cases: a) $\frac{\so_i}{\so_{i+1}} < 5\sqrt{k}$, b) $\frac{\so_i}{\so_{i+1}} \geq 5\sqrt{k}$.\\
%{\bf Case (a)}: In this case, using \label{eq:sam_err_svp1} directly gives error bound, i.e, $\|M-\Uw^T_{1:i}\Vw^T_{1:i}\|_F^2=O(k\soii)$. \\
%{\bf Case (b)}: For this case, we can use the proof of Lemma~\ref{lemma:cond-num-bad-altmin-performance}. In particular, the desired result is obtained by using  \eqref{eq:sam_err_bc_t}. 
\end{proof}
We now present two lemmas used by the above given proof. See Appendix~\ref{app:stagewise_sensing} for a proof of each of the lemmas. 
\begin{lemma}\label{lemma:sam_svp_err}
Let assumptions of Theorem~\ref{thm:sam_err} be satisfied. Also, let $\Uw^0_{1:i}$, $\Vw^0_{1:i}$ be the output of Step~\ref{stp:sam_svp} of Algorithm~\ref{algo:sensing-altmin-improved}. Then, assuming that $\|M-\Uw^T_{1:i-1}\Vw^T_{1:i-1}\|_F^2\leq 16k(\so_{i})^2$, we obtain:
\begin{align*}
  \left\|M - \widehat{U}^0_{1:i}(\widehat{V}_{1:i}^0)^{\dag}\right\|_F^2 \leq \sum_{j=i+1}^k (\so_j)^2 + \frac{1}{100} (\so_{i})^2. 
\end{align*}
\end{lemma}
\begin{lemma}\label{lemma:sam_altmin}
Let assumptions of Theorem~\ref{thm:sam_err} be satisfied. Also, let $\Uw^T_{1:i}$, $\Vw^T_{1:i}$ be the $T$-th step iterates of the $i$-th stage of Algorithm~\ref{algo:sensing-altmin-improved}. Then, assuming that $\|M-\Uw^0_{1:i}V^0_{1:i}\|_F^2\leq \sum_{j=i+1}^k (\so_j)^2 + \frac{1}{100} (\so_{i})^2$, we obtain:
\begin{align*}
  \left\|M - \widehat{U}^T_{1:i}(\widehat{V}_{1:i}^T)^{\dag}\right\|_F^2 \leq \max(\epsilon, 16k(\so_{i+1})^2),
\end{align*}
where $T=\Omega(\log(\|M\|_F/\epsilon))$. 
\end{lemma}
%In this section, we will state and prove our main result which says that Algorithm \ref{algo:sensing-altmin-improved}
%has linear convergence to the optimum. The improvement in this result compared to that for Algorithm \ref{algo:sensing-altmin}
%is that our result in this section does not require $\delta$ to depend on the condition number of the underlying matrix.
% \begin{lemma}\label{lemma:cond-num-good}
% Assume conditions given in Theorem~\ref{thm:sam_err} are satisfied and  let $\frac{\so_i}{\so_{i+1}} < 5\sqrt{k}$. Also, let $\left\|M - \widehat{U}^0_{1:i}(\widehat{V}^0_{1:i})^{\dag}\right\|_F^2 \leq \sum_{j=i+1}^k (\so_j)^2 + \frac{1}{100} (\so_i)^2.$ Then, $\Uw^T_{1:i}$, $V^T_{1:i}$ satisfy:\begin{align*}
%   \left\|M - \widehat{U}^0_{1:i+1}(\widehat{V}^0_{1:i+1})^{\dag}\right\|_F^2 \leq \sum_{j=i+2}^k (\so_j)^2 + \frac{1}{100} (\so_{i+1})^2.
% \end{align*}
% %  Fix $i\in [k]$. Suppose in the $i^{\mbox{th}}$ iteration, after the execution of line $5$, we have the following:
% %Suppose further that $\frac{\sigma_i}{\sigma_{i+1}} < 5\sqrt{k}$. Then, in the $(i+1)^{\mbox{th}}$ iteration, after the
% %execution of line $5$, we have the following:
% %\begin{align*}
% %  \left\|M - \widehat{U}\widehat{V}^{\dag}\right\|_F^2 \leq \sum_{j=i+2}^k \sigma_j^2 + \frac{1}{100} \sigma_{i+1}^2.
% %\end{align*}
% \end{lemma}
% \begin{proof}
% %We prove the lemma using mathematical induction. That is, let the lemma is true for $i$-th stage. Then, we show that the Lemma holds for $i+1$-th stage as well. 
%
% Note that, the alternative minimization steps (Step~\ref{stp:vt}, \ref{stp:ut}) of Algorithm~\ref{algo:sensing-altmin-improved} does not increase the error. %$\left\|M - \widehat{U}^0_{1:i}(\widehat{V}^0_{1:i})^{\dag}\right\|_F^2 \leq \sum_{j=i+1}^k (\so_j)^2 + \frac{1}{100} (\so_i)^2$. 
% Hence, the iterates at the end of $i$-th stage, $\Uw^T_{1:i}$, $\Vw^T_{1:i}$ satisfy: 
% \begin{align*}
%   \left\|M - \widehat{U}^T_{1:i}(\widehat{V}^T_{1:i})^{\dag}\right\|_F^2 \leq \sum_{j=i+1}^k (\so_j)^2 + \frac{1}{100} (\so_i)^2.
% \end{align*}
% Then, after executing Step~\ref{stp:sam_svp} (Algorithm~\ref{algo:sensing-altmin-improved}) of the $(i+1)^{\mbox{th}}$ stage, we have:
% \begin{align*}
%   \left\|M - \widehat{U}^0_{1:i+1}(\widehat{V}^0_{1:i+1})^{\dag}\right\|_F^2 &\leq \left(1+\delta_{2k}\right) \sum_{j=i+2}^k (\so_j)^2 + \delta_{2k} \left(\sum_{j=i+1}^{k}(\so_{j})^2 + \frac{1}{100} \soi\right)\\
% &\leq \sum_{j=i+2}^k \soj + 2\delta_{2k}\left(\sum_{j=i+1}^k \soj + \frac{1}{200} \soi\right) \\
% 	&\leq \sum_{j=i+2}^k \soj + 2\delta_{2k}\left( k \soii + \frac{25k}{200} \soii\right) \leq \sum_{j=i+2}^k \soj +  \frac{1}{100}\soii,
% \end{align*}
% where the first inequality follows using Lemma 2.1 of \cite{JainMD10}. The third inequality follows using the fact that $\frac{\so_i}{\so_{i+1}} < 5\sqrt{k}$.% and $\delta < \frac{1}{200 k}$.
%
% Hence proved. 
% \end{proof}
%\begin{lemma}\label{lemma:cond-num-bad-altmin-performance}
%Assume conditions given in Theorem~\ref{thm:sam_err} are satisfied and  let $\frac{\so_i}{\so_{i+1}} \geq 5\sqrt{k}$. Also, let $$\left\|M - \widehat{U}^0_{1:i}(\widehat{V}^0_{1:i})^{\dag}\right\|_F^2 \leq \sum_{j=i+1}^k (\so_j)^2 + \frac{1}{100} (\so_i)^2.$$ Then, $U^T_{1:i}$, $V^T_{1:i}$ satisfy:$$\|M-\Uw^T_{1:i}V^T_{1:i}\|_F^2\leq \max(\epsilon, 16k(\so_{i+1})^2),$$
%  Fix $i\in [k]$. Suppose in the $i^{\mbox{th}}$ iteration, after the execution of line $5$, we have the following:
%\begin{align*}
%  \left\|M - \widehat{U}^0_{1:i+1}(\widehat{V}^0_{1:i+1})^{\dag}\right\|_F^2 \leq \sum_{j=i+2}^k (\so_j)^2 + \frac{1}{100} (\so_{i+1})^2.
%\end{align*}
%
%  Fix $i\in[k]$. Suppose in the $i^{\mbox{th}}$ iteration, after the execution of line $5$, we have the following:
%\begin{align*}
%  \left\|M - \widehat{U}\widehat{V}^{\dag}\right\|_F^2 \leq \sum_{j=i+1}^k \sigma_j^2 + \frac{1}{100} \sigma_i^2.
%\end{align*}
%Suppose further that $\frac{\sigma_i}{\sigma_{i+1}} \geq 5\sqrt{k}$. Then, at the end of lines $6-9$ in the
%$i^{\mbox{th}}$ iteration, we have the following:
%\begin{align*}
%  \left\|M - \widehat{U}\widehat{V}^{\dag}\right\|_F^2 \leq \sum_{j=i+2}^k \sigma_j^2 + 9k \sigma_{i+1}^2.
%\end{align*}
%Furthermore, in the $(i+1)^{\mbox{th}}$ iteration, after the
%execution of line $5$, we have the following:
%\begin{align*}
%  \left\|M - \widehat{U}\widehat{V}^{\dag}\right\|_F^2 \leq \sum_{j=i+2}^k \sigma_j^2 + \frac{1}{100} \sigma_{i+1}^2.
%\end{align*}
%\end{lemma}
%See Appendix~\ref{app:stagewise_sensing} for a detailed proof of the above lemma.
%From Lemmas \ref{lemma:cond-num-good} and \ref{lemma:cond-num-bad-altmin-performance}, we see that the following
%hypothesis holds after the execution of line $5$ in iteration $i$:
%\begin{align*}
%  \left\|M - \widehat{U}\widehat{V}^{\dag}\right\|_F^2 \leq \max\left(\epsilon,\sum_{j=i+2}^k \sigma_j^2 + \frac{1}{100} \sigma_{i+1}^2\right).
%\end{align*}
%This proves the convergence of Algorithm \ref{algo:sensing-altmin-improved}.
%%% Local Variables: 
%%% mode: latex
%%% TeX-master: "altmin"
%%% End: 

%\input{newcompletion}
\section{Summary and Discussion}

% {\bf \tt @@ points to make:

% 1. ours is first analysis

% 2. we prove success under standard conditions

% 3. computational complexity

% 4. appeal as a distributed / parallelizable algorithm

% 5. our analysis sheds light on alt min - e.g. initialization, stage wise

% }

Alternating minimization provides an empirically appealing and popular approach to solving several different low-rank matrix recovery problems. The main motivation, and result, of this paper was to provide the {\bf first theoretical guarantees on the global optimality of alternating minimization}, for matrix completion and the related problem of matrix sensing. We would like to note the following aspects of our results and proofs:% we like to note:
\begin{itemize}\setlength{\itemindent}{0pt}\setlength{\topsep}{-5pt}\setlength{\itemsep}{-3pt}
\item For both the problems, we show that alternating minimization recovers the true matrix under {\em similar problem conditions} (RIP, incoherence) to those used by existing algorithms (based on convex optimization or iterated SVDs); computationally, our results show faster convergence to the global optima, but with possibly higher statistical (i.e. sample) complexity.
\item We develop a new framework for analyzing alternating minimization for low-rank problems. Key observation of our framework is that for some problems (under standard problem conditions) alternating minimization can be viewed as a perturbed version of the power method. In our case, we can control the perturbation error based on the extent of RIP / incoherence demonstrated by the problem. This idea is likely to have applications to other similar problems where  trace-norm based convex relaxation techniques have rigorous theoretical results but  alternating minimization has enjoyed more empirical success. For example, robust PCA \cite{ChandrasekaranSPW11,CandesLMW11}, spectral clustering \cite{JalaliCSX11} etc. %in what instances of non-negative matrix factorization, k-means, etc. problems would our ``approximate power method" approach work ? Are these the problem instances of practical interest ?
\item Our analysis also sheds light on two key aspects of the  alternating minimization approach:\\
%\begin{itemize}\setlength{\itemsep}{-3pt}\setlength{\topsep}{-5pt}
%  \item	
{\bf Initialization}: Due to its connection  to power method, it is now easy to see  that for alternating minimization to succeed, the initial iterate should not be orthogonal to the target vector. Our results indeed show that alternating minimization  succeeds if the
initial iterate is not ``almost orthogonal'' to the target subspace. This suggests that, selecting initial iterate smartly is preferable to random initialization.\\ %critical to the success of alternating minimization. %o choose the initial iterate smartly instead of choosing it randomly.
{\bf Dependence on the condition number}: Our results for the alternating minimization algorithm depend on the condition
number. However, using a stagewise adaptation of alternating minimization, we can remove this dependence
for the matrix sensing problem. This suggests that (problem specific) modifications of the basic alternating minimization
algorithm may in fact perform better than the original one, while (mostly) retaining the computational / implementational
simplicity of the underlying method.
%\end{itemize}
\end{itemize}



% \subsection{Matrix Sensing}
% \textbf{QR decomposition}: Algorithms \ref{algo:sensing-altmin} and \ref{algo:sensing-altmin-improved} do QR decomposition after each step of alternating minimization.
% This is infact not necessary and is useful only in simplifying the analysis. To see this, consider the following updates:
% \begin{align*}
%   \Vht &= \argmin_{\widehat{V}\in \R^{n\times k}} \twonorm{\aff\left(\Ut\left(\widehat{V}\right)^{\dag}\right)-b} \\
%   \widetilde{V}^{(t+1)} &= \argmin_{\widehat{V}\in \R^{n\times k}} \twonorm{\aff\left(\Uht\left(\widehat{V}\right)^{\dag}\right)-b},
% \end{align*}
% Since $\Uht = \Ut R^{(t)}$, where $R^{(t)}$ is invertible, we have $\widetilde{V}^{(t+1)} = \Vht \left(R^{(t)}\right)^{-1}$ and hence
% $\widetilde{V}^{(t+1)}$ and $\Vht$ span the same subspace. This implies that $\dist\left(\widetilde{V}^{(t+1)},V^*\right) = \dist\left(\Vht,V^*\right)$ and hence
% $\dist\left(\widetilde{V}^{(t+1)},V^*\right)$ still converges to zero geometrically.\\
% \textbf{Computational Complexity}: One of the most attractive features of alternating minimization is its computational efficiency. Each step of alternating
% minimization corresponds to solving a least squares problem with $nk$ variables and $d$ observations - the complexity of this is $O\left(k^2n^2d\right)$.
% Since there are $\log\frac{1}{\epsilon}$ such steps and $d=O\left(\frac{kn\log n}{\delta_{2k}^2}\right)$, the computational complexity of all the alternating
% minimization steps is $O\left(\frac{k^3n^3\log\frac{n}{\epsilon}}{\delta_{2k}^2}\right)$.\\
% \underline{Algorithm \ref{algo:sensing-altmin}}: Since we need $\delta_{2k} < \frac{\sigma_k^*}{k \sigma_1^*}$, the computational complexity of Algorithm \ref{algo:sensing-altmin}
% is $O\left(\left(\frac{\sigma_1^*}{\sigma_k^*}\right)^2k^5n^3\log\frac{n}{\epsilon}\right)$.
% \underline{Algorithm \ref{algo:sensing-altmin-improved}}: Since we need $\delta_{2k} < \frac{1}{k^2}$, the computational complexity of Algorithm \ref{algo:sensing-altmin-improved}
% is $O\left(k^7n^3\log\frac{n}{\epsilon}\right)$.
% 
% @@@Add comparison to \cite{RechtFP2007} and \cite{JainMD10}.
% 
% \subsection{Matrix Completion}
% \textbf{QR decomposition}: Similar to the algorithms in matrix sensing, Algorithm \ref{algo:completion-altmin} also does QR decomposition after each step of alternating minimization.
% Even in this case, using the same argument as above, we can see that the geometric decay of distance holds for iterates in the absence of QR decomposition.\\
% \textbf{Computational Complexity}: The computational complexity of alternating minimization for matrix completion is much better than that for matrix sensing because of the
% simplicity of the least squares problems encountered : each alternating minimization step decouples into $n$ least squares problems each with $k$ variables and
% $np = O\left(\left(\frac{\sigma_1^*}{\sigma_k^*}\right)^4 k^{4.5} \log n\right)$ and hence has computational complexity $O\left(\left(\frac{\sigma_1^*}{\sigma_k^*}\right)^4 k^{6.5} n\log n\right)$.
% Since the number of alternating minimization steps needed is $O\left(\frac{1}{\epsilon}\right)$ and the computational complexity of computing the top-$k$ singular components of (the sparse matrix)
% $P_{\Omega_0}(M)$ is $\order{kmnp}= \order{\left(\frac{\sigma_1^*}{\sigma_k^*}\right)^4 k^{5.5} n\log n}$, the computational complexity of Algorithm \ref{algo:completion-altmin} is
% $\order{\left(\frac{\sigma_1^*}{\sigma_k^*}\right)^4 k^{6.5} n\log \frac{n}{\epsilon}}$.
% @@Add comparison to \cite{CandesT2009} and \cite{KeshavanOM2009}.
%%% Local Variables: 
%%% mode: latex
%%% TeX-master: "altmin"
%%% End: 

%\section{Conclusion}



%
% The following two commands are all you need in the
% initial runs of your .tex file to
% produce the bibliography for the citations in your paper.
% \bibliographystyle{abbrv}
% \bibliography{refs}%\balancecolumns
\begin{thebibliography}{10}

\bibitem{CaiCS2008}
J.-F. Cai, E.~J. Cand\`{e}s, and Z.~Shen.
\newblock A singular value thresholding algorithm for matrix completion.
\newblock {\em SIAM Journal on Optimization}, 20(4):1956--1982, 2010.

\bibitem{CandesLMW11}
E.~J. Cand{\`e}s, X.~Li, Y.~Ma, and J.~Wright.
\newblock Robust principal component analysis?
\newblock {\em J. ACM}, 58(3):11, 2011.

\bibitem{CandesR2007}
E.~J. Cand{\`e}s and B.~Recht.
\newblock Exact matrix completion via convex optimization.
\newblock {\em Foundations of Computational Mathematics}, 9(6):717--772,
  December 2009.

\bibitem{CandesT2005}
E.~J. Cand{\`e}s and T.~Tao.
\newblock Decoding by linear programming.
\newblock {\em IEEE Transactions on Information Theory}, 51(12):4203--4215,
  2005.

\bibitem{CandesT2009}
E.~J. Cand{\`e}s and T.~Tao.
\newblock The power of convex relaxation: Near-optimal matrix completion.
\newblock {\em IEEE Trans. Inform. Theory}, 56(5):2053--2080, 2009.

\bibitem{ChandrasekaranSPW11}
V.~Chandrasekaran, S.~Sanghavi, P.~A. Parrilo, and A.~S. Willsky.
\newblock Rank-sparsity incoherence for matrix decomposition.
\newblock {\em SIAM Journal on Optimization}, 21(2):572--596, 2011.

\bibitem{chen2012matrix}
C.~Chen, B.~He, and X.~Yuan.
\newblock Matrix completion via an alternating direction method.
\newblock {\em IMA Journal of Numerical Analysis}, 32(1):227--245, 2012.

\bibitem{GolVan96}
G.~H. Golub and C.~F.~V. Loan.
\newblock {\em Matrix Computations}.
\newblock The Johns Hopkins University Press, 3rd edition, 1996.

\bibitem{HsiehCD12}
C.-J. Hsieh, K.-Y. Chiang, and I.~S. Dhillon.
\newblock Low rank modeling of signed networks.
\newblock In {\em KDD}, pages 507--515, 2012.

\bibitem{JainMD10}
P.~Jain, R.~Meka, and I.~S. Dhillon.
\newblock Guaranteed rank minimization via singular value projection.
\newblock In {\em NIPS}, pages 937--945, 2010.

\bibitem{JalaliCSX11}
A.~Jalali, Y.~Chen, S.~Sanghavi, and H.~Xu.
\newblock Clustering partially observed graphs via convex optimization.
\newblock In {\em ICML}, pages 1001--1008, 2011.

\bibitem{KeshavanOM2009}
R.~H. Keshavan, A.~Montanari, and S.~Oh.
\newblock Matrix completion from a few entries.
\newblock {\em IEEE Transactions on Information Theory}, 56(6):2980--2998,
  2010.

\bibitem{KimPar08}
H.~Kim and H.~Park.
\newblock Nonnegative matrix factorization based on alternating nonnegativity
  constrained least squares and active set method.
\newblock {\em SIAM J. Matrix Anal. Appl.}, 30(2):713--730, July 2008.

\bibitem{KimPark08b}
J.~Kim and H.~Park.
\newblock Sparse nonnegative matrix factorization for clustering.
\newblock Technical Report GT-CSE-08-01, Georgia Institute of Technology, 2008.

\bibitem{Koren09}
Y.~Koren.
\newblock The {B}ell{K}or solution to the {N}etflix grand prize, 2009.

\bibitem{KorenBV09}
Y.~Koren, R.~M. Bell, and C.~Volinsky.
\newblock Matrix factorization techniques for recommender systems.
\newblock {\em IEEE Computer}, 42(8):30--37, 2009.

\bibitem{LeeB10}
K.~Lee and Y.~Bresler.
\newblock Admira: atomic decomposition for minimum rank approximation.
\newblock {\em IEEE Transactions on Information Theory}, 56(9):4402--4416,
  2010.

\bibitem{MekaJCD2008}
R.~Meka, P.~Jain, C.~Caramanis, and I.~S. Dhillon.
\newblock Rank minimization via online learning.
\newblock In {\em ICML}, pages 656--663, 2008.

\bibitem{RechtFP2007}
B.~Recht, M.~Fazel, and P.~A. Parrilo.
\newblock Guaranteed minimum-rank solutions of linear matrix equations via
  nuclear norm minimization.
\newblock {\em SIAM Review}, 52(3):471--501, 2010.

\bibitem{Zangwill69}
W.~I. Zangwill.
\newblock {\em Nonlinear Programming: A Unified Approach}.
\newblock Englewood Cliffs: Prentice-Hall, 1969.

\bibitem{Zhouetal08}
Y.~Zhou, D.~M. Wilkinson, R.~Schreiber, and R.~Pan.
\newblock Large-scale parallel collaborative filtering for the netflix prize.
\newblock In {\em AAIM}, pages 337--348, 2008.

\bibitem{ZouHT06}
H.~Zou, T.~Hastie, and R.~Tibshirani.
\newblock Sparse principal component analysis.
\newblock {\em JCGS}, 15(2):262--286, 2006.

\end{thebibliography}

  % sigproc.bib is the name of the Bibliography in this case
% You must have a proper ".bib" file
%  and remember to run:
% latex bibtex latex latex
% to resolve all references
%
% ACM needs 'a single self-contained file'!
%
%APPENDICES are optional
%\balancecolumns
\onecolumn
\appendix
\section{Preliminaries}
\begin{lemma}[Lemma 2.1 of \cite{JainMD10}]\label{lem:jmd}
  Let $b=\aff(M)+e$, where $e$ is a bounded error vector, $M$ is a rank-$k$ matrix and $\aff$ is a linear measurement operator that satisfies $2k$-RIP with constant $\delta_{2k}$ (assume $\delta_{2k}<1/3$). Let $X^{t+1}$ be the $t+1$-th step iterate of SVP, then the following holds:
$$\|\aff(X^{t+1})-b\|_2^2\leq \|\aff(M)-b\|_2^2 + 2\delta_{2k}\|\aff(X^t)-b\|_2^2.$$
\end{lemma}
In our analysis, we heavily use the following two results. The first result is the well-known  Bernstein's inequality.
\begin{lemma}\label{lemma:bernstein}[Bernstein's inequality]
  Let $X_1,X_2,\cdots,X_n$ be independent random variables. Also, let $\left|X_i\right| \leq L \in \R\; \forall \; i$ w.p. $1$. Then, we have the following inequality:{\small
\begin{align}\label{eqn:bernstein}
  \prob{\left|\sum_{i=1}^n X_i - \sum_{i=1}^n \expec{X_i}\right| > t} \leq 2\exp\left(\frac{-t^2/2}{\sum_{i=1}^n \Var\left(X_i\right) + Lt / 3}\right).
\end{align}}
\end{lemma}
The second result is a restatement of Theorem 3.1 from \cite{KeshavanOM2009}.% (using our notation).
\begin{theorem}\label{thm:KOM-firststepSVD}(Restatement of Theorem 3.1 from \cite{KeshavanOM2009})
  Suppose $M$ is an incoherent rank-$k$ matrix and let $p,\Omega$ be as in Theorem~\ref{thm:sampling-altmin-main}. Further, let $M_k$ be the best rank-$k$ approximation  of $\frac{1}{p}P_{\Omega}\left(M\right)$. Then, w.h.p. we have:
\begin{align}\label{eqn:KOM-firststepSVD}
  \twonorm{M - M_k} \leq C \sqrt{\frac{k}{p\sqrt{mn}}} \frob{M}.
\end{align}
\end{theorem}
\textbf{Remark}: Note that Theorem 3.1 from \cite{KeshavanOM2009} holds only for $T_r\left(P_{\Omega}(M)\right)$  where
$T_r\left(P_{\Omega}(M)\right)$ is a trimmed version of $P_{\Omega}(M)$ obtained by setting all rows and columns of $P_{\Omega}(M)$ with too many observed entries  to zero.
However, using standard Chernoff bound we can argue that for our choice of $p$, none of the rows and columns of $P_{\Omega}(M)$ have too many observed entries and hence $T_r\left(P_{\Omega}(M)\right)=P_{\Omega}(M)$, whp.
\section{Matrix Sensing: Proofs}\label{app:sensing}
The following is an alternate characterization of RIP that we use heavily in our proofs.
At a conceptual level, it says that if $\aff$ satisfies RIP, then it also preserves inner-product between any two rank-$k$ matrices (upto some additive error).

\begin{lemma}\label{lem:alt_rip}
Suppose $\calA(\cdot)$ satisfies $2k$-RIP with constant $\delta_{2k}$. Then, for any $U_1,U_2\in \R^{m\times k}$ and
$V_1,V_2\in \R^{n\times k}$, we have the following:
\begin{align}\label{eqn:sensing-2k-RIP2}
  \left| \left\langle \calA\left(U_1V_1^\dag \right), \calA\left(U_2V_2^\dag \right)\right\rangle  - \Tr \left(U_2^\dag U_1V_1^\dag V_2\right) \right|
    \leq 3\delta_{2k} \left\|U_1V_1^\dag \right\|_F \left\|U_2V_2^\dag \right\|_F
\end{align}
\end{lemma}
\begin{proof}
  Consider the matrices $X_1\eqdef U_1V_1^T$, $X_2\eqdef U_2V_2^T$ and $X=X_1+X_2$.
Since the rank of $X$ is at most $2k$, we obtain the following using the RIP of $\calA$:
\begin{align*}
  (1-\delta) \left\|U_1V_1^T + U_2V_2^T\right\|_F^2 \leq \left\|\calA(X)\right\|_2^2
	\leq (1+\delta) \left\|U_1V_1^T+U_2V_2^T\right\|_F^2.
\end{align*}
Concentrating on the second inequality, we obtain
\begin{align}
  &\sum_i \left(\Tr\left(A_i U_1V_1^T\right)+\Tr\left(A_i U_2V_2^T\right)\right)^2
	\leq (1+\delta) \left(\left\|U_1V_1^T\right\|_F^2 + \left\|U_2V_2^T\right\|_F^2 +
		\Tr\left(U_1V_1^TV_2U_2^T\right)\right) \nonumber \\
  \stackrel{(\zeta_1)}{\Rightarrow} &\sum_i \Tr\left(A_i U_1V_1^T\right)\Tr\left(A_i U_2V_2^T\right) - \Tr\left(U_1V_1^TV_2U_2^T\right)
	\leq \delta\left(\left\|U_1V_1^T\right\|_F^2 + \left\|U_2V_2^T\right\|_F^2 +
		\Tr\left(U_1V_1^TV_2U_2^T\right)\right) \nonumber \\
  \stackrel{(\zeta_2)}{\Rightarrow} &\sum_i \Tr\left(A_i U_1V_1^T\right)\Tr\left(A_i U_2V_2^T\right) - \Tr\left(U_1V_1^TV_2U_2^T\right)
	\leq \delta\left(\left\|U_1V_1^T\right\|_F^2 + \left\|U_2V_2^T\right\|_F^2 +
		\left\|U_1V_1^T\right\|_F \left\|U_2V_2^T\right\|_F \right) \label{eqn:sensing-RIP-step1}
\end{align}
where $(\zeta_1)$ follows from the fact that $X_1$ and $X_2$ are rank-$k$ matrices and hence $\calA(\cdot)$ satisfies RIP w.r.t.
those matrices and $(\zeta_2)$ follows from the fact that $\Tr\left(U_1V_1^TV_2U_2^T\right) \leq 
\left\|U_1V_1^T\right\|_F \left\|U_2V_2^T\right\|_F$. Note that if we replace $U_1V_1^T$ by $\lambda U_1V_1^T$
and $U_2V_2^T$ by $\frac{1}{\lambda} U_2V_2^T$
in \eqref{eqn:sensing-RIP-step1} for some non-zero $\lambda \in \R$, the LHS of \eqref{eqn:sensing-RIP-step1} does not
change where as the RHS of \eqref{eqn:sensing-RIP-step1} changes. Optimizing the RHS w.r.t. $\lambda$, we obtain
\begin{align*}
  \sum_i \Tr\left(A_i U_1V_1^T\right)\Tr\left(A_i U_2V_2^T\right) - \Tr\left(U_2^TU_1V_1^TV_2\right)
	\leq 3\delta \left\|U_1V_1^T\right\|_F \left\|U_2V_2^T\right\|_F.
\end{align*}
A similar argument proves the other side of the inequality. This proves the lemma.
\end{proof}


\begin{proof}[Proof of Lemma~\ref{lemma:sensing-updateeqn1}]
We first show that the update \eqref{eq:update_k} reduces to:
\begin{align}\label{eqn:sensing-updateeqn}
  \sum_{q=1}^k \left(\sum_{i=1}^s A_i u_p^{(t)} u_q^{(t)^\dag } A_i^\dag \right)\widehat{v}_q^{(t+1)}
	= \sum_{q=1}^k \left(\sum_{i=1}^s A_i u_p^{(t)} u_q^{*^\dag } A_i^\dag \right) v_q^* \;\;\;\;\; \forall \; p\in[k].
\end{align}
  Let $Err(V)\eqdef \sum_i \left(\Tr\left(A_i M\right) - \Tr\left(A_i U^{(t)}V^\dag \right)\right)^2$. Since $\widehat{V}^{(t+1)}$
minimizes $E(V)$, we have $\nabla_{V} E(\widehat{V}^{(t+1)})=0$.
\begin{align*}
  &\nabla_{v_p} Err(\widehat{V}^{(t+1)})=0 \\
  \Rightarrow &\sum_{i=1}^s \left(\sum_{l=1}^k {v_q^{(t)}}^\dag  A_i u_q^{(t)} - \sum_{l=1}^k \sigma_q^*{v_q^*}^\dag  A_i u_q^* \right)
	A_i u_p =0 \\
  \Rightarrow &\sum_{l=1}^k \sum_{i=1}^s A_i u_p \left({v_q^{(t+1)}}^\dag  A_i u_q^{(t)}\right)
	= \sum_{l=1}^k \sum_{i=1}^s A_i u_p \left(\sigma_q^*{v_q^*}^\dag  A_i u_q^*\right) \\
  \Rightarrow &\sum_{l=1}^k \sum_{i=1}^s A_i u_p \left({u_q^{(t)}}^\dag  A_i^\dag  v_q^{(t+1)}\right)
	= \sum_{l=1}^k \sum_{i=1}^s A_i u_p \left({u_q^*}^\dag  A_i^\dag  \sigma_q^*v_q^*\right) \\
  \Rightarrow &\sum_{l=1}^k \left(\sum_{i=1}^s A_i u_p {u_q^{(t)}}^\dag  A_i^\dag \right) v_q^{(t+1)}
	= \sum_{l=1}^k \left(\sum_{i=1}^s A_i u_p {u_q^*}^\dag  A_i^\dag \right) \sigma_q^*v_q^*
\end{align*}
Define $$S=\left[\begin{matrix}\sigma_1^*I_{n}&\dots&0_{n}\\\vdots&\vdots&\vdots\\0_n&\dots&\sigma_k^*I_n\end{matrix}\right],\ \ v^*=\left[\begin{matrix}v_1^*\\\vdots\\v_k^*\end{matrix}\right],\mbox{ and }\ \  \widehat{v}_1^{(t+1)}=\left[\begin{matrix}\widehat{v}_1^{(t+1)}\\\vdots\\\widehat{v}_k^{(t+1)}\end{matrix}\right].$$ 
Then, 
\begin{align*}
   \widehat{v}_1^{(t+1)}&= B^{-1}CSv^* \\
	&= DSv^* - B^{-1}\left(BD-C\right)Sv^*
\end{align*}
where inverting $B$ is valid since the minimum singular value of $B$ is strictly positive (please refer Lemma
\ref{lemma:sensing-minsingval-B}). Considering the $p^{\textrm{th}}$ block of $\widehat{v}^{(t)}$, we obtain
\begin{align*}
  \widehat{v}_p^{(t+1)} &= \left(\sum_q \iprod{u_p^{(t)}}{u_q^*} \sigma_q^*v_q^*\right) - \left(B^{-1}\left(BD-C\right)Sv^*\right)_p \\
	&= \left(\sum_q \sigma_q^*v_q^* {u_q^*}^\dag  \right)u_p^{(t)} - \left(B^{-1}\left(BD-C\right)Sv^*\right)_p.
\end{align*}
This gives us the following equation for $\widehat{V}^{(t)}$:
\begin{align*}
  \widehat{V}^{(t+1)} = V^*\Sigma^*{U^*}^\dag  U^{(t)} - F
\end{align*}
where $F =\left[ \begin{array}{cccc}
\left(B^{-1}\left(BD-C\right)Sv^*\right)_1 & \left(B^{-1}\left(BD-C\right)Sv^*\right)_2
& \cdots & \left(B^{-1}\left(BD-C\right)Sv^*\right)_k
\end{array} \right].$
 
Hence Proved.
\end{proof}

\subsection{Rank-$1$ Matrix Sensing: Proofs}\label{app:sensing1}
\begin{proof}[Proof of Lemma~\ref{lem:error_1}]
Using definition of the spectral norm: 
\begin{equation}\label{eq:error1_l2}\|B^{-1}\left(\ip{\uo}{\ut}B-C\right)\vo\|\leq \|B^{-1}\|_2\cdot \|\ip{\uo}{\ut}B-C\|_2\cdot \|\vo\|_2.\end{equation}
Consider $B=\sum_i A_i \ut \utt \Ait$. Now, smallest eigenvalue of $B$, i.e., $\lambda_{min}(B)$ is given by: 
\begin{align}
\lambda_{min}(B)&=\min_{\|z\|=1}z^\dag B z=\min_{\|z\|=1} \sum_i z^\dag A_i \ut \utt \Ait z=\min_{\|z\|=1} \sum_i \Tr(A_i \ut z^\dag) \Tr(A_i \ut z^\dag),\nonumber\\
&=\min_{\|z\|=1} \ip{\aff(\ut z^\dag)}{\aff(\ut z^\dag)}\geq 1-3\delta_{2},\label{eq:bmin}
\end{align}
where the last inequality follows using Lemma~\ref{lem:alt_rip}. 
Using \eqref{eq:bmin}, \begin{equation}\|B^{-1}\|_2\leq \frac{1}{1-3\delta_2}.\label{eq:bimax}\end{equation}
\noindent Now, consider $G=\ip{\uo}{\ut}B-C=\sum_i A_i \left(\ip{\uo}{\ut}\ut\utt-\ut\uot\right)\Ait=\sum_i A_i \ut\left(\ip{\uo}{\ut}\ut-\uo\right)^\dag\Ait$. Using definition of the spectral norm: 
\begin{align}
  \|G\|_2&=\max_{\|z\|=1, \|y\|=1} z^\dag G y, \nonumber\\
&=\max_{\|z\|=1, \|y\|=1} \sum_i z^\dag A_i \ut \left(\ip{\uo}{\ut}\ut-\uo\right)^\dag\Ait y, \nonumber\\
&=\max_{\|z\|=1, \|y\|=1} \ip{\aff(\ut z^\dag)}{\aff\left( \left(\ip{\uo}{\ut}\ut-\uo\right) y^\dag\right)}, \nonumber\\
&\leq 3\delta_2\sqrt{1-\ip{\ut}{\uo}^2},\label{eq:gmax}
\end{align}
where the last inequality follows by using Lemma~\ref{lem:alt_rip} and the fact that $\ip{\ut}{\left(\ip{\uo}{\ut}\ut-\uo\right)}=0$. 

Lemma now follows using \eqref{eq:error1_l2}, \eqref{eq:bimax}, \eqref{eq:gmax}. 
\end{proof}

\subsection{Rank-$k$ Matrix Sensing}\label{app:sensingk}
\begin{proof}[Proof of Lemma \ref{lemma:QR-irrelevance}]
Since $\Uht$ and $\widetilde{U}^t$ have full rank and span the same subspace, there exists a $k\times k$, full rank matrix $R$ such that $\Uht = \widetilde{U}^t R = U^t R_U^t R$.
  We have:
  \begin{align*}
    \twonorm{\aff\left(\Uht V^{\dag}\right)-b} &= \twonorm{\aff\left(U^t \left(V\left(R_U^tR\right)^{\dag}\right)^{\dag}\right)-b}
    \geq \twonorm{\aff\left(U^t \left(\widetilde{V}^{t+1}\right)^{\dag}\right)-b}
  \end{align*}
with equality holding in the last step for $V=\widetilde{V}^{t+1}\left(\left(R_U^t R\right)^{\dag}\right)^{-1}$.
The proof of Theorem \ref{thm:sensing} shows that $\widetilde{V}^{t+1}$ is unique and has full rank (since $\dist\left(\widetilde{V}^{t+1},V^*\right)<1$).
This means that $\Vht$ is also unique and is equal to $\widetilde{V}^{t+1}\left(\left(R_U^t R\right)^{\dag}\right)^{-1}$.
This shows that $Span\left(\Vht\right) = Span\left(\widetilde{V}^{t+1}\right)$ and that both $\Vht$ and $\widetilde{V}^{t+1}$ have full rank.
\end{proof}

\begin{lemma}\label{lemma:sensing-minsingval-B}
Let linear measurement $\aff$ satisfy RIP for all $2k$-rank matrices and let $b=\aff(M)$ with $M\in \rmn$ being a rank-$k$ matrix.  Let $\delta_{2k}$ be the RIP constant for rank $2k$-matrices. Then, we have the following bound on the minimum singular value of $B$:
\begin{align}\label{eqn:sensing-minsingval-B}
\sigma_{\textrm{min}}(B) \geq 1-\delta_{2k}.
\end{align}
\end{lemma}
\begin{proof}
  Select any $w\in \R^{nk}$ such that $\twonorm{w}=1$. Let
\begin{align*}
w = \left[ \begin{array}{c}
            w_1 \\
	    w_2 \\
	    \vdots \\
	    w_k
           \end{array}
    \right]
\end{align*}
where each $w_p\in \rn$. Also denote $W \eqdef \left[w_1 w_2 \cdots w_k\right] \in \R^{n\times k}$, i.e., $w=vec(W)$. 

We have,
\begin{align*}
  w^\dag  B w &= \sum_{p,q=1}^k w_p^\dag  B_{pq} w_q
	= \sum_{p,q=1}^k w_p^\dag  \left(\sum_{i=1}^d A_i u_p^t (u_q^t)^\dag  A_i^\dag  \right)w_q 
	= \sum_{i=1}^d \sum_{p,q=1}^k w_p^\dag  A_i u_p^t (u_q^t)^\dag  A_i^\dag  w_q \\
	&= \sum_{i=1}^d \left(\sum_{p=1}^k w_p^\dag  A_i u_p^t\right) \left(\sum_{q=1}^k w_q^\dag  A_i u_q^t\right)=\sum_{i=1}^d \Tr\left(A_iU^tW^\dag \right)^2.
\end{align*}
Now, using RIP (see Definition~\ref{defn:rip}) along with the above equation, we get:
\begin{align*}
	w^\dag  B w&=\sum_{i=1}^d \Tr\left(A_iU^tW^\dag \right)^2\geq \left(1-\delta\right) \frob{U^tW^\dag }^2=\left(1-\delta_{2k}\right) \|W\|_F^2=(1-\delta_{2k})\|w\|^2=(1-\delta_{2k}). 
\end{align*}
Since $w$ was arbitrary, this proves the lemma.
\end{proof}

% \begin{lemma}\label{lemma:sensing-updateeqn}
%   Suppose we have $U^{(t)}$ and $\widehat{V}^{(t)}$ in time step $t$ of Algorithm \ref{algo:sensing-altmin}.
% Let $U^{(t)}=\left[u_1^{(t)},\cdots,u_k^{(t)}\right]$ and $\widehat{V}^{(t)}=\left[\widehat{v}_1^{(t)},\cdots,
% \widehat{v}_k^{(t)}\right]$ where each $u_p^{(t)}$ and $\widehat{v}_p^{(t)} \in \R^n$. Then $\widehat{v}_q^{(t)}$
% satisfy the following set of equations:
% \begin{align}\label{eqn:sensing-updateeqn}
%   \sum_{q=1}^k \left(\sum_{i=1}^s A_i u_p^{(t)} u_q^{(t)^\dag } A_i^\dag \right)\widehat{v}_q^{(t)}
% 	= \sum_{q=1}^k \left(\sum_{i=1}^s A_i u_p^{(t)} u_q^{*^\dag } A_i^\dag \right) v_q^* \;\;\;\;\; \forall \; p\in[k].
% \end{align}
% \end{lemma}

% \begin{proof}
%   Let $Err(V)\eqdef \sum_i \left(\Tr\left(A_i X\right) - \Tr\left(A_i U^{(t)}V^\dag \right)\right)^2$. Since $\widehat{V}^{(t)}$
% minimizes $E(V)$, we have $\Diff E(\widehat{V}^{(t)})=0$.
% \begin{align*}
%   &\Diff_{v_p} Err(\widehat{V}^{(t)})=0 \\
%   \Rightarrow &\sum_{i=1}^s \left(\sum_{l=1}^k {v_q^{(t)}}^\dag  A_i u_q^{(t)} - \sum_{l=1}^k {v_q^*}^\dag  A_i u_q^* \right)
% 	A_i u_p =0 \\
%   \Rightarrow &\sum_{l=1}^k \sum_{i=1}^s A_i u_p \left({v_q^{(t)}}^\dag  A_i u_q^{(t)}\right)
% 	= \sum_{l=1}^k \sum_{i=1}^s A_i u_p \left({v_q^*}^\dag  A_i u_q^*\right) \\
%   \Rightarrow &\sum_{l=1}^k \sum_{i=1}^s A_i u_p \left({u_q^{(t)}}^\dag  A_i^\dag  v_q^{(t)}\right)
% 	= \sum_{l=1}^k \sum_{i=1}^s A_i u_p \left({u_q^*}^\dag  A_i^\dag  v_q^*\right) \\
%   \Rightarrow &\sum_{l=1}^k \left(\sum_{i=1}^s A_i u_p {u_q^{(t)}}^\dag  A_i^\dag \right) v_q^{(t)}
% 	= \sum_{l=1}^k \left(\sum_{i=1}^s A_i u_p {u_q^*}^\dag  A_i^\dag \right) v_q^*
% \end{align*}
% This proves the lemma.
% \end{proof}
\begin{proof}[Proof of Lemma~\ref{lemma:sensing-twonorm-F}]
Note that,
\begin{align}
  \twonorm{F(\Sigma^*)^{-1}} \leq \frob{F(\Sigma^*)^{-1}} &= \twonorm{B^{-1} \left(BD-C\right)v^*} \nonumber\\
	&\leq \twonorm{B^{-1}}\twonorm{(BD-C)} \twonorm{v^*} \nonumber\\
	&\leq \frac{\sqrt{k}}{1-\delta_{2k}}\twonorm{(BD-C)} \label{eqn:sensing-twonorm-F-1}
\end{align}
where the last step follows from Lemma~\ref{lemma:sensing-minsingval-B}. %\eqref{eqn:sensing-minsingval-B} and \eqref{eqn:sensing-vp*-norm}.
Now we need to bound $\twonorm{(BD-C)}$. Choose any $w,z \in \R^{nk}$ such that $\twonorm{w}=\twonorm{z}=1$.
As in Lemma \ref{lemma:sensing-minsingval-B}, define the following components of $w$ and $z$:
\begin{align*}
w = \left[ \begin{array}{c}
            w_1 \\
	    w_2 \\
	    \vdots \\
	    w_k
           \end{array}
    \right] \mbox{ and }
z = \left[ \begin{array}{c}
            z_1 \\
	    z_2 \\
	    \vdots \\
	    z_k
           \end{array}
    \right]
\end{align*}
where each $w_p,z_p\in \R^n$ and $W \eqdef \left[w_1 w_2 \cdots w_k\right]$ and $Z \eqdef \left[z_1 z_2
\cdots z_k\right] \in \R^{n\times k}$. We have,
\begin{align*}
  w^\dag \left(BD-C\right)z = \sum_{p,q=1}^k w_p^\dag  \left(BD-C\right)_{pq} z_q
\end{align*}
We calculate $\left(BD-C\right)_{pq}$ as follows:
\begin{align*}
  \left(BD-C\right)_{pq} &= \sum_{l=1}^k B_{pl} D_{lq} - C_{pq} = \left(\sum_{l=1}^k B_{pl} \iprod{u_l^t}{u_q^*}\mathbb{I}_{n\times n}\right) - C_{pq} %= \left(\sum_{l=1}^k {u_q^*}^\dag u_l^t B_{pl}\right) - C_{pq} \\
	= \left(\sum_{l=1}^k {u_q^*}^\dag u_l^t \sum_{i=1}^d A_iu_p^t(u_l^t)^\dag A_i^\dag \right) - C_{pq} \\
	&= \sum_{i=1}^d A_iu_p^t {u_q^*}^\dag  \sum_{l=1}^k u_l^t(u_l^t)^\dag  A_i^\dag  -  \sum_{i=1}^d A_iu_p^t(u_q^*)^\dag A_i^\dag  
	= \sum_{i=1}^d A_iu_p^t {u_q^*}^\dag  \left(U^t(U^t)^\dag -\mathbb{I}_{n\times n}\right) A_i^\dag.
\end{align*}
So we have,
\begin{align*}
  w^\dag \left(BD-C\right)z &= \sum_{p,q=1}^k
    w_p^\dag  \sum_{i=1}^d A_iu_p^t {u_q^*}^\dag  \left(U^t(U^t)^\dag -\mathbb{I}_{n\times n}\right) A_i^\dag  z_q = \sum_{i=1}^d \sum_{p,q=1}^k w_p^\dag  A_iu_p^t {u_q^*}^\dag  \left(U^t(U^t)^\dag -\mathbb{I}_{n\times n}\right) A_i^\dag  z_q \\
%  &= \sum_{i=1}^d \left(\sum_{p=1}^k w_p^\dag  A_iu_p^t \right) \left(\sum_{q=1}^k z_q^\dag  A_i \left(U^t(U^t)^\dag -\mathbb{I}_{n\times n}\right)u_q^* \right) \\
  &= \sum_{i=1}^d \Tr\left(A_iU^tW^\dag \right) \Tr\left(A_i\left(U^t(U^t)^\dag -\mathbb{I}_{n\times n}\right) U^* Z^\dag \right) \\
  &\stackrel{(\zeta_1)}{\leq} \Tr\left({U^*}^\dag \left(U^t(U^t)^\dag -\mathbb{I}_{n\times n}\right)U^tW^\dag Z\right)
	+ \delta_{2k} \frob{U^tW^\dag } \frob{\left(U^t(U^t)^\dag -\mathbb{I}_{n\times n}\right)U^*Z^\dag } \\
  &\stackrel{(\zeta_2)}{\leq} \delta_{2k} \|W\|_F \sqrt{\frob{(U^*)^\dag \left(U^t(U^t)^\dag -\mathbb{I}_{n\times n}\right)^2U^*}\frob{Z^\dag Z}}\\ %\sqrt{\Tr\left(\left(UU^\dag -\mathbb{I}_{n\times n}\right)U^*Z^\dag 
&\stackrel{(\zeta_3)}{\leq} \delta_{2k}\sqrt{k}\cdot dist(U^t, U^*),
%  &\leq Z {U^*}^\dag \left(UU^\dag -\mathbb{I}_{n\times n}\right)\right)} \\ 
 % &= \delta \sqrt{\Tr\left(U^\dag UW^\dag W\right)} \sqrt{\Tr\left({U^*}^\dag \left(UU^\dag -\mathbb{I}_{n\times n}\right)^2
%		U^*Z^\dag Z\right)} \\
%  &= \delta \sqrt{\Tr\left(W^\dag W\right)} \sqrt{\frob{{U^*}^\dag \left(UU^\dag -\mathbb{I}_{n\times n}\right)^2U^*}
%	\frob{Z^\dag Z}} \\
%  &\stackrel{(\zeta_3)}{=} \delta \sqrt{\frob{{U^*}^\dag \left(UU^\dag -\mathbb{I}_{n\times n}\right)^2U^*}} \\
%  &\stackrel{(\zeta_4)}{=} \delta \sqrt{k \Gamma} \sqrt{1-\lambda^2}
\end{align*}
where $(\zeta_1)$ follows from the fact that $\calA$ satisfies $2k$-RIP and Lemma \ref{lem:alt_rip}, $(\zeta_2)$ follows from the fact that
$\left(U^t(U^t)^\dag -\mathbb{I}_{n\times n}\right)U^t=0$, $(\zeta_3)$ follows from the following: $\|W\|_F=\|w\|_2=1$, $\|Z^\dag Z\|_F\leq \|Z\|_F^2=1$ and 
% \begin{align*}
%   \Tr\left(W^\dag W\right) = \sum_{p=1}^k \twonorm{w_p}^2 = \twonorm{w}^2 = 1 \mbox{ and }
% \end{align*}
% \begin{align*}
%   \frob{Z^\dag Z}^2	 = \sum_{p,q=1}^k \left(\iprod{z_p}{z_q}\right)^2
% 	&\leq \sum_{p,q=1}^k \twonorm{z_p}^2\twonorm{z_q}^2 \\
% 	&= \left(\sum_{p=1}^k \twonorm{z_p}^2\right)^2 \\
% 	&= 1.
% \end{align*}
and finally : $\frob{\left(U^t(U^t)^\dag -\mathbb{I}_{n\times n}\right)U^*}\leq \sqrt{k}\twonorm{\left(U^t(U^t)^\dag -\mathbb{I}_{n\times n}\right)U^*}$. 
%\begin{align*}
%  \frob{{U^*}^\dag \left(UU^\dag -\mathbb{I}_{n\times n}\right)^2U^*}
%	&\leq \left(\sum_{p=1}^k \twonorm{\left(UU^\dag -\mathbb{I}_{n\times n}\right) u_p^*}^2\right) \\
%	&\leq k \Gamma \left(1-\lambda^2\right).
%\end{align*}

Since $w$ and $z$ were arbitrary unit vectors, we can conclude that $\twonorm{BD-C} \leq \delta_{2k} \sqrt{k}\cdot dist(U^t, U^*)$. 
%\sqrt{1-\lambda^2}$. 
Plugging this bound in \eqref{eqn:sensing-twonorm-F-1} proves the lemma.
\end{proof}

\begin{proof}[Proof of Lemma~\ref{lemma:sensing-minsingval-Rt}]
Note that $\|\Sigma^*(R^{(t+1)})^{-1}\|_2\leq \frac{\sigma^*_1}{\sigma_{\text{min}}(R^{(t+1)})}$. Now,
\begin{align}
  \sigma_{\text{min}}(R^{(t+1)})&=\min_{z, \|z\|_2=1}\|R^{(t+1)}z\|_2=\min_{z, \|z\|_2=1}\|V^{(t+1)}R^{(t+1)}z\|_2,\nonumber\\
&=\min_{z, \|z\|_2=1}\|V^*\Sigma^*{U^*}^\dag  U^{(t)}z-Fz\|_2, \nonumber\\
&\geq \min_{z, \|z\|_2=1}\|V^*\Sigma^*{U^*}^\dag  U^{(t)}z\|_2-\|Fz\|_2,\nonumber\\
&\geq \min_{z, \|z\|_2=1}\|V^*\Sigma^*{U^*}^\dag  U^{(t)}z\|_2-\|F\|_2,\nonumber\\
&\geq \sigma_{k}^*\sigma_{\text{min}}({U^*}^\dag  U^{(t)})-\|F\|_2,\nonumber\\
&\geq \sigma_{k}^*\sqrt{1-\twonorm{{U_{\perp}^*}^\dag U^{(t)}}^2}-\sigma^*_{1}\|F(\Sigma^*)^{-1}\|_2,\nonumber\\
&=\sigma_{k}^*\sqrt{1-\dist(U^*, U^{(t)})^2}-\sigma^*_{1}\|F(\Sigma^*)^{-1}\|_2. 
\end{align}
Lemma now follows using above inequality with Lemma~\ref{lemma:sensing-twonorm-F}. 

% Define $E=V^*\Sigma^*{U^*}^\dag  U^{(t)}(\Sigma^*)^{-1} - F(\Sigma^*)^{-1}=U_E\Sigma_EV_E^T$, then using \eqref{eqn:sensing-orthogonaliter-error}, we get: 
%   \begin{align}
%   E\Sigma^*(R^{(t+1)})^{-1}&=V^{(t+1)},\nonumber\\
% \Sigma^*(R^{(t+1)})^{-1}&=V_E\Sigma_E^{-1}U_E^TV^{t+1},\nonumber\\
% \|\Sigma^*(R^{(t+1)})^{-1}\|_2&=\|V_E\Sigma_E^{-1}U_E^TV^{t+1}\|_2\leq \|\Sigma_E^{-1}\|_2=\frac{1}{\sigma_{\text{min}}(E)}.\label{eq:Rs1}
%   \end{align}
% Now, we bound $\sigma_{\text{min}}(E)$:
% \begin{align}
%   \sigma_{\text{min}}(E)&=\min_{z, \|z\|_2=1}\|V^*\Sigma^*{U^*}^\dag  U^{(t)}(\Sigma^*)^{-1}z - F(\Sigma^*)^{-1}z\|_2,\nonumber\\
% &\geq\min_{z, \|z\|_2=1}\|V^*\Sigma^*{U^*}^\dag  U^{(t)}(\Sigma^*)^{-1}z\|_2-\|F(\Sigma^*)^{-1}z\|_2,\nonumber\\
% &\geq \min_{z, \|z\|_2=1}\|V^*\Sigma^*{U^*}^\dag  U^{(t)}(\Sigma^*)^{-1}z\|_2-\max_{z, \|z\|_2=1}\|F(\Sigma^*)^{-1}z\|_2,\nonumber\\
% &\geq \min_{z, \|z\|_2=1}\|\Sigma^*{U^*}^\dag  U^{(t)}(\Sigma^*)^{-1}z\|_2-\|F(\Sigma^*)^{-1}\|_2.\label{eq:Esigma}
% \end{align}
% Also, $\min_{z, \|z\|_2=1}\|\Sigma^*{U^*}^\dag  U^{(t)}(\Sigma^*)^{-1}z\|_2=\sigma_{\text{min}}(\Sigma^*{U^*}^\dag  U^{(t)}(\Sigma^*)^{-1})=\sigma_{\text{min}}({U^*}^\dag  U^{(t)})$ where last equality follows using the fact that $A^{-1}BA$ has same eigenvalues as $A$. Therefore, using \eqref{eq:Esigma} and Theorem 2.6.1 in \cite{GolVan96} ($U_{\perp}^*$ is an orthonormal basis of $\textrm{Span}\left(u_1^*,u_2^*,\cdots,u_k^*\right)^{\perp}$), we get:
% \begin{align}\label{eq:Esigma1}
%   \sigma_{\text{min}}(E)&\geq \sigma_{\text{min}}({U^*}^\dag  U^{(t)})-\|F(\Sigma^*)^{-1}\|_2\geq \sqrt{1-\twonorm{{U_{\perp}^*}^\dag U}^2}-\|F(\Sigma^*)^{-1}\|_2,\nonumber\\
% &\geq \sqrt{1-\dist^2\left(U^{(t)},U^*\right)}-\|F(\Sigma^*)^{-1}\|_2.
% \end{align}
% Lemma now follows using \eqref{eq:Rs1}, \eqref{eq:Esigma1} and Lemma~\ref{lemma:sensing-twonorm-F}. 
%  We have the following set of inequalities:
% \begin{align*}
%   \sigma_{\textrm{min}}\left(R^{(t)}\right) = \min_{\twonorm{z}=1} \twonorm{R^{(t)}z}
% 	&\stackrel{(\zeta_1)}{=} \min_{\twonorm{z}=1} \twonorm{V^{(t)}R^{(t)}z} \\
% 	&= \min_{\twonorm{z}=1} \twonorm{V^*{U^*}^\dag Uz - Fz} \\
% 	&\geq \min_{\twonorm{z}=1} \twonorm{V^*{U^*}^\dag Uz} - \twonorm{F} \\
% 	&\geq \sigma_{\textrm{min}}\left(V^*\right)\min_{\twonorm{z}=1} \twonorm{{U^*}^\dag Uz} - \twonorm{F} \\
% 	&\stackrel{(\zeta_2)}{\geq} \sqrt{\gamma}\sigma_{\textrm{min}}\left({U^*}^\dag U\right) - \twonorm{F} \\
% 	&\stackrel{(\zeta_3)}{\geq} \gamma\sqrt{1-\twonorm{{U_{\perp}^*}^\dag U}^2} - \twonorm{F} \\
% 	&= \gamma\lambda - \twonorm{F}
% \end{align*}
% where $(\zeta_1)$ follows from the fact that $V^{(t)}$ is an orthonormal column matrix, $(\zeta_2)$ follows from the fact that
% $V^*$ is an orthogonal column matrix and \eqref{eqn:sensing-vp*-norm} and finally $(\zeta_3)$ follows from
% \eqref{eqn:sensing-up*-norm} and the proof of Theorem 2.6.1 in \cite{GolVan96} ($U_{\perp}^*$ is an orthonormal
% basis of $\textrm{Span}\left(u_1^*,u_2^*,\cdots,u_k^*\right)^{\perp}$). This proves the lemma.
\end{proof}
\subsection{Noisy Matrix Sensing: Proofs}\label{app:noisy_sensing}
We now consider an extension of the matrix sensing problem where measurements can be corrupted arbitrarily using a bounded noise. That is, we 
%We now extend Theorem \ref{thm:sensing} in the presence of noise. In particular, we consider the case where we 
observe $b=\aff\left(M+N\right)$, where $N$ is the noise matrix. % s.t.,
%\begin{align}\label{eqn:matrix-noisy}
%  M = U^* \Sigma^* {V^*}^{\dag},\ \ \ 
%  N = \Un \Sn (\Vn)^{\dag},
%\end{align}
%nd $N=\Un \Sn (\Vn)^{\dag}$ is the SVD of the noise matrix $N$. 
For this noisy case as well, we show that \as recovers $M$ upto an additive approximation depending on the Frobenius norm of  $N$. 
\begin{theorem}\label{thm:sensing-noisycase}
 Let $M$ and $\calA(\cdot)$ be as defined in Theorem~\ref{thm:sensing_err}. Suppose, \as algorithm (Algorithm \ref{algo:sensing-altmin}) is supplied inputs $\aff$, $b=\aff(M+N)$, where $N$ is the noise matrix s.t. $\left\|N\right\|_F  < \frac{1}{100}\so_k$.  Then, after $T=4\log(2/\epsilon)$ steps, iterates $\widehat{U}^{T}$, $\widehat{V}^{T}$ of \as  satisfy: 
%   Let $M=U^*\Sigma^*V^{*^\dag }$ be a rank-$k$ matrix with non zero singular values $\sigma_1^*\geq \sigma_2^*\dots\geq \sigma_k^*$ and $N=\Un \Sn (\Vn)^{\dag}$ be a matrix with singular values $\sn_1 \geq \sn_2 \geq
%\cdots \sn_n$ such that $\left\|N\right\|_F  < \frac{1}{100}\sigma_k$. Also, let $\calA(\cdot):\rmn \rightarrow \mathbb{R}^d$ be a linear measurement
%operator that satisfies $2k$-RIP with RIP constant $\delta_{2k}<\frac{(\sigma^*_k)^2}{100k(\sigma_1^*)^2}$.  Suppose, AltMin algorithm (Algorithm \ref{algo:sensing-altmin}) is supplied inputs $\aff$, $b=\aff(M+N)$. Then, after $T=4\log(2/\epsilon)$ steps, iterates $U^{T}$, $V^{T}$ of the AltMin algorithm  satisfy: 
%with $\dist\left(U^{(0)},U^*\right) \leq \frac{1}{2}$. Then $U^{(t)}$ and $V^{(t)}$ satisfy:
\begin{align*}
  \dist\left(\Vw^{T},V^*\right) \leq \frac{10\|N\|_F}{\so_k} + \epsilon, \ 
  \dist\left(\Uw^{T},U^*\right) \leq \frac{10\|N\|_F}{\so_k} + \epsilon.
\end{align*}
See Definition~\ref{defn:dist} for definition of $\dist\left(U, W\right)$. 
%where $\alpha_V, \alpha_U < 1$ are some constants depending only on $\delta_{2k}$. 
% denotes the principal angle based distance between subspaces spanned by the columns of $U\in \mathbb{R}^{m\times k}$ and $W\in \mathbb{R}^{m\times k}$ (See Definition~\ref{defn:dist}).
\end{theorem}
%\begin{theorem}\label{thm:sensing-noisy-main}
%   Let $M$ and $\calA(\cdot)$ be as defined in Theorem~\ref{thm:sensing_err}. 
%=U^*\Sigma^*V^{*^\dag }$ be a rank-$k$ matrix with non zero singular values $\sigma_1^*\geq \sigma_2^*\dots\geq \sigma_k^*$ and $N=\Un \Sn (\Vn)^{\dag}$ be the noise matrix. Also, let $\calA(\cdot):\rmn \rightarrow \mathbb{R}^d$ be a linear measurement operator that satisfies $2k$-RIP with RIP constant $\delta_{2k}<\frac{(\sigma^*_k)^2}{100k(\sigma_1^*)^2}$.  
%Suppose, \as algorithm (Algorithm \ref{algo:sensing-altmin}) is supplied inputs $\aff$, $b=\aff(M+N)$. Then, after $T=4\log(2/\epsilon)$ steps, iterates $\widehat{U}^{T}$, $\widehat{V}^{T}$ of \as  satisfy: 
%$$\|M-\widehat{U}\widehat{V}^T\|_F\leq C \|N\|_F,$$
%where $C=(\frac{\so_1}{\so_k}+10\delta_{2k})$. 
%\end{theorem}
\begin{proof}
At a high level, our proof for noisy case follows closely, the exact case proof given in Section~\ref{sec:sensing}. That is, we show that the update of \as algorithm is similar to power-method type update but with two errors terms: one due to incomplete measurements and another due to the noise matrix.

Similar to our proof for sensing problem (Section~\ref{sec:sensing}), we analyze QR-decomposition based updates. That is,
\begin{align*}
\widehat{U}^{t}&= U^{t}R_U^{t}\ \ \text{ (QR decomposition), }\\
\widehat{V}^{t+1}&=\argmin_{V} \  \|\aff(U^t V^\dag )-b\|_2^2,\\ 
\widehat{V}^{t+1}&=V^{t+1}R_V^{t+1}.\ \ \text{ (QR decomposition) }
\end{align*}
Similar to Lemma~\ref{lemma:sensing-updateeqn1}, we can re-write the above given update equation as:
\begin{align}
  \Vht&=\Vo\So(\Uo)^\dag\Ut-F+\Vn \Sn (\Un)^{\dag}U^{(t)}-G,\nonumber\\
  \Vt&=\Vht(\Rt)^{-1},\label{eq:noisy_update_k}
\end{align}
where, $F$ is the error matrix and is as defined in \eqref{eq:F_k} and $G$ is the error matrix due to noise and is given by:
\begin{equation}\label{eq:G_k}
G \eqdef \left[ 
\left(B^{-1}\left(BD^N-C^N\right)S^Nv^N\right)_1\ \ \cdots\ 
 \left(B^{-1}\left(BD^N-C^N\right)S^Nv^N\right)_k
 \right],
\end{equation}
where $B$, $C$ and $D$ defined in the previous section (See \eqref{eq:BCD_k}) and $C^N$ and $D^N$ are defined below:
\begin{align*}
 C^N \eqdef \left[\begin{array}{ccc}
                 C_{11}^N & \cdots & C_{1m}^N \\
		 \vdots &  \ddots &\vdots \\
                 C_{k1}^N & \cdots & C_{km}^N \\
                \end{array} \right]
\mbox{ , }
 D^N \eqdef \left[\begin{array}{ccc}
                 D_{11}^N & \cdots & D_{1m}^N \\
		 \vdots &  \ddots &\vdots \\
                 D_{k1}^N &  \cdots & D_{km}^N \\
                \end{array} \right],
\end{align*}
with $C_{pq}^N \eqdef \sum_{i=1}^d A_i u_p^{(t)}{(u^N_q)}^\dag  A_i^\dag $ and $D_{pq}^N \eqdef \iprod{u_p^{(t)}}{u^N_q}\mathbb{I}_{n\times n}$.
Also, $$S^N=\left[\begin{matrix}\sn_1I_{n}&\dots&0_{n}\\\vdots&\vdots&\vdots\\0_n&\dots&\sn_NI_n\end{matrix}\right],\ \  v^N=\left[\begin{matrix}v^N_1\\\vdots\\v_k^N\end{matrix}\right].$$
Now, multiplying \eqref{eq:noisy_update_k} with $\Vo_\perp$, we get:
\begin{equation*}
{V_{\perp}^*}^\dag \Vt =  ({V_{\perp}^*}^\dag \Vn \Sn U^{N{\dag}}U^{(t)}- {V_{\perp}^*}^\dag  F - {V_{\perp}^*}^\dag  G) {\Rt}^{-1}.%\label{eqn:sensing-orthogonaliter-error-noisycase}
\end{equation*}
That is,%Using \eqref{eqn:sensing-orthogonaliter-error-noisycase},
\begin{align}
&\hspace*{-15pt}\dist(\Vo, \Vt)=\|{V_{\perp}^*}^\dag \Vt\|_2\nonumber\\&\hspace*{-15pt}\leq (\| \Vn \Sn (\Un)^{\dag}U^{(t)}\|_2+ \|F |_2 + \| G\|_2) \|{(\Rt)}^{-1}\|_2,\nonumber\\
&\hspace*{-15pt}\leq \left(\sn_1+ \|F(\So)^{-1}\|_2\|\So\|_2 + \| G\|_2\right) \|({\Rt})^{-1}\|_2, \nonumber\\
&\hspace*{-15pt}\leq \left(\sn_1+\frac{\sigma_1^*\delta_{2k}k}{1-\delta_{2k}}\dist(\Ut, \Uo) + \| G\|_2\right) \|({\Rt})^{-1}\|_2,\label{eq:dist_noisy_1}
\end{align}
where the last inequality follows using Lemma~\ref{lemma:sensing-twonorm-F}. 

Now, we break down the proof in the following two steps:
\begin{itemize}
%\item Bound $\twonorm{Y \Lambda X^{\dag}U^{(t)}}$ (Lemma \ref{lemma:sensing-twonorm-noise-noisycase})
\item Bound $\twonorm{G}$ (Lemma \ref{lemma:sensing-twonorm-G-noisycase}, analogous to Lemma \ref{lemma:sensing-twonorm-F})
\item Bound $\|({\Rt})^{-1}\|_2$ (Lemma \ref{lemma:sensing-minsingval-Rt-noisycase}, similar to Lemma \ref{lemma:sensing-minsingval-Rt})
\end{itemize}
Later in this section, we provide the above mentioned lemmas and their detailed proof. 

Now, by assumption, $\sn_1\leq \frob{N}\leq \so_k$. Also, as $\delta_{2k}\leq 1/2$, $\frac{1}{1-\delta_{2k}}\leq 2$. Finally, assume $\dist(\Vo, \Vt)\geq \max(10\cdot\frac{\sn_1}{\so_k}, 10\frac{\|N\|_F}{\sigma_1^*})$. Using these observations and lemmas~ \ref{lemma:sensing-twonorm-G-noisycase}, \ref{lemma:sensing-minsingval-Rt-noisycase} along with \eqref{eq:dist_noisy_1}, we get:
%Now, using 
%\begin{align}
%  \dist(\Vo, \Vt)\leq \frac{\left(\sn_1+\frac{\sigma^*_1\delta_{2k}k}{1-\delta_{2k}}\dist(\Ut, \Uo)+\frac{\delta_{2k} \frob{N}}{1-\delta_{2k}}\right)}{\left(\sigma_k^*\sqrt{1-\dist(\Ut,\Uo)^2}-\sn_1-\frac{\sigma^*_1\delta_{2k}k}{1-\delta_{2k}}\dist(\Ut, \Uo)-\frac{\delta_{2k} \frob{N}}{1-\delta_{2k}}\right)}
%\end{align}
%Now, by assumption, $\sn_1\leq \frob{N}\leq \so_k$. Also, as $\delta_{2k}\leq 1/2$, $\frac{1}{1-\delta_{2k}}\leq 2$. Furthermore, assume $\dist(\Vo, \Vt)\geq \max(10\cdot\frac{\sn_1}{\so_k}, 10\frac{\|N\|_F}{\sigma_1^*})$. 
%
%Hence, using above equation, 
\begin{align}
  \dist(\Vo, \Vt)\leq \frac{0.5\dist(\Uo, \Ut)}{\sqrt{1-\dist(\Ut,\Uo)^2}-0.5\dist(\Uo, \Ut)}. \label{eq:dist_noisy_f}
\end{align}
As, $U^{0}$ is obtained using SVD of $\sum_i A_i b_i$. Hence, using Lemma~\ref{lem:jmd}, we have: %2.1 of \cite{JainMD10}, we have:
\begin{align*}
 &\|\aff(U^{0}\Sigma^{0}V^{0}-U^*\Sigma^*(V^*)^\dag)\|_2^2\leq 0.5 \|\aff(N)\|_{2}^2+4\delta_{2k}\|\aff(U^*\Sigma^*(V^*)^\dag)\|_2^2,\\
&\Rightarrow \|U^{0}\Sigma^{0}V^{0}-U^*\Sigma^*(V^*)^\dag\|_F^2\leq \|N\|_{F}^2+4\delta_{2k}(1+\delta_{2k})\|\Sigma^*\|_F^2,\\
%\Rightarrow&\|U^{(0)}(U^{(0)})^\dag U^*\Sigma^*(V^*)^\dag-U^*\Sigma^*(V^*)^\dag\|_F^2\leq 4\delta_{2k}(1+3\delta_{2k})\|\Sigma^*\|_F^2,\nonumber\\
&\Rightarrow(\sigma_k^*)^2\|(U^{0}(U^{0})^\dag-I)U^*\|_F^2\leq \|N\|_{F}^2+4\delta_{2k}(1+\delta_{2k})k(\sigma_1^*)^2,\\
&\Rightarrow\dist(U^{0}, U^*)\leq \|(U^{0}(U^{0})^\dag-I)U^*\|_F^2\leq \frac{\|N\|_{F}^2}{(\so_k)^2}+6\delta_{2k}k\left(\frac{\sigma_1^*}{\sigma_k^*}\right)^2<\frac{1}{2},
\end{align*}
where last inequality follows using $\frac{\|N\|_{F}}{\so_k}<1/100$. 

Theorem now follows using above equation with \eqref{eq:dist_noisy_f}. 
%similar to the proof of Theorem~\ref{thm:sensing-intro} for exact case, our proof reduces 
%\begin{align}
%  \dist(\Vo, \Vt)&=\|{V_{\perp}^*}^\dag \Vt\|_2\leq \| Y \Lambda X^{\dag}U^{(t)}{(\Rt)}^{-1}\|_2+ \|F ({\Rt})^{-1}\|_2 + \| G ({\Rt})^{-1}\|_2,\nonumber\\
%&\leq \left(\lambda_1+ \|F(\So)^{-1}\|_2\|\So\|_2 + \| G(\Lambda)^{-1}\|_2\|\Lambda\|_2\right) \|({\Rt})^{-1}\|_2, \nonumber
%\end{align}
% Proceeding in a similar way as in Section \ref{sec:sensing_k}, we obtain the following equation which is an analogue of \eqref{eqn:sensing-orthogonaliter-error}.
% \begin{align}
%   {V_{\perp}^*}^\dag V^{(t)} = Y \Lambda X^{\dag}U^{(t)}{R^{(t)}}^{-1} - {V_{\perp}^*}^\dag  F {R^{(t)}}^{-1} - {V_{\perp}^*}^\dag  G {R^{(t)}}^{-1}.\label{eqn:sensing-orthogonaliter-error-noisycase}
% \end{align}
% where $F \eqdef \left[ \begin{array}{cccc}
% \left(B^{-1}\left(BD-C\right)v^*\right)_1 & \left(B^{-1}\left(BD-C\right)v^*\right)_2
% & \cdots & \left(B^{-1}\left(BD-C\right)v^*\right)_k
% \end{array} \right]$ (as before) and 
%With \eqref{eqn:sensing-orthogonaliter-error-noisycase} in hand, we set about proving the following along the lines of lemmas in Section \ref{sec:sensing_k}:
%\begin{itemize}
%  \item $\twonorm{Y \Lambda X^{\dag}U^{(t)}}$ is very small (Lemma \ref{lemma:sensing-twonorm-noise-noisycase}),
%  \item $\twonorm{F}$ is very small (Lemma \ref{lemma:sensing-twonorm-F}),
%  \item $\twonorm{G}$ is very small (Lemma \ref{lemma:sensing-twonorm-G-noisycase}, analogous to Lemma \ref{lemma:sensing-twonorm-F}) and
%  \item the minimum singular value of $R^{(t)}$ is not very small (Lemma \ref{lemma:sensing-minsingval-Rt-noisycase}, analogous to Lemma \ref{lemma:sensing-minsingval-Rt}).
%\end{itemize}
\end{proof}
% \begin{lemma}\label{lemma:sensing-twonorm-noise-noisycase}
%   We have the following bound on $\twonorm{Y \Lambda X^{\dag}U^{(t)}}$:
% \begin{align}\label{eqn:sensing-twonorm-noise-noisycase}
%   \twonorm{Y \Lambda X^{\dag}U^{(t)}} \leq \lambda_1.
% \end{align}
% \end{lemma}
% \begin{proof}
%   The proof is straightforward:
% \begin{align*}
%   \twonorm{Y \Lambda X^{\dag}U^{(t)}} \leq \twonorm{Y \Lambda X^{\dag}} \twonorm{U^{(t)}} = \lambda_1 \cdot \twonorm{U^{(t)}} = \lambda_1.
% \end{align*}
% where the last step follows from the fact that $U^{(t)}$ is an orthonormal matrix.
% \end{proof}
\begin{lemma}\label{lemma:sensing-twonorm-G-noisycase}
Let linear measurement $\aff$ satisfy RIP for all  $2k$-rank matrices and let $b=\aff(M+N)$ with $M\in \rmn$ being a rank-$k$ matrix and let $N=\Un\Sn (\Vn)^\dag$.  Let $\delta_{2k}$ be the RIP constant for rank $2k$-matrices. Then, we have the following bound on $\twonorm{G}$:
\begin{align}\label{eqn:sensing-twonorm-G-noisycase}
    \twonorm{G} \leq \frac{\delta_{2k} \frob{N}}{1-\delta_{2k}}.
\end{align}
\end{lemma}
\begin{proof}
Note that, 
\begin{align}
  \|G\|_2\leq \|G\|_F&=\|B^{-1}(BD^N-C^N)S^Nv^N\|_2\nonumber\\&\leq \|B^{-1}\|_2\|(BD^N-C^N)S^N\|_2\|S^Nv^N\|_2\leq \frac{\sqrt{k}}{1-\delta_{2k}} \|(BD^N-C^N)S^N\|_2,
\end{align}
where the last inequality follows using Lemma~\ref{lemma:sensing-minsingval-B} and the fact that $\|\Vn\|_F=\sqrt{k}$. 
%  Proceeding along the lines of the proof of Lemma \ref{lemma:sensing-twonorm-F}, we obtain the following equation (which is the analogue of \eqref{eqn:eqn:sensing-twonorm-F-1}):
%\begin{align*}
%  \twonorm{G} \leq \frac{1}{1-\delta}\twonorm{(BD^N-C^N)}.
%\end{align*}
Now let $w=[w_1^\dag\ w_2^\dag\ \dots\ w_k^\dag]^\dag\in \mathbb{R}^{nk}$ and $z=[z_1^\dag\ z_2^\dag\ \dots\ z_n^\dag]^\dag\in \mathbb{R}^{n^2}$ be any two arbitrary vectors such that $\twonorm{w}=\twonorm{z}=1$. Then, 
%Now for all $w\in \mathbb{R}^{nk}$ and $z\in \mathbb{R}^{nk}$ such that $\twonorm{w}=\twonorm{z}=1$, % (as in Lemma \ref{lemma:sensing-twonorm-F}; please refer there for a definition of all notation),
%we have the following:
\begin{align*}
  w^\dag \left(BD^N-C^N\right)S^Nz &= \sum_{p=1}^k \sum_{q=1}^n
    w_p^\dag  \sum_{i=1}^d A_i\ut_p {\un_q}^\dag  \left(\Ut(\Ut)^\dag -\mathbb{I}_{n\times n}\right) A_i^\dag  \sn_qz_q \\
  &= \sum_{i=1}^d \sum_{p=1}^k \sum_{q=1}^n w_p^\dag  A_i\ut_p {\un_q}^\dag  \left(\Ut(\Ut)^\dag -\mathbb{I}_{n\times n}\right) A_i^\dag  \sn_qz_q \\
  &= \sum_{i=1}^d \left(\sum_{p=1}^k w_p^\dag  A_i\ut_p \right) \left(\sum_{q=1}^n \sn_qz_q^\dag  A_i \left(\Ut(\Ut)^\dag -\mathbb{I}_{n\times n}\right)\un_q \right) \\
  &= \sum_{q=1}^n \left(\sum_{i=1}^d \Tr\left(A_iUW^\dag \right) \Tr\left(A_i\left(\Ut(\Ut)^\dag -\mathbb{I}_{n\times n}\right) \un_q \sn_qz_q^\dag \right)\right).
\end{align*}
Now, using RIP, we get:
\begin{align*}
  w^\dag \left(BD^N-C^N\right)z&\leq \sum_{q=1}^n {\un_q}^\dag \left(\Ut{\Ut}^\dag -\mathbb{I}_{n\times n}\right)\Ut W^\dag \sn_qz_q
	+ \delta_{2k} \frob{\Ut W^\dag } \frob{\left(\Ut{\Ut}^\dag -\mathbb{I}_{n\times n}\right)\un_q\sn_q{z_q}^\dag } \\
  &\leq \sum_{q=1}^n \delta_{2k} \|W^\dag\|_F \|(\Ut(\Ut)^\dag -\mathbb{I}_{n\times n})\un_q\|_2\|\sn_q{z_q}\|_2,\\%\sqrt{\Tr\left(\left(\Ut(\Ut)^\dag -\mathbb{I}_{n\times n}\right)x_q z_q^\dag 		z_q {x_q}^\dag \left(\Ut(\Ut)^\dag -\mathbb{I}_{n\times n}\right)\right)} \\ 
%  &= \delta \sqrt{\Tr\left(U^\dag UW^\dag W\right)} \sum_{q=1}^n \sqrt{\Tr\left({x_q}^\dag \left(\Ut(\Ut)^\dag -\mathbb{I}_{n\times n}\right)^2
%		x_qz_q^\dag z_q\right)} \\
  &\leq \delta_{2k} \ \sum_{q=1}^n \twonorm{\un_q} \twonorm{\sn_qz_q}= \delta_{2k} \sum_{q=1}^n \sn_q\twonorm{z_q},\\&\leq \delta_{2k} \sqrt{\sum_{q=1}^n \twonorm{z_q}^2} \sqrt{\sum_{q=1}^n (\sn_q)^2}\leq \delta_{2k} \frob{N}.
\end{align*}
This finishes the proof.
\end{proof}
\begin{lemma}\label{lemma:sensing-minsingval-Rt-noisycase}
Assuming conditions of Lemma~\ref{lemma:sensing-twonorm-G-noisycase}, we have the following bound on the minimum singular value of $R^{(t)}$:
\begin{align*}%\label{eqn:sensing-minsingval-Rt-noisycase}
  \sigma_{\textrm{min}}\left(\Rt\right) \geq \so_k\sqrt{1-\dist(\Ut, \Uo)^2} - \sn_1 - \twonorm{F} - \twonorm{G}.
\end{align*}
\end{lemma}
\begin{proof}
  Similar to the proof of Lemma \ref{lemma:sensing-minsingval-Rt}, we have the following set of inequalities:
\begin{align*}
  \sigma_{\textrm{min}}\left(\Rt\right) = \min_{\twonorm{z}=1} \twonorm{\Rt z}
	&= \min_{\twonorm{z}=1} \twonorm{V^{(t)}\Rt z} \\
	&= \min_{\twonorm{z}=1} \twonorm{V^*\So{U^*}^\dag \Ut z + \Vn\Sn(\Un)^{\dag}U^{(t)} z - Fz - Gz} \\
	&\geq \min_{\twonorm{z}=1} \twonorm{V^*\So{U^*}^\dag \Ut z} - \twonorm{\Vn\Sn(\Un)^{\dag}} - \twonorm{F} - \twonorm{G} \\
	&\geq \so_k\min_{\twonorm{z}=1} \twonorm{{U^*}^\dag \Ut z} - \sn_1 - \twonorm{F} - \twonorm{G} \\
%	&\geq \sqrt{\gamma}\sigma_{\textrm{min}}\left({U^*}^\dag U\right) - \sn_1 - \twonorm{F} - \twonorm{G} \\
	&\geq \so_k\sqrt{1-\twonorm{{U_{\perp}^*}^\dag U}^2} - \sn_1 - \twonorm{F} - \twonorm{G} \\
	&= \so_k\sqrt{1-\dist(\Ut,\Uo)^2} - \sn_1 - \twonorm{F} - \twonorm{G}.
\end{align*}
This proves the lemma.
\end{proof}
\subsection{Stagewise Alternating Minimization for Matrix Sensing: Proofs}\label{app:stagewise_sensing}
\begin{proof}[Proof of Lemma~\ref{lemma:sam_svp_err}]
As the initial point of the $i$-th stage is obtained by one step of SVP \cite{JainMD10}, using Lemma~\ref{lem:jmd}, we obtain:
$$\left\|M - \widehat{U}^0_{1:i}(\widehat{V}_{1:i}^0)^{\dag}\right\|_F^2 \leq \sum_{j=i+1}^k(\so_j)^2+2\delta_{2k}\|M-\Uw^T_{1:i-1}V^T_{1:i-1}\|_F^2.$$
Now, by assumption over the $(i-1)$-th stage error (this assumption follows from the inductive hypothesis in proof of Theorem~\ref{thm:sam_err}), 
$$\left\|M - \widehat{U}^0_{1:i}(\widehat{V}_{1:i}^0)^{\dag}\right\|_F^2 \leq \sum_{j=i+1}^k(\so_j)^2+2\delta_{2k} 16k(\so_i)^2.$$
Lemma now follows by setting $\delta_{2k}\leq \frac{1}{3200k}$. 
%We prove the lemma using mathematical induction. \\
%{\bf Base Case}: In the first stage, $\Uw^0_{1:1}$, $\Vw^0_{1:1}$ are given by the svd of $\aff^T(\aff(M))$. Now, using Lemma 2.1 of \cite{JainMD10}, we get:
%$$\|M-\Uw^0_{1:1}(\Vw^0_{1:1})^\dag\|_F^2\leq \sum_{j=2}^k\soj+\delta_{2k}\left(\sum_{j=1}^k\soj\right)\leq \sum_{j=2}^k\soj+\frac{1}{100}\soi.$$
%{\bf Induction Step}: We prove this step for two mutually exhaustive cases: a) $\frac{\so_i}{\so_{i+1}} < 5\sqrt{k}$, b) $\frac{\so_i}{\so_{i+1}} \geq 5\sqrt{k}$. Case (a) is proved in Lemma~\ref{lemma:cond-num-good} and Case (b) is proved in Lemma~\ref{lemma:cond-num-bad-altmin-performance}. 
\end{proof}
\begin{proof}[Proof of Lemma~\ref{lemma:sam_altmin}]
For our proof, we consider two cases: a) $\frac{\so_i}{\so_{i+1}} < 5\sqrt{k}$, b) $\frac{\so_i}{\so_{i+1}} \geq 5\sqrt{k}$.\\
{\bf Case (a)}: In this case, using monotonicity of the AltMin algorithm directly gives error bound. That is, $\|M - \widehat{U}^T_{1:i}(\widehat{V}_{1:i}^T)^{\dag}\|_F^2\leq \|M-\Uw^0_{1:i}V^0_{1:i}\|_F^2 \leq k(\so_{i+1})^2+\frac{25k}{100}(\so_{i+1})^2.$\\%See Lemma~\ref{lemma:cond-num-good} for a detailed proof.\\
{\bf Case (b)}: At a high level,  if $\frac{\so_i}{\so_{i+1}} \geq 5\sqrt{k}$ then $U^0_{1:i}$ is ``close'' to $U^*_{1:i}$ and hence the error bound follows by using an analysis similar to the noisy case. Note that $\so_{i+1}$ being small implies that the ``noise'' is small. See Lemma~\ref{lemma:cond-num-bad-altmin-performance} for a formal proof of this case. \hfill 
\end{proof}

\begin{lemma}\label{lemma:cond-num-bad-altmin-performance}
Assume conditions given in Theorem~\ref{thm:sam_err} are satisfied and  let $\frac{\so_i}{\so_{i+1}} \geq 5\sqrt{k}$. Also, let $$\left\|M - \widehat{U}^0_{1:i}(\widehat{V}^0_{1:i})^{\dag}\right\|_F^2 \leq \sum_{j=i+1}^k (\so_j)^2 + \frac{1}{100} (\so_i)^2.$$ Then, $U^T_{1:i}$, $V^T_{1:i}$ satisfy:$$\|M-\Uw^T_{1:i}V^T_{1:i}\|_F^2\leq \max(\epsilon, 16k(\so_{i+1})^2),$$
\end{lemma}
\begin{proof}
%We prove the lemma using mathematical induction. That is, let the lemma be true for the $i$-th stage. Then, we show that the Lemma holds for $i+1$-th stage as well. 
We first show that if $\sigma_i$ and $\sigma_{i+1}$ have large gap then $\forall\; t$, the $t^{\textrm{th}}$ iterate of the $i$-th stage, $\Uw^t_{1:i}$ is close to $\Uo_{1:i}$. Let $\Up^t$ be a basis of the subspace orthogonal to $\Uw^t_{1:i}$. 
\begin{align}
  \|\left(\Up^t\right)^\dag(M-\Uw^t_{1:i}(\Vw^t_{1:i})^\dag)\|_2&=\|\left(\Up^t\right)^\dag M\|_2\geq \|\left(\Up^t\right)^\dag\Uo_{1:i}\So_{1:i}(\Vo_{1:i})^\dag\|_2-\|\left(\Up^t\right)^\dag\Uo_{i+1:k}\So_{i+1:k}(\Vo_{i+1:k})^\dag\|_2,\nonumber\\
&\geq \so_i\|\left(\Up^t\right)^\dag\Uo_{1:i}\|_2-\so_{i+1}\geq \so_i(\|\left(\Up^t\right)^\dag\Uo_{1:i}\|_2-\frac{1}{5\sqrt{k}}). \label{eq:sam_bc_uperp1}
\end{align}
We also have:
\begin{align}
\|\left(\Up^t\right)^\dag(M-\Uw^t_{1:i}(\Vw^t_{1:i})^\dag)\|_2^2 &
	\leq \|M-\Uw^t_{1:i}(\Vw^t_{1:i})^\dag\|_F^2\leq \frac{1}{1-\delta_{2k}}\twonorm{\aff\left(M-\Uw^t_{1:i}(\Vw^t_{1:i})^\dag\right)}^2 \nonumber\\
	&\stackrel{(\zeta_1)}{\leq} \frac{1}{1-\delta_{2k}}\twonorm{\aff\left(M-\Uw^0_{1:i}(\Vw^0_{1:i})^\dag\right)}^2 \leq \frac{1+\delta_{2k}}{1-\delta_{2k}}\|M-\Uw^0_{1:i}(\Vw^0_{1:i})^\dag\|_F^2 \nonumber\\
	&\leq \frac{1+\delta_{2k}}{1-\delta_{2k}} \left(\sum_{j=i+1}^k \soj + \frac{1}{100} \soi\right) \nonumber\\
	&\leq \frac{1+\delta_{2k}}{1-\delta_{2k}} \left(k\soii+\frac{1}{100}\soi\right), \label{eq:sam_bc_uperp2}
\end{align}
where $(\zeta_1)$ follows from the fact that lines $5-8$ of Algorithm \ref{algo:sensing-altmin-improved} never increases $\twonorm{\aff\left(M-\Uw^t_{1:i}(\Vw^t_{1:i})^\dag\right)}$.
Using \eqref{eq:sam_bc_uperp1}, \eqref{eq:sam_bc_uperp2}, and $\frac{\so_i}{\so_{i+1}} \geq 5\sqrt{k}$, we obtain the following bound:
\begin{equation}
  \label{eq:sam_bc_uperp3}
  \|\left(\Up^t\right)^\dag \Uo_{1:i}\|_2\leq \frac{1}{2} \; \forall \;t.
\end{equation}
% Now, we first study the error incurred by the $T$-th iterate of the $i$-th stage:
% \begin{equation}
% \label{eq:sam_bc_err_i}
% \left\|M - \Uw^T_{1:i}(\Vw^T_{1:i})^\dag\right\|_F^2 \leq 2\left(\frob{U_{\perp}^{\dag} \Uo_{1:i}\So_{1:i}}^2
% 	+ \frob{\Uo_{i+1:k}\So_{i+1:k}\left(\Vo_{i+1:k}\right)^{\dag}}^2\right),
% \end{equation}
% where $U_{\perp}=I-\Uw^T_{1:i}(\Uw^T_{1:i})^\dag$ is a basis of the subpace perpendicular to $\Uw^T_{1:i}$.  
% We now bound the first term of the RHS of \eqref{eq:sam_bc_err_i}. Now, 
% \begin{align}
%   \|\Up^\dag M\|_2=\|\Up^\dag M\|_2
% \end{align}
% To this end, we first show that: 
% \begin{align}\label{eqn:UperpUstartwonorm}
%   \twonorm{U_{\perp}^{\dag}\Uo_{1:i}} \leq \frac{1}{2}.
% \end{align}
% %We first show that under the given assumptions, we have the following bound on the distance between the subspaces spanned by $\widehat{U}$ and $\Uo_{1-i}$% (where $\Uo_{1-i}$ representsthe top $i$ singular vectors of $M$):
% %\begin{align}\label{eqn:UperpUstartwonorm}
% %  \twonorm{U_{\perp}^{\dag}\Uo_{1-i}} \leq \frac{1}{2}.
% %\end{align}
% To see that, consider the matrix $N \eqdef M - \widehat{U}\widehat{V}^{\dag}$. We know that:
% \begin{align}\label{eqn:Ntwonorm}
%   \twonorm{N}^2 \leq \frob{N}^2 &\leq \sum_{j=i+1}^k \sigma_j^2 + \frac{1}{100} \sigma_i^2 \nonumber\\
% 	&\leq k \sigma_{i+1}^2 + \frac{1}{100} \sigma_i^2 \nonumber\\
% 	&\leq \frac{1}{20} \sigma_i^2,
% \end{align}
% where we used the fact that $\frac{\sigma_i}{\sigma_{i+1}} \geq 5\sqrt{k}$. Now choose a unit vector $z\in U_{\perp}$ such that
% $\twonorm{U_{\perp}^{\dag}\Uo_{1-i}} = \twonorm{z^{\dag} \Uo_{1-i}}$. We have, 
% \begin{align}\label{eqn:UperpUstartwonorm-intermediate}
%   \twonorm{z^{\dag} N} &= \twonorm{z^{\dag} \left(\Uo_{1-i}\Sigma_{1-i} (\Vo)_{1-i}^{\dag} + \Uo_{i+1-k}\Sigma_{i+1-k} (\Vo)_{i+1-k}^{\dag}\right)} \nonumber\\
% 	&\geq \twonorm{z^{\dag} \Uo_{1-i}\Sigma_{1-i} (\Vo)^{\dag}_{1-i}} - \twonorm{z^{\dag}\Uo_{i+1-k}\Sigma_{i+1-k} (\Vo)^{\dag}_{i+1-k}} \nonumber\\
% 	&\geq \sigma_i \twonorm{U_{\perp}^{\dag}\Uo_{1-i}} - \sigma_{i+1} \nonumber\\
% 	&\geq \sigma_i \left(\twonorm{U_{\perp}^{\dag}\Uo_{1-i}} - \frac{1}{5\sqrt{k}}\right).
% \end{align}
% Putting \eqref{eqn:Ntwonorm} and \eqref{eqn:UperpUstartwonorm-intermediate} together gives us \eqref{eqn:UperpUstartwonorm}.
Now, we consider the update equation for $\Vht$:
$$\Vht=\argmin_{\Vw}\|\aff(\Uw^t_{1:i}\Vw-\Uo_{1:i}\So_{1:i}(\Vo_{1:i})^\dag-\Uo_{i+1:k}\So_{i+1:k}(\Vo_{i+1:k})^\dag)\|_2^2.$$
Note that, the update is same as noisy case with noise matrix $N=\Uo_{i+1:k}\So_{i+1:k}(\Vo_{i+1:k})^\dag$ from \eqref{eq:noisy_update_k}:
\begin{align}\label{eqn:updateeqn-stagewise}
  \Vht = \Vo_{1:i}\So_{1:i}(\Uo_{1:i})^\dag\Ut_{1:i}-F+ \Vo_{i+1:k}\So_{i+1:k}(\Uo_{i+1:k})^\dag\Ut_{1:i} - G,
\end{align}
where $F$ and $G$ are given by \eqref{eq:F_k}, \eqref{eq:G_k}. %the terms corresponding to $\Vo_{1-i}\So_{1-i}(\Uo_{1-i})^\dag$ and $\Vo_{i+1-k}\So_{i+1-k}(\Uo_{i+1-k})^\dag$
%respectively. We will first show a decay for the quantity $\frob{U_{\perp}^{\dag}\Uo_{1-i}\So_{1-i}}$. 
Multiplying \eqref{eqn:updateeqn-stagewise}
from the left by $V_{\perp}^{\dag}=I-\Vt(\Vt)^\dag$, we obtain:
\begin{align}
  &0 = V_{\perp}^{\dag}\Vht_{1:i} = V_{\perp}^{\dag} \left(\Vo_{1:i}\So_{1:i}(\Uo_{1:i})^\dag\Ut_{1:i} - F + \Vo_{i+1:k}\So_{i+1:k}(\Uo_{i+1:k})^\dag\Ut_{1:i} - G\right) \nonumber\\
\Rightarrow &V_{\perp}^{\dag} \Vo_{1:i}\So_{1:i}(\Uo_{1:i})^\dag\Ut_{1:i} = V_{\perp}^{\dag} \left( F - \Vo_{i+1:k}\So_{i+1:k}(\Uo_{i+1:k})^\dag\Ut_{1:i} + G\right) \nonumber\\
\Rightarrow &\frob{V_{\perp}^{\dag} \Vo_{1:i}\So_{1:i}(\Uo_{1:i})^\dag\Ut_{1:i}} \leq \frob{F} + \frob{V_{\perp}^{\dag}\Vo_{i+1:k}\So_{i+1:k}(\Uo_{i+1:k})^\dag\Ut_{1:i}} + \frob{G} \nonumber\\
\Rightarrow &\frob{V_{\perp}^{\dag} \Vo_{1:i}\So_{1:i}} \leq \frac{1}{\sigma_{\textrm{min}}\left((\Uo_{1:i})^\dag\Ut_{1:i}\right)} \left(\frob{F} + \frob{V_{\perp}^{\dag}\Vo_{i+1:k}\So_{i+1:k}(\Uo_{i+1:k})^\dag\Ut_{1:i}} + \frob{G}\right), \label{eqn:weighted-dist-decay1}
\end{align}
where the last inequality follows using the fact that $\sigma_{\min}(A)\|B\|_F\leq \|AB\|_F$. Using Lemma~\ref{lemma:sensing-twonorm-G-noisycase}, and a modification of Lemma~\ref{lemma:sensing-twonorm-F}, we get:
%To show the decay of $\frob{V_{\perp}^{\dag} \Vo_{1-i}\So_{1-i}}$, we use the following inequalities:
\begin{align}
  &\frob{F} \leq \delta_{2k} \frob{U_{\perp}^{\dag} \Uo_{1:i}\So_{1:i}},\ \ \ \frob{G} \leq \delta_{2k} \frob{U_{\perp}^{\dag} \Uo_{i+1:k}\So_{i+1:k}} \leq \delta_{2k}\sqrt{k}\sigma_{i+1} \label{eqn:frob-G-bound}.
%  &\frob{V_{\perp}^{\dag}\Vo_{i+1-k}\So_{i+1-k}(\Uo_{i+1-k})^\dag\Ut} \leq \frob{\Vo_{i+1-k}\So_{i+1-k}(\Uo_{i+1-k})^\dag} = \sqrt{\sum_{j=i+1}^k \sigma_j^2}, \label{eqn:frob-smallersingvalues-bound}
\end{align}
Using \eqref{eqn:weighted-dist-decay1}, \eqref{eqn:frob-G-bound}, and the fact that $\sigma_{\min}(U_\perp^\dag\Uo_{1:i})=\sqrt{1-\|U_\perp^\dag\Uo_{1:i}\|_2^2}$,
%%where \eqref{eqn:frob-F-bound} (and similarly \eqref{eqn:frob-G-bound}) follows from the proof of Lemma \ref{lemma:sensing-twonorm-F}.
%Plugging \eqref{eqn:UperpUstartwonorm}, \eqref{eqn:frob-F-bound}, \eqref{eqn:frob-G-bound} and \eqref{eqn:frob-smallersingvalues-bound}
%into \eqref{eqn:weighted-dist-decay1}, we obtain the following:
\begin{align*}
  \frob{V_{\perp}^{\dag} \Vo_{1:i}\So_{1:i}} 
	&\leq \frac{2}{\sqrt{3}} \left(\delta_{2k} \frob{U_{\perp}^{\dag} \Uo_{1:i}\So_{1:i}}
		+ \sqrt{\sum_{j=i+1}^k \soj} + \delta_{2k}\sqrt{k}\sigma_{i+1} \right).
%	&\leq \frac{2}{\sqrt{3}} \left(\delta \frob{U_{\perp}^{\dag} \Uo_{1-i}\So_{1-i}}%
%		+ \sqrt{\sum_{j=i+1}^k \sigma_j^2} + \frac{1}{200\sqrt{k}} \sqrt{\sum_{j=i+1}^k \sigma_j^2} \right).
\end{align*}
Assuming $\frob{U_{\perp}^{\dag} \Uo_{1:i}\So_{1:i}} > 2 \sqrt{\sum_{j=i+1}^k \sigma_j^2}$, we obtain:
\begin{align}
  \frob{V_{\perp}^{\dag} \Vo_{1:i}\So_{1:i}} 
	&\leq \frac{2}{3} \frob{U_{\perp}^{\dag} \Uo_{1:i}\So_{1:i}}. 
\end{align}
Using similar analysis, we can show that, $$\frob{U_{\perp}^{\dag} \Uo_{1:i}\So_{1:i}} 
	\leq \frac{2}{3} \frob{V_{\perp}^{\dag} \Vo_{1:i}\So_{1:i}}. $$
So after $T\geq 8\log(k\so_i)$ iterations, we have:
\begin{align*}
  \frob{U_{\perp}^{\dag} \Uo_{1:i}\So_{1:i}}^2 \leq 4\sum_{j=i+1}^k \soj. 
\end{align*}
Using the above inequality, we now bound the error after $T\geq  8\log(k\so_i)$ iterations of the $i$-th stage: 
\begin{align}\label{eq:sam_err_bc_t_1}
  \left\|M - \Uw^T_{1:i}(\Vw^T_{1:i})^{\dag}\right\|_F
	\leq \left\| \Uo_{1:i}\So_{1:i}\left(\Vo_{1:i}\right)^{\dag} - \Uw^T_{1:i}(\Vw^T_{1:i})^{\dag}\right\|_F + \frob{\Uo_{i+1:k}\So_{i+1:k}\left(\Vo_{i+1:k}\right)^{\dag}}.
\end{align}

For the first term, we have:
\begin{align}
	&\left\| \Uo_{1:i}\So_{1:i}\left(\Vo_{1:i}\right)^{\dag} - \Uw^T_{1:i}(\Vw^T_{1:i})^{\dag}\right\|_F^2 \nonumber\\
	&= \frob{\Uo_{1:i}\So_{1:i}\left(\Vo_{1:i}\right)^{\dag} - U^T_{1:i}\left(U^T_{1:i}\right)^{\dag}\Uo_{1:i}\So_{1:i}\left(\Vo_{1:i}\right)^{\dag}
		+ U^T_{1:i}\left(U^T_{1:i}\right)^{\dag}\Uo_{1:i}\So_{1:i}\left(\Vo_{1:i}\right)^{\dag} - \Uw^T_{1:i}(\Vw^T_{1:i})^{\dag}}^2 \nonumber\\
	&= \frob{\left(I-U^T_{1:i}\left(U^T_{1:i}\right)^{\dag}\right)\Uo_{1:i}\So_{1:i}\left(\Vo_{1:i}\right)^{\dag} }^2
		+ \frob{U^T_{1:i}\left(U^T_{1:i}\right)^{\dag}\Uo_{1:i}\So_{1:i}\left(\Vo_{1:i}\right)^{\dag} - \Uw^T_{1:i}(\Vw^T_{1:i})^{\dag}}^2 \nonumber\\
	&\stackrel{(\zeta_1)}{\leq} \frob{U_{\perp}^{\dag} \Uo_{1:i}\So_{1:i}}^2 + \frob{U_{1:i}^T\left(F+G-\left(\Ut_{1:i}\right)^{\dag}\Uo_{i+1:k}\So_{i+1:k}(\Vo_{i+1:k})^\dag\right)}^2 \nonumber\\
	&= \frob{U_{\perp}^{\dag} \Uo_{1:i}\So_{1:i}}^2 + \frob{F+G-\left(\Ut_{1:i}\right)^{\dag}\Uo_{i+1:k}\So_{i+1:k}(\Vo_{i+1:k})^\dag}^2 \nonumber\\
	&= \frob{U_{\perp}^{\dag} \Uo_{1:i}\So_{1:i}}^2 + 3\frob{F}^2+3\frob{G}^2+3\frob{\Uo_{i+1:k}\So_{i+1:k}(\Vo_{i+1:k})^\dag}^2 \nonumber\\
	&\stackrel{(\zeta_2)}{\leq} (1+3\delta_{2k}^2) \frob{U_{\perp}^{\dag} \Uo_{1:i}\So_{1:i}}^2 + 3(1+\delta_{2k}^2) \frob{\Uo_{i+1:k}\So_{i+1:k}(\Vo_{i+1:k})^\dag}^2 \nonumber\\
	&\leq 8k (\sigma_{i+1}^*)^2,
%	&\leq \max\left(\epsilon,\sum_{j=i+2}^k \sigma_j^2 + 9k \sigma_{i+1}^2\right) + \sum_{j=i+1}^k \sigma_j^2 \\
%	&\leq 2\max\left(\epsilon,\sum_{j=i+2}^k \sigma_j^2 + 9k \sigma_{i+1}^2\right).
\label{eq:sam_err_bc_t_2}
\end{align}
where $(\zeta_1)$ follows from \eqref{eqn:updateeqn-stagewise} and $(\zeta_2)$ follows from \eqref{eqn:frob-G-bound}.
Using \eqref{eq:sam_err_bc_t_1} and \eqref{eq:sam_err_bc_t_2}, we obtain the following bound:
\begin{align}\label{eq:sam_err_bc_t}
  \left\|M - \Uw^T_{1:i}(\Vw^T_{1:i})^{\dag}\right\|_F \leq 4\sqrt{k} \sigma_{i+1}^*.
\end{align}
Hence Proved.
%Finally, we note that $\Uw^0_{1:i+1}$, $\Vw^0_{1:i+1}$, the initial points of $i+1$-th stage, are obtained using one step of SVP algorithm \cite{JainMD10} with the latest iterate being $X_t= \Uw^T_{1:i}(\Vw^T_{1:i})^{\dag}$. Lemma now follows by applying Lemma 2.1 of \cite{JainMD10}. %The second part of the lemma now follows from the linear convergence of SVP.
\end{proof}
\section{Matrix Completion: Proofs}\label{app:mc}
\begin{proof}[Proof Of Theorem~\ref{thm:sampling-altmin-main}]
Using Theorem~\ref{thm:sampling-altmin-geomconv}, after $O(\log (1/\epsilon))$ iterations, we get: $$\dist(\Ut, \Uo)\leq \epsilon,\ \dist(\Vt, \Vo)\leq \epsilon.$$
Now, using  \eqref{eq:updatecomp_k}, the residual after $t$-th step is given by:
$$M-\Ut(\Vht)^\dag=(I-\Ut(\Ut)^\dag)M-\Ut F^\dag.$$
%Combining the above two equations, we get: 
That is, 
$$\|M-\Ut(\Vht)^\dag\|_F\leq \|(I-\Ut(\Ut)^\dag)M\|_F+\|F\|_F\leq \sqrt{k}\|(I-\Ut(\Ut)^\dag)U^*\So\|_2+\|F\|_F\leq \sqrt{k}\so_1\dist(\widehat{U}^t, \Uo).$$%\sqrt{k}\sigma_1^*\epsilon+\sqrt{k}\sigma_1^*\epsilon\leq 2\sqrt{k}\|M\|_F\epsilon,$$
Now, using the fact that $\dist(\Ut, \Uo)\leq \epsilon$ and the above equation, we get: 
$$\|M-\Ut(\Vht)^\dag\|_F\leq \sqrt{k}\so_1\epsilon+\|F\|_F\stackrel{\zeta_1}{\leq}\sqrt{k}\so_1\epsilon+\so_1\sqrt{k}\epsilon\leq 2\so_1\sqrt{k}\epsilon,$$
where $\zeta_1$ follows by Lemma~\ref{lemma:Fbound_comp} and setting $\delta_{2k}$ appropriately. Theorem \ref{thm:sampling-altmin-main} now follows by setting $\epsilon'=2\sqrt{k}\|M\|_F\epsilon$.
\end{proof}
\subsection{Initialization: Proofs}\label{app:comp_init}
\begin{proof}[Proof Of Lemma~\ref{lemma:svd-clipping-jointguarantee}]
  From Lemma \ref{lemma:first-step-svd}, we see that $U^0$ obtained after step $3$ of Algorithm~\ref{algo:completion-altmin} satisfies: $\dist\left(U^{0},U^*\right)\leq \frac{1}{64 k}$. Lemma now follows by using the above mentioned observation with Lemma~\ref{lemma:clipping-incoherence-distance}.\end{proof}
% gives us the desired Hence, $U^{0}$ satisfies the hypothesis of Lemma \ref{lemma:clipping-incoherence-distance}. Using Lemma \ref{lemma:clipping-incoherence-distance} gives us the desired result.
We now provide the two results used in the above lemma. %See Appendix~\ref{app:comp_init} for a detailed proof of both the lemmas. 
\begin{lemma}\label{lemma:first-step-svd}
After step $3$ in Algorithm \ref{algo:completion-altmin}, whp we have:
  \begin{align*}
    \dist\left(U^{0},U^*\right) \leq \frac{1}{64 k}
  \end{align*}
\end{lemma}
\begin{proof}
  From Theorem 3.1 in \cite{KeshavanOM2009}, we have the following result:
  \begin{align*}
    \twonorm{M - M_k} \leq C \left(\frac{k}{p\sqrt{mn}}\right)^{\frac{1}{2}}\|M\|_F.
  \end{align*}
Let $U^{(0)}\Sigma V^{\dag}$ be the top $k$ singular components of $M_k$.
We also have:
\begin{align*}
  \twonorm{M - M_k}^2 &= \twonorm{U^* \Sigma^* (V^*)^{\dag} - U^{(0)}\Sigma V^{\dag}}^2\\
	&= \left\|U^* \Sigma^* (V^*)^{\dag} - U^{(0)}\left(U^{(0)}\right)^{\dag}U^* \Sigma^* (V^*)^{\dag}
	  + U^{(0)}\left(U^{(0)}\right)^{\dag}U^* \Sigma^* (V^*)^{\dag} - U^{(0)}\Sigma V^{\dag}\right\|_2^2 \\
	&= \left\|\left(I-U^{(0)}\left(U^{(0)}\right)^{\dag}\right)U^* \Sigma^* (V^*)^{\dag}
	  + U^{(0)} \left(\left(U^{(0)}\right)^{\dag}U^* \Sigma^* (V^*)^{\dag} - \Sigma V^{\dag}\right)\right\|_2^2 \\
	&\stackrel{(\zeta_1)}{\geq} \twonorm{\left(I-U^{(0)}\left(U^{(0)}\right)^{\dag}\right)U^* \Sigma^* (V^*)^{\dag} }^2 = \twonorm{\left(U_{\perp}^{(0)}\right)^{\dag}U^* \Sigma^*}^2 \geq \left(\sigma_k^*\right)^2 \twonorm{\left(U_{\perp}^{(0)}\right)^{\dag}U^*}^2,
\end{align*}
where $(\zeta_1)$ follows from the fact that the column space of the first two terms in the equation is $U_{\perp}^{(0)}$ where as the column space
of the last two terms is $U^{(0)}$. Using the above two inequalities, we get:
\begin{align*}
  \twonorm{\left(U_{\perp}^{(0)}\right)^{\dag}U^*} \leq C\cdot\frac{\so_1}{\sigma_k^*} \cdot \frac{k}{\sqrt{mp}}
	\leq \frac{1}{10^4 k},
\end{align*}
if $p > \frac{C'k^4\log n}{m}\cdot \frac{(\so_1)^2}{(\so_k)^2}$ for a large enough constant $C'$.
\end{proof}
\begin{lemma}\label{lemma:clipping-incoherence-distance}(Analysis of step $4$ of Algorithm \ref{algo:completion-altmin})
  Suppose $U^*$ is incoherent with parameter $\mu$ and $U$ is an orthonormal column matrix such that $\dist\left(U,U^*\right) \leq \frac{1}{64k}$.
Let $U^c$ be obtained from $U$ by setting all entries greater than $\frac{2\mu\sqrt{k}}{\sqrt{n}}$  to zero. Let $\widetilde{U}$ be an orthonormal basis of $U^c$. Then,\vspace*{-4pt}\setlength{\topsep}{-5pt}\setlength{\itemsep}{-5pt}
\begin{itemize}\setlength{\topsep}{-5pt}\setlength{\itemsep}{-5pt}
  \item	$\dist\left(\widetilde{U},U^*\right) \leq 1/2$ and
  \item	$\widetilde{U}$ is incoherent with parameter $4\mu \sqrt{k}$.\vspace*{-5pt}
\end{itemize}
\end{lemma}
\begin{proof}
  Since $\dist\left(U,U^*\right) \leq d$, we have that for every $i$, $\exists \breve{u}_i \in \textrm{Span}(U^*), \twonorm{\breve{u}_i}=1$ such that
 $\iprod{u_i}{\breve{u}_i} \geq \sqrt{1-d^2}$. Also, since $\breve{u}_i \in \textrm{Span}(U^*)$, we have that $\breve{u}_i$
 is incoherent with parameter $\mu \sqrt{k}$:
\begin{align*}
  \twonorm{\breve{u}_i}=1 \mbox{ and }
  \infnorm{\breve{u}_i} \leq \frac{\mu \sqrt{k}}{\sqrt{m}}.
\end{align*}
Let $u_i^c$ be the vector obtained by setting all the elements of $u_i$ with magnitude greater than $\frac{2\mu\sqrt{k}}{\sqrt{m}}$ to zero and let
$u_i^{\overline{c}} \eqdef u_i - u_i^c$. 
Now, note that if for element $j$ of $u_i$ we have $\left|u_{i}^j\right| > \frac{2\mu\sqrt{k}}{\sqrt{m}}$, then, $|(u_{i}^c)^j-\breve{u}_i^j|=\left|\breve{u}_i^j\right|\leq \frac{\mu\sqrt{k}}{\sqrt{m}} \leq \left|u_i^j - \breve{u}_i^j\right|$. Hence, 
\begin{align*}
  \twonorm{u_i^c - \breve{u}_i} \stackrel{(\zeta_1)}{\leq} \twonorm{u_i - \breve{u}_i}
	= \left(\twonorm{u_i}^2 + \twonorm{\breve{u}_i}^2 - 2 \iprod{u_i}{\breve{u}_i}\right)^{\frac{1}{2}} \leq \sqrt{2}d,
\end{align*}
%where $(\zeta_1)$ follows from the fact that if for element $j$ of $u_i$ we have $\left|u_{i}^j\right| > \frac{2\mu\sqrt{k}}{\sqrt{n}}$,
%then $\left|u_i^j - \breve{u}_i^j\right| \geq \left|u_i^j\right| - \left|\breve{u}_i^j\right| \geq \frac{\mu\sqrt{k}}{\sqrt{n}} \geq \left|\breve{u}_i^j\right|$.
This also implies the following:
\begin{align*}
  \twonorm{u_i^c} &\geq \twonorm{\breve{u}_i} - \sqrt{2}d = 1 - \sqrt{2}d\ \ \mbox{, and }\\
  \twonorm{u_i^{\overline{c}}} &\leq \sqrt{1 - \twonorm{u_i^c}^2} \leq \sqrt{2d(\sqrt{2}-d)} \leq 2 \sqrt{d},\ \  \mbox{ for } d < \frac{1}{\sqrt{2}}.
\end{align*}

Let $U^c = \widetilde{U}\Lambda^{-1}$ (QR decomposition). Then, for any $u_{\perp}^* \in \textrm{Span}(U_{\perp}^*)$ we have:
\begin{align*}
  \twonorm{\left(u_{\perp}^*\right)^{\dag} \widetilde{U}} &= \twonorm{\left(u_{\perp}^*\right)^{\dag} U^c \Lambda} \leq \twonorm{\left(u_{\perp}^*\right)^{\dag} U^c} \twonorm{\Lambda} \leq \left(\twonorm{\left(u_{\perp}^*\right)^{\dag} U} + \twonorm{\left(u_{\perp}^*\right)^{\dag} U^{\overline{c}} }\right) \twonorm{\Lambda} \\
	&\leq \left(d + \twonorm{ U^{\overline{c}} }\right)\twonorm{\Lambda}
	\leq \left(d + \frob{ U^{\overline{c}} }\right)\twonorm{\Lambda}
	\leq \left(d + 2\sqrt{kd} \right)\twonorm{\Lambda}
	\leq 3\sqrt{kd} \twonorm{\Lambda}.
\end{align*}
We now bound $\twonorm{\Lambda}$ as follows:
\begin{align*}
  \twonorm{\Lambda}^2 = \frac{1}{\sigma_{\textrm{min}}\left(\Lambda^{-1}\right)^2} 
	= \frac{1}{\sigma_{\textrm{min}}\left(\widetilde{U}\Lambda^{-1}\right)^2} 
	= \frac{1}{\sigma_{\textrm{min}}\left(U^c\right)^2} 
	\leq \frac{1}{1 - \twonorm{U^{\overline{c}}}^2} 
	\leq \frac{1}{1 - 4kd} \leq 4/3,
\end{align*}
where we used the fact that $d < \frac{1}{16k}$.
So we have:
\begin{align*}
  \twonorm{\left(u_{\perp}^*\right)^{\dag} \widetilde{U}} \leq  3 \sqrt{kd} \cdot 4/3 = 4\sqrt{kd}.
\end{align*}
This proves the first part of the lemma.

Incoherence of $\widetilde{U}$ follows using the following set of inequalities: 
%For the second part about incoherence of $\widetilde{U}$, we observe the following bound on each entry of $\widetilde{U}$:
\begin{align*}
\mu(\widetilde{U})=\frac{\sqrt{m}}{\sqrt{k}}\max_i \|e_i^\dag \widetilde{U}\|_2 \leq \frac{\sqrt{m}}{\sqrt{k}}\max_i \|e_i^\dag U^c \Lambda\| \leq \frac{\sqrt{m}}{\sqrt{k}}\max_i \|e_i^\dag U^c\|_2\|\Lambda\|_2\leq 4\mu \sqrt{k}. 
\end{align*}
\end{proof}
\subsection{Rank-$1$ Matrix Completion: Proofs}\label{app:rank1_comp}
\begin{proof}[Proof Of Lemma~\ref{lem:errb_comp1}]
Using the definition of spectral norm, 
\begin{align*}
  \|B^{-1}\left(\ip{\uo}{\ut}B-C\right)\vo\|_2\leq \|B^{-1}\|_2\|(\ip{\uo}{\ut}B-C)\vo\|_2.%\label{eq:error1_comp}
\end{align*}
As $B$ is a diagonal matrix, $\|B^{-1}\|_2= \frac{1}{\min_iB_{ii}}\leq \frac{1}{1-\delta_2}$, where the last inequality follows using Lemma~\ref{lem:errb_comp2}. The lemma now follows using the above observation and Lemma~\ref{lem:bccompbound}.% with \eqref{eq:error1_comp}. 
\end{proof}
\begin{lemma}\label{lem:errb_comp2}
Let $M=\so\uo(\vo)^\dag$, $p$, $\Omega$, $\ut$ be as defined in Lemma~\ref{lem:errb_comp1}. 
%  Let $\ut$ be a $\mu_1$ incoherent unit vector and let $\Omega$ be sampled uniformly and independently of $\ut$, with $p\geq \frac{9\mu_1^2\log n}{n\delta_2^2}$. 
Then, w.p. at least $1-\frac{1}{n^3}$,
$$\left|\frac{\sum_{i:(i,j)\in \Omega} (\ut_i)^2}{p}-1\right| \leq \delta_2,\ \left|\frac{\sum_{i:(i,j)\in \Omega} \ut_i \uo_i}{p}-\ip{\ut}{\uo}\right|\leq \delta_2.$$ %$$$$ %Moreover, %if $\uo$ is also a $\mu_1$ incoherent unit vector then with probability greater than $1-\frac{1}{n^3}$, we have$\ip{\ut}{\uo}-\delta_2\leq \frac{\sum_{i:(i,j)\in \Omega} \ut_i \uo_i}{p} \leq \ip{\ut}{\uo}+\delta_2.$
\end{lemma}
\begin{proof}
Since the first part of the lemma is a direct consequence of the second part, we will prove only the second part.
Let $\delta_{ij}$ be a Bernoulli random variable that indicates membership of index $(i,j)\in \Omega$. That is, $\delta_{ij}=1$ w.p. $p$ and $0$ otherwise. Define $Z_j=\frac{1}{p}\sum_{i} \delta_{ij} \ut_i \uo_i$. Note that $\mathbb{E}[Z_j]=\ip{\ut}{\uo}$. Furthermore, $\mathbb{E}[Z_j^2]=\left(\frac{1}{p}-1\right)\sum_{i}(\ut_i \uo_i)^2\leq \frac{\mu_1^2}{mp}$ and $\max_i |\ut_i \uo_i|\leq \frac{\mu_1^2}{m}$. Using Bernstein's inequality, we get:  
\begin{align}
  \Pr(\left|Z_j-\ip{\ut}{\uo}\right|>\delta_2)\leq \exp\left(-\frac{\delta_2^2mp/2}{\mu_1^2+\mu_1^2\delta_2/3}\right). 
\end{align}
Using union bound (for all $j$) and for $p\geq \frac{9\mu_1^2\log n}{m\delta_2^2}$, w.p. $1-\frac{1}{n^3}$: $\forall j, \ip{\ut}{\uo}-\delta_2\leq Z_j \leq \ip{\ut}{\uo}+\delta_2$.
\end{proof}
\begin{lemma}\label{lem:bccompbound}
Let $M=\so\uo(\vo)^\dag$, $p$, $\Omega$, $\ut$ be as defined in Lemma~\ref{lem:errb_comp1}. 
%  Let $u^*, v^*$ be $\mu$ incoherent unit vectors and $\ut$ be a $\mu_1$ incoherent unit vector (assume $\mu_1\geq \mu$). Let $\Omega$ be sampled uniformly and independently of $\ut$, with $p\geq \frac{C\mu_1^2\log n}{n\delta_2^2}$ ($C>0$ is a global constant). 
Then, w.p. at least $1-\frac{1}{n^3}$,
$$\|(\ip{\uo}{\ut}B-C)\vo\|_2\leq \delta_2\sqrt{1-\ip{\uo}{\ut}^2}. $$
\end{lemma}
\begin{proof}
Let $x\in \rn$ be a unit vector. Then, $\forall x$: 
\begin{align}
x^\dag(\ip{\uo}{\ut}B-C)\vo&=\frac{1}{p}\sum_{ij\in \Omega}x_j\vo_j(\ip{\uo}{\ut}(\ut_i)^2-\ut_i\uo_i)\nonumber\\
&\stackrel{(\zeta_1)}{\leq} \frac{1}{p}C\sqrt{mp}\sqrt{\sum_j x_j^2(\vo_j)^2}\sqrt{\sum_i(\ip{\uo}{\ut}(\ut_i)^2-\ut_i\uo_i)^2}, \nonumber\\
&\stackrel{(\zeta_2)}{\leq}\frac{1}{p}C\frac{\sqrt{mp}\mu_1^2}{n} \sqrt{1-\ip{\uo}{\ut}^2},
%=\sum_{(i,j)\in \Omega}(\vo_j)^2(\ip{\uo}{\ut}(\ut_i)^2-\ut_i\uo_i)^2\leq \sum_{(i,j)\in \Omega}(\vo_j)^2(\ut_i)^2(\ip{\uo}{\ut}\ut_i-\uo_i)^2\nonumber\\&\leq \sum_{(i,j)\in \Omega}(\vo_j)^2(\ip{\uo}{\ut}\ut_i-\ut_i)^2=\|P_\Omega(z(\vo)^\dag\|_F^2,
\end{align}
where $C>0$ is a global constant and $(\zeta_1)$ follows by using a modified version of Lemma 6.1 by \cite{KeshavanOM2009} (see Lemma~\ref{lem:komlem}) and $(\zeta_2)$ follows by using incoherence of $\vo$ and $\ut$. Lemma now follows by observing that $\max_{x, \|x\|_2=1} x^\dag(\ip{\uo}{\ut}B-C)\vo=\|(\ip{\uo}{\ut}B-C)\vo\|_2$ and $p>\frac{C\mu_1^2\log n}{m\delta_2^2}$.
\end{proof}
\begin{proof}[Proof of Lemma~\ref{lem:incoherence_vt}]
Using \eqref{eq:compvtupdate} and using the fact that $B, C$ are diagonal matrices:
$$\widehat{v}^{t+1}_j=\so\ip{\ut}{\uo}\vo_j-\frac{\so}{B_{jj}} \left(\ip{\ut}{\uo}B_{jj}-C_{jj}\right)\vo_j.$$%\leq \frac{\so\ip{\ut}{\uo}}{\sqrt{n}}-C\delta_2\sqrt{1-\ip{\ut}{\uo}^2}.$$
We bound the largest magnitude of elements in $\widehat{v}^{t+1}$ as follows. For every $j\in[n]$, we have:
\begin{align*}
  \left|\widehat{v}^{t+1}_j\right| &\leq \left|\so\ip{\ut}{\uo}\vo_j\right| + \left|\frac{\so}{B_{jj}} \left(\ip{\ut}{\uo}B_{jj}-C_{jj}\right)\vo_j\right| \\
	&\stackrel{(\zeta_1)}{\leq} \so \ip{\ut}{\uo} \frac{\mu}{\sqrt{n}} + \frac{\so}{1-\delta_2} \left(\ip{\ut}{\uo} \left(1+\delta_2\right) + \left(\ip{\ut}{\uo}+\delta_2\right)\right) \frac{\mu}{\sqrt{n}} \\
	&\leq \frac{\frac{3\so(1+\delta_2)\mu}{1-\delta_2}}{\sqrt{n}} \leq \frac{\so\mu_1}{2\sqrt{n}},
\end{align*}
where $(\zeta_1)$ follows from the fact that $1-\delta_2 \leq B_{jj} \leq 1+\delta_2$ and $\left|C_{jj}\right| \leq \left(\left|\ip{\ut}{\uo}\right|+\delta_2\right)$ (please refer Lemma \ref{lem:errb_comp2}).

Also, from \eqref{eq:ipvtvo_comp} we see that:
\begin{align*}
  \twonorm{\widehat{v}^{t+1}} \geq \ip{\widehat{v}^{t+1}}{\vo} &\geq \so\ip{\ut}{\uo}-2\so\delta_2\sqrt{1-\ip{\ut}{\uo}^2} \\
	&\geq \so\ip{u^0}{\uo}-2\so\delta_2\sqrt{1-\ip{u^0}{\uo}^2} \\
	&\stackrel{(\zeta_1)}{\geq} \frac{\so}{2},
\end{align*}
where $(\zeta_1)$ follows from the fact that $\dist\left(u^0,u^*\right) \leq \frac{3}{50}$ (please refer Lemma \ref{lemma:svd-clipping-jointguarantee}).
Using the above two inequalities, we obtain:
\begin{align*}
  \infnorm{v^{t+1}} = \frac{\infnorm{\widehat{v}^{t+1}}}{\twonorm{\widehat{v}^{t+1}}} \leq \frac{\left(\frac{\so\mu_1}{2\sqrt{n}}\right)}{\left(\frac{\so}{2}\right)}
	= \frac{\mu_1}{\sqrt{n}}.
\end{align*}
This finishes the proof.
%=\so \frac{\sum_{(i,j)\in \Omega}\ut_i\uo_i}{\sum_{(i,j)\in \Omega} (\ut_i)^2} \vo_j.$$
%Now, $Z=\sum_{(i,j)\in \Omega}\ut_i\uo_i=\frac{1}{p}\sum_{i}\delta_{ij}\ut_i\uo_i.$ Note that, $\E[Z]=\ip{\ut}{\uo}$, $\E[Z^2]=\sum_{i}(\ut_i\uo_i)^2/p\leq \mu^2/(np)\leq \frac{1}{\log n}$ and $\max_{i}|\ut_i\uo_i|/p\leq $. Using Bernstein's inequality: $$\Pr(|Z-\E[Z]|\geq \delta)\leq \exp\left(\frac{-\delta^2}{\frac{1}{\log n}+}\right).$$
\end{proof}

\begin{lemma}[Modified version of Lemma 6.1 of \cite{KeshavanOM2009}]
\label{lem:komlem}
Let $\Omega$ be a set of indices sampled uniformly at random from $[m]\times[n]$ with each element of $[m]\times[n]$ sampled
independently with probability $p\geq \frac{C\log n}{m}$. Then, w.p. at least $1-\frac{1}{n^3}$, $\forall x \in \mathbb{R}^m, y \in \rn$ s.t. $\sum_i x_i=0$, we have: %Also, let $\x\in \rm$ s.t. $\sum_i x_i =0$. Then, $
$\sum_{ij\in \Omega}x_iy_j\leq C\sqrt{\sqrt{mn}p}\|x\|_2\|y\|_2,$ where $C>0$ is a global constant. 
\end{lemma}
\subsection{General Rank-$k$ Matrix Completion: Proofs}\label{app:compk}
\begin{proof}[Proof of Lemma~\ref{lemma:completion-incoherence-Vt}]
  From the decoupled update equation, \eqref{eqn:sampling-updateeqn-decoupled}, we obtain:
\begin{align*}
    (\Vt)^{(j)} = (R^{(t+1)})^{-1}(D^j-(B^j)^{-1}(B^jD^j-C^j)) \Sigma^*(\Vo)^{(j)},\;\;\;\;  \;\; 1\leq j\leq n.
\end{align*}
We bound the two norm of the $(\Vt)^{(j)}$ as follows:
\begin{align*}
  \twonorm{(\Vt)^{(j)}}
	&\leq \frac{\sigma_1\twonorm{(\Vo)^{(j)}}}{\sigma_{\textrm{min}}\left(R^{(t+1)}\right)} \left( \twonorm{D^j} + \twonorm{(B^j)^{-1}(B^jD^j-C^j)} \right) \\
	&\leq \frac{\sigma_1\twonorm{(\Vo)^{(j)}}}{\sigma_{\textrm{min}}\left(R^{(t+1)}\right)} \left( \twonorm{D^j} + \frac{\twonorm{B^jD^j}+\twonorm{C^j}}{\sigma_{\textrm{min}}\left(B^j\right)} \right) \\
	  &\stackrel{(\zeta_1)}{\leq} \frac{\sigma_1 \frac{\mu \sqrt{k}}{\sqrt{n}}}{\sigma_k^*\sqrt{1-\dist^2\left(U^{(t)},U^*\right)}-\frac{\sigma_1^* \delta_{2k}k\dist(U^{(t)},U^*)}{1-\delta_{2k}}} \left( 1 + \frac{(1+\delta_{2k})+ (1+\delta_{2k}) }{1-\delta_{2k}} \right) \\
	  &\leq \frac{4\sigma_1 \frac{\mu \sqrt{k}}{\sqrt{n}}}{\sigma_k^*\sqrt{1-\dist^2\left(U^{(0)},U^*\right)}-\frac{\sigma_1^* \delta_{2k}k\dist(U^{(0)},U^*)}{1-\delta_{2k}}}
	  \leq \frac{\left(\frac{16\sigma_1^*\mu}{\sigma_k^*}\right)\sqrt{k}}{\sqrt{n}},
\end{align*}
where we used the following inequalities in $(\zeta_1)$:
\begin{align}
  \twonorm{(\Vo)^j} &\leq \frac{\mu \sqrt{k}}{\sqrt{n}}, \label{eqn:Vstar-incoherence-int} \\
  \sigma_{\textrm{min}}\left(R^{(t+1)}\right) &\geq 
	\sigma^*_k\sqrt{1-\dist^2\left(U^{(t)},U^*\right)}-\sigma^*_1 \delta_{2k}k\dist(U^{(t)},U^*),
		\label{eqn:Rsigmamin-int} \\
  \sigma_{\textrm{min}}\left(B^j\right) \geq 1-\delta_{2k} &\mbox{ and }
	\sigma_{\textrm{max}}\left(B^j\right) \leq 1+\delta_{2k}, \label{eqn:Bjeigenvaluebound-int} \\
  \sigma_{\textrm{max}}\left(C^j\right) &\leq 1+\delta_{2k} \mbox{ and } \label{eqn:Cjeigenvaluebound-int} \\
  \sigma_{\textrm{max}}\left(D^j\right) &\leq 1, \label{eqn:Djeigenvaluebound-int}
\end{align}
where
\eqref{eqn:Vstar-incoherence-int} follows from the incoherence of $\Vo$,
\eqref{eqn:Rsigmamin-int} follows from from an analysis similar to the proof of Lemma \ref{lemma:sensing-minsingval-Rt},
\eqref{eqn:Bjeigenvaluebound-int} follows from (the proof of) Lemma \ref{lemma:sampling-minsingval-Bl},
\eqref{eqn:Cjeigenvaluebound-int} follows from Lemma \ref{lemma:sampling-minsingval-Cjl} and finally
\eqref{eqn:Djeigenvaluebound-int} follows from the fact that $D^j = \left(U^t\right)^{\dag}U^*$ with
$U^t$ and $U^*$ being orthonormal column matrices.
\end{proof}
\begin{proof}[Proof of Lemma~\ref{lemma:Fbound_comp}]
  Note that,
\begin{align}
  \twonorm{F(\Sigma^*)^{-1}} \leq \frob{F(\Sigma^*)^{-1}} &= \twonorm{B^{-1} \left(BD-C\right)v^*} \nonumber\\
	&\leq \twonorm{B^{-1}}\twonorm{(BD-C)v^*} \nonumber\\
	&\leq \frac{\delta_{2k}}{1-\delta_{2k}}\dist(\Ut, \Uo),   \label{eqn:sampling-twonorm-F-1}
\end{align}
where the last inequality follows using Lemma~\ref{lemma:sampling-minsingval-Bl} and Lemma~\ref{lemma:sampling-minsingval-Cl}. 
\end{proof}
We now bound $\|B^{-1}\|_2$ and $\|C^j\|_2$, which is required by our bound for $F$ as well as for our incoherence proof. 
\begin{lemma}\label{lemma:sampling-minsingval-Bl}
Let $M, \Omega, p,$ and $\Ut$ be as defined in Theorem~\ref{thm:sampling-altmin-main} and Lemma~\ref{lemma:Fbound_comp}. 
%Suppose $p>\frac{6\mu_1^2k\log n}{n\delta_{2k}^2}$. Let $\Ut$ be the $t$-th step iterate; assume it to be $\mu_1$-incoherent. Also, let $\Uo$ be $\mu$-incoherent ($\mu_1\geq \mu$). 
Then, w.p. at least $1-\frac{1}{n^3}$: 
%  Let $B^j \in \mathbb{R}^{k\times k}$ be defined as: $B^j=\frac{1}{p}\sum_{(i,j)\in \Omega} (\Ut)^{(i)}{(\Ut)^{(i)}}^\dag$ where $\Ut$ is $t$-th step iterate and is $\mu_1$-incoherent. Then, for every $1\leq j\leq n$, w.p. at least $1-\frac{1}{n^3}$: 
\begin{align}\label{eqn:sampling-minsingval-Bl}
\|B^{-1}\|_2\leq \frac{1}{1-\delta_{2k}}.
%  \sigma_{\textrm{min}}\left(B^j\right) \geq \left(1-k\delta\right).
\end{align}
\end{lemma}
\begin{proof}[Proof of Lemma~\ref{lemma:sampling-minsingval-Bl}]
% Let $B^j \in \mathbb{R}^{k\times k}$ be defined as: $B^j=\frac{1}{p}\sum_{(i,j)\in \Omega} (\Ut)^{(i)}{(\Ut)^{(i)}}^\dag$. Then, 
We have:
$$\|B^{-1}\|_2=\frac{1}{\sigma_{min}(B)}=\frac{1}{\min_{x, \|x\|=1}x^\dag B x},$$
where $x\in \mathbb{R}^{nk}$. Let $x=vec(X)$, i.e.,  $x_p$ is the $p$-th column of $X$ and $x^j$ is the $j$-th row of $X$. Now, $\forall x$,
$$x^\dag B x=\sum_{j} (x^j)^\dag B^j (x^j)\geq min_j \sigma_{min}(B^j).$$
Lemma would follow using the bound on $\sigma_{min}(B^j), \forall j$ that we show below. 

\noindent {\bf Lower bound on $\sigma_{min}(B^j)$}: Consider any $w\in \R^k$ such that $\twonorm{w}=1$.
We have:
\begin{align*}
  Z=w^\dag  B^j w = \frac{1}{p}\sum_{i:(i,j)\in \Omega} \ip{w}{(\Ut)^{(i)}}^2= \frac{1}{p}\sum_{i}\delta_{ij}\ip{w}{(\Ut)^{(i)}}^2.
\end{align*}
Note that, $\E[Z]=w^\dag UU^\dag w=w^\dag w=1$ and $\E[Z^2]=\frac{1}{p}\sum_{i}\ip{w}{(\Ut)^{(i)}}^4\leq \frac{\mu_1^2k}{mp}\sum_{i}\ip{w}{(\Ut)^{(i)}}^2=\frac{\mu_1^2k}{mp},$ where the second last inequality follows using incoherence of $\Ut$. Similarly, $\max_{i}|\ip{w}{(\Ut)^{(i)}}^2|\leq \frac{\mu_1^2k}{mp}$. Hence, using Bernstein's inequality: 
$$\Pr(|Z-\E[Z]|\geq \delta_{2k})\leq \exp(-\frac{\delta_{2k}^2/2}{1+\delta_{2k}/3}\frac{mp}{\mu_1^2k}).$$
That is, by using $p$ as in the statement of the lemma with the above equation and using union bound, we get (w.p. $>1-1/n^3$): $\forall w,j\ \ w^\dag B^j w \geq 1-\delta_{2k}$. That is, $\forall j, \sigma_{min}(B^j)\geq (1-\delta_{2k})$. 
%Lemma now follows from variational definition of smallest singular value. 
\end{proof}
\begin{lemma}\label{lemma:sampling-minsingval-Cjl}
%Suppose $p>\frac{6\mu_1^2k\log n}{n\delta_{2k}^2}$. 
Let $M, \Omega, p,$ and $\Ut$ be as defined in Theorem~\ref{thm:sampling-altmin-main} and Lemma~\ref{lemma:Fbound_comp}. Also, let $C^j \in \mathbb{R}^{k\times k}$ be defined as: $C^j=\frac{1}{p}\sum_{i:(i,j)\in \Omega} (\Ut)^{(i)}{(\Uo)^{(i)}}^\dag$. Then, w.p. at least $1-\frac{1}{n^3}$: 
\begin{align}\label{eqn:sampling-minsingval-Cjl}
\|C^j\|_2\leq 1+\delta_{2k}, \forall j
%\|B^{-1}\|_2\leq \frac{1}{1-\delta_{2k}}.
%  \sigma_{\textrm{min}}\left(B^j\right) \geq \left(1-k\delta\right).
\end{align}
\end{lemma}
\begin{proof}[Proof of Lemma~\ref{lemma:sampling-minsingval-Cjl}]
Let $x\in \mathbb{R}^{k}$ and $y\in \mathbb{R}^k$ be two arbitrary unit vectors. Then, 
$$x^TC^jy=\frac{1}{p}\sum_{i:(i,j)\in \Omega}(x^\dag(\Ut)^{(i)})(y^\dag(\Uo)^{(i)}).$$
That is, $Z=x^TC^jy=\frac{1}{p}\sum_{i}\delta_{ij}(x^\dag(\Ut)^{(i)})(y^\dag(\Uo)^{(i)})$. Note that, $\E[Z]=x^\dag (\Ut)^\dag\Uo y$, $\E[Z^2]=\frac{1}{p}\sum_{i}(x^\dag(\Ut)^{(i)})^2(y^\dag(\Uo)^{(i)})^2\leq \frac{\mu^2}{mp}x^\dag(\Ut)^\dag\Ut x=\frac{\mu^2k}{mp}$ and $\max_i |(x^\dag(\Ut)^{(i)})(y^\dag(\Uo)^{(i)})|\leq \frac{\mu_1^2k}{m}$. Lemma now follows using Bernstein's inequality and using bound for $p$ given in the lemma statement. 
\end{proof}
Finally, we provide a lemma to bound the second part of the error term ($F$).
\begin{lemma}\label{lemma:sampling-minsingval-Cl}
Let $M, \Omega, p,$ and $\Ut$ be as defined in Theorem~\ref{thm:sampling-altmin-main} and Lemma~\ref{lemma:Fbound_comp}.
%Suppose $p>\frac{6\mu_1^2k^{\frac{5}{2}}\log n}{n\delta_{2k}^2}$. 
%  Let $C^j \in \mathbb{R}^{k\times k}$ be defined as: $C^j=\frac{1}{p}\sum_{(i,j)\in \Omega} (\Ut)^{(i)}{(\Uo)^{(i)}}^\dag$ where 
%Let $\Ut$ be the $t$-th step iterate; assume it to be $\mu_1$-incoherent. Also, let $\Uo$ be $\mu$-incoherent ($\mu_1\geq \mu$). 
Then, w.p. at least $1-\frac{1}{n^3}$: 
\begin{align}\label{eqn:sampling-minsingval-Cl}
\|(BD-C)v^*\|_2\leq  \delta_{2k} \dist(\Vt, \Vo),
\end{align}
where $v^*=vec(\Vo)$, i.e. $v^*=\left[\begin{matrix}\Vo_1\\ \vdots \\ \Vo_k\end{matrix}\right]$. 
\end{lemma}
\begin{proof}[Proof of Lemma~\ref{lemma:sampling-minsingval-Cl}]
Let $X\in \mathbb{R}^{n\times k}$ and let $x=vec(X)\in \mathbb{R}^{nk}$ s.t. $\|x\|_2=1$. Also, let $x_p$ be the $p$-th column of $X$ and $x^j$ be the $j$-th column of $X$. 

Let $u^i=(\Ut)^{(i)}$ and $u^{*(i)}=(\Uo)^{(i)}$. Also, let $H^j=(B^jD-C^j)$, i.e., 
$$H^j=\frac{1}{p}\sum_{i:(i,j)\in \Omega}u^i(u^i)^\dag (\Ut)^\dag \Uo - u^i (u^{*(i)})^\dag =\frac{1}{p}\sum_{i:(i,j)\in \Omega} H^j_i,$$ where $H^j_i\in \mathbb{R}^{k\times k}$. Note that, 
\begin{equation}\label{eq:ch_k}\sum_{i}H^j_i=(\Ut)^\dag \Ut (\Ut)^\dag\Uo - (\Ut)^\dag \Uo = 0.\end{equation}
Now, $x^\dag(BD-C)v^*= \sum_j(x^j)^\dag(B^jD-C^j)(\Vo)^{(j)}=\frac{1}{p}\sum_{pq}\sum_{(i,j)\in \Omega}x^j_p\Vo_{jq}(H^j_i)_{pq}$. Also, using \eqref{eq:ch_k}, $\forall (p,q)$: $$\sum_{i}(H^j_i)_{pq}=0.$$
Hence, applying Lemma~\ref{lem:komlem}, we get w.p. at least $1-\frac{1}{n^3}$:
\begin{equation}\label{eq:bcdbound_compk}x^\dag (BD-C)v^*=\sum_j(x^j)^\dag(B^jD-C^j)(\Vo)^{(j)}\leq \frac{1}{p}\sum_{pq}\sqrt{\sum_j(x^j_p)^2(\Vo_{jq})^2}\sqrt{\sum_i(H^j_i)_{pq}^2}.\end{equation}
Also, 
\begin{align}
\sum_i(H^j_i)_{pq}^2&=\sum_i(u^i_p)^2((u^i)^\dag (\Ut)^\dag\Uo_q-\Uo_{iq})^2\leq \max_i(u^i_p)^2\sum_i((u^i)^\dag (\Ut)^\dag\Uo_q-\Uo_{iq})^2\nonumber\\&=\max_i(u^i_p)^2 (1-\|\Ut\Uo_q\|_2^2) \leq \frac{\mu_1^2k}{m}\dist(\Ut, \Uo)^2.\label{eq:hbound_compk}\end{align}
Using \eqref{eq:bcdbound_compk}, \eqref{eq:hbound_compk} and incoherence of $\Vo$, we get (w.p. $1-1/n^3$), $\forall x$: 
$$x^\dag(BD-C)v^*\leq \sum_{pq}\frac{\mu_1^2k}{mp}\dist(\Ut, \Uo)\|x_p\|_2\leq \delta_{2k}\dist(\Ut, \Uo),$$
where we used the fact that $\sum_p \twonorm{x_p} \leq \sqrt{k}\twonorm{x} = \sqrt{k}$ in the last step.
Lemma now follows by observing $\max_{x, \|x\|=1} x^\dag(BD-C)v^*=\|(BD-C)v^*\|_2$. 
%and $H^j_i=$Now, $\sum_j(x^j)^\dag(B^jD-C^j)(\Ut)^{(j)}=\sum_j\sum_{p,q}x^j_p(\Ut)^{(j)}_qH^j_{pq}$ Note that $H^j_{pq}=\sum_{i:(i,j)\in \Omega}\Ut_{ip}$
%$\sum_j(x^j)^\dag(\sum_i\delta_{ij}(\Ut)^{(i)}(\Ut)^{(i)\dag} (\Ut)^\dag\Uo(\Vt)^{(j)}-(\Ut)^{(i)}(\Uo)^{(i)\dag}(\Vt)^{(j)}$
\end{proof}

%%% Local Variables: 
%%% mode: latex
%%% TeX-master: "altmin"
%%% End: 

%Appendix A


% That's all folks!
\end{document}
